% MRD Stage 0--2: what we now know about the questions from 2026-10-08.
% Build: bash build.sh --full (figures + pdflatex) writes report_full.pdf; plain bash build.sh builds the short
% core report (report_core.tex -> report_core.pdf). See README.md.
%
% FILTERING SECTIONS: every section is one \input line below. Comment a line out (prefix %) to drop
% that section; numbering, the table of contents and the "Answers at a glance" table update on the
% next build (a dropped section shows as "not included" there; delete its row if you prefer).
% Sections refer to each other by name, not number, so any subset compiles.
\documentclass[11pt]{article}
% Shared preamble for the 2026-10-08 questions report. Section files use only these macros.
\usepackage[T1]{fontenc}
\usepackage[utf8]{inputenc}
\usepackage{lmodern}
\usepackage{microtype}
\usepackage[letterpaper,margin=1in]{geometry}
\usepackage{graphicx}
\usepackage{booktabs}
\usepackage{tabularx}
\usepackage{array}
\usepackage{longtable}
\usepackage{multirow}
\usepackage{amsmath}
\usepackage{amssymb}
\usepackage{enumitem}
\usepackage{caption}
\usepackage{float}
\usepackage{parskip}
\usepackage[dvipsnames]{xcolor}
\usepackage{tcolorbox}
\usepackage{fancyhdr}
\usepackage{titlesec}
\usepackage[hidelinks]{hyperref}

\definecolor{inkmuted}{HTML}{52514E}
\definecolor{boxtake}{HTML}{EAF2FC}
\definecolor{boxtakeedge}{HTML}{2A78D6}
\definecolor{boxcaveat}{HTML}{F4F3EF}
\definecolor{boxcaveatedge}{HTML}{898781}
\definecolor{boxteam}{HTML}{FDF0EA}
\definecolor{boxteamedge}{HTML}{EB6834}

\setlist{itemsep=2pt,topsep=3pt,parsep=0pt}
\captionsetup{font=small,labelfont=bf,margin=0pt,skip=4pt}
\setlength{\parskip}{5pt}

\pagestyle{fancy}
\fancyhf{}
\fancyhead[L]{\small\color{inkmuted}MRD Stage 0--2: answers to the 2026-10-08 questions}
\fancyhead[R]{\small\color{inkmuted}\thepage}
\renewcommand{\headrulewidth}{0.3pt}

\titleformat{\section}{\Large\bfseries}{\thesection}{0.6em}{}
\titleformat{\subsection}{\large\bfseries}{\thesubsection}{0.6em}{}

% The short answer to the section's question (one per section, at the top).
\newtcolorbox{answerbox}{colback=boxtake,colframe=boxtakeedge,boxrule=0.6pt,arc=2pt,left=6pt,right=6pt,top=4pt,bottom=4pt,fonttitle=\bfseries\small,title=Short answer}
% What this result does NOT show (claim boundary).
\newtcolorbox{caveatbox}{colback=boxcaveat,colframe=boxcaveatedge,boxrule=0.6pt,arc=2pt,left=6pt,right=6pt,top=4pt,bottom=4pt,fonttitle=\bfseries\small,coltitle=black,title=What this does not show}
% A sentence the user can say to the team, and the likely follow-up question.
\newtcolorbox{teambox}{colback=boxteam,colframe=boxteamedge,boxrule=0.6pt,arc=2pt,left=6pt,right=6pt,top=4pt,bottom=4pt,fonttitle=\bfseries\small,title=For the team}
% A concept explainer (intuition first, then the formula). Drawn as a left rule with a plain title, so it
% cannot be mistaken for the framed gray "What this does not show" box.
\newtcolorbox{explainbox}[1]{colback=white,colframe=inkmuted!55,colbacktitle=white,coltitle=black,boxrule=0pt,leftrule=2.5pt,titlerule=0pt,arc=0pt,left=8pt,right=6pt,top=3pt,bottom=4pt,toptitle=2pt,bottomtitle=0pt,before upper={\setlength{\parskip}{3pt}},fonttitle=\bfseries\small,title={#1}}

% Source pointer: devlog id / analysis id, set small and muted.
\newcommand{\src}[1]{{\small\color{inkmuted}[Source: #1]}}
\newcommand{\pp}{\,pp}
\newcommand{\lift}{\mathrm{lift}}


% Long code paths: let lines stretch rather than overflow the margin.
\setlength{\emergencystretch}{3em}

% Let the answer/caveat/team/explainer boxes split across pages, so a tall box no longer leaves a
% half-empty page before it (the box definitions in preamble.tex are unchanged).
\tcbuselibrary{breakable}
\tcbset{breakable, lines before break=4}
% ... except the short answer and the team box, which are short and read as one unit: they never split
% (a box that does not fit moves to the next page whole). A caveat box does not start in the last six
% lines of a page. (\AtBeginEnvironment and \BeforeBeginEnvironment are from etoolbox, which tcolorbox
% loads; \AtBeginEnvironment runs inside the environment's group, so the setting stays local to that box.)
% \softneedspace{len} starts a new page only when less than len is left. Unlike \needspace, it adds no
% \penalty-100 break point, which on a ragged-bottom page could end the page early and leave it half
% empty (version 2 lost about half of a page before a caveat box this way).
\makeatletter
\newcommand{\softneedspace}[1]{\par\begingroup\dimen@ii\pagegoal\advance\dimen@ii-\pagetotal
  \ifdim #1>\dimen@ii\ifdim\dimen@ii>\z@\newpage\fi\fi\endgroup}
\makeatother
% Unique link targets for tables. Every longtable steps the table counter, and the uncaptioned ones
% (revision list, glance table, section 8's list) step it back, so two tables could share the PDF
% link target "table.N" (pdfTeX warned "duplicate ignored" and a link could jump to the wrong table).
% Counting longtables in the hyperref name keeps each target unique; printed numbers are unchanged.
\newcounter{reportlongtables}
\renewcommand{\theHtable}{\arabic{table}.\arabic{reportlongtables}}
\AtEndEnvironment{longtable}{\stepcounter{reportlongtables}}
\AtBeginEnvironment{answerbox}{\tcbset{unbreakable}}
\AtBeginEnvironment{teambox}{\tcbset{unbreakable}}
\BeforeBeginEnvironment{caveatbox}{\softneedspace{6\baselineskip}}
% No single stranded line at the top or bottom of a page, inside the boxes too.
\widowpenalty=10000
\clubpenalty=10000

% Float placement for the assembled report. The section files ask for [H] (fixed position), which
% is right for a standalone section but leaves half-empty pages when a tall figure misses the end
% of a page. Here every figure and table floats (!htbp) instead, and \FloatBarrier at each section
% start (see \labelnextsection) keeps a section's floats inside that section.
\usepackage{placeins}
\usepackage{needspace}
\let\reportorigfigure\figure \let\endreportorigfigure\endfigure
\renewenvironment{figure}[1][]{\reportorigfigure[!htbp]}{\endreportorigfigure}
\let\reportorigtable\table \let\endreportorigtable\endtable
\renewenvironment{table}[1][]{\reportorigtable[!htbp]}{\endreportorigtable}
\renewcommand{\topfraction}{0.9}
\renewcommand{\bottomfraction}{0.6}
\renewcommand{\textfraction}{0.08}
\renewcommand{\floatpagefraction}{0.75}
\setcounter{topnumber}{3}
\setcounter{totalnumber}{4}

% \labelnextsection{key}: label the NEXT \section with key, so the glance table can cite each section
% without depending on the section files' own labels. Kept on the same line as the section's \input.
% It also flushes the previous section's floats and starts a new page when less than about a third
% of the page is left, so a section heading and its question never sit alone at the bottom of a page.
\makeatletter
\newcommand{\labelnextsection}[1]{%
  \FloatBarrier
  \needspace{0.33\textheight}%
  \begingroup
  \phantomsection
  \edef\@currentlabel{\the\numexpr\value{section}+1\relax}%
  \label{#1}%
  \endgroup}
% \labelnextsectionb{key}: the same for a companion section numbered after the previous one with a
% "b" (section 2b follows section 2), so the numbers of the later sections stay those of version 2.
% \plainsectionnumbers restores plain numbering after it.
\newcommand{\labelnextsectionb}[1]{%
  \FloatBarrier
  \needspace{0.33\textheight}%
  \addtocounter{section}{-1}%
  \renewcommand{\thesection}{\arabic{section}b}%
  \renewcommand{\theHsection}{\arabic{section}b}%
  \begingroup
  \phantomsection
  \edef\@currentlabel{\the\numexpr\value{section}+1\relax b}%
  \label{#1}%
  \endgroup}
\newcommand{\plainsectionnumbers}{%
  \renewcommand{\thesection}{\arabic{section}}%
  \renewcommand{\theHsection}{\arabic{section}}}
% \glanceref{key}: the section number if that section is included, otherwise a dash.
\newcommand{\glanceref}[1]{\@ifundefined{r@#1}{\textcolor{inkmuted}{--}}{\ref{#1}}}
\makeatother

\newcommand{\swatch}[2]{\fcolorbox{#2}{#1}{\rule{0pt}{5pt}\rule{5pt}{0pt}}}
\newcommand{\barswatch}{{\color{inkmuted!55}\rule[-1pt]{2.5pt}{9pt}}\hspace{6.8pt}}

% Team build: `bash build.sh --team` defines \TEAMBUILD and writes report_full_team.pdf without Appendix C.

\hypersetup{
  pdftitle={MRD Stage 0--2: what we now know about the questions from 2026-10-08},
  pdfsubject={Working report for the project lead; descriptive results, seven supported patients},
}

\begin{document}

% ---------------------------------------------------------------- title block
\thispagestyle{plain}
\begin{flushleft}
{\LARGE\bfseries MRD Stage 0--2: what we now know about the questions from 2026-10-08\par}
\vspace{6pt}
{\large\color{inkmuted} A working report for the project lead; descriptive results, seven supported patients\par}
\vspace{10pt}
{\small
\begin{tabular}{@{}l@{\hspace{1em}}p{0.78\linewidth}@{}}
\textbf{Date} & Version 3.2, 10 October 2026 (results through the D26 leverage run of 2026-10-10, now its own section 2b, your decisions of that day, the surface-protein explainers of that evening and the CD63 and total-ADT follow-up of that night, D29, in section 5); version 2 of 9--10 October changed the report in answer to your reading notes; version 1 of 8--9 October is kept as \texttt{report\_v1\_2026-10-08.pdf} \\
\textbf{Data} & Experiment \nolinkurl{20260525_060508_stage0_mrd_old34_broad_screen_82db5093} (EXP; 60,118 MRD CITE-seq cells, 13 patients) \\
\textbf{Code} & Repository \texttt{sc\_classification}, branch \texttt{sc\_classification\_gpu}, commit \texttt{f5a8f4d}, plus uncommitted changes: the full-design patient-count curve entry (D15b) with its registry rows, the 2026-10-09 and 2026-10-10 analyses (D26, D27, D28) and this report \\
\textbf{Length} & This is the long companion (\texttt{report\_full.pdf}): every section in full, for checking and for reference. The short read is \texttt{report\_core.pdf} (one page per question, at most 12 pages); its pointers use this report's section numbers. Each section here stands alone \\
\textbf{Status} & Draft for your review; filter sections before sharing (see \texttt{README.md}). Appendix C is for you only: \texttt{bash build.sh -{}-team} writes \texttt{report\_full\_team.pdf} without it \\
\end{tabular}\par}
\end{flushleft}
\vspace{4pt}

\subsection*{How to read this report}

Each numbered section answers one question from the 2026-10-08 team report, quoted at its top, and stands alone: you can drop any section without breaking the others. The primer section, \emph{The cohort and the yardstick}, fixes the cohort facts and the effect size (\emph{lift}: held-out AUPRC minus that patient's malignant prevalence $\pi$, in percentage points, pp); Appendix~A defines every other term. Inside a section the coloured boxes always mean the same thing:

\begin{itemize}[leftmargin=1.4em]
\item[\swatch{boxtake}{boxtakeedge}] \textbf{Short answer}: the answer, its key numbers, and how sure we are.
\item[\barswatch] \textbf{Explainer}: a concept introduced in that section, intuition first, then the exact quantity.
\item[\swatch{boxcaveat}{boxcaveatedge}] \textbf{What this does not show}: the claim boundary; read it before quoting the answer.
\item[\swatch{boxteam}{boxteamedge}] \textbf{For the team}: one sentence you can say, the question it will likely draw, and a crisp answer.
\end{itemize}

All numbers are \emph{descriptive}: effect sizes measured on these data, with no calibrated test, no patient-level interval and no outside cohort behind them, so they describe these patients and bases and claim nothing beyond them. They come from seven supported patients (at least 10 malignant cells and both classes), mostly on a Stage~1 basis fit transductively on all cells. The study's claim rules (THEME \S6, with the claim levels of the devlog rules) name two higher levels that no result here reaches: \emph{inferential} (a calibrated null, patient-level uncertainty and nested selection) and \emph{publication} (every reviewer alternative addressed, every blocking task done, evidence from committed code). When a section compares classifiers it quotes pooled AP, median lift and P20 lift together, because they can rank methods differently. Grey bracketed pointers such as \src{2026-10-08\_d19-regularization-contract} name the devlog entry that holds each number; Appendix~B maps sections to code and artifacts.

\subsection*{Revision 2026-10-10, late evening (version 3.2)}

\begin{itemize}[leftmargin=1.4em]
\item \textbf{Your questions on Figures 20--23} (asks A-20261010-26 to -37; devlog 2026-10-10\_d29-cd63-depth-and-every-card) are answered in section \emph{What is PC8}. Figure~\ref{fig:05rho}'s caption says in bold that each gray dot is one malignant cell, and why the cloud looks random. ``Every card'' now says where the 289,961 pairs come from, where Stage~2 enters, what $\eta^2$ of patient is and how to read panels c--e, with a statement we can make (box); Figure~\ref{fig:05f} no longer colours by method. A new subsection, \emph{CD63, followed up}, adds analysis D29: CD63 cell by cell (Figures~\ref{fig:05n}--\ref{fig:05o}), a direct malignant-vs-healthy protein contrast inside one sample and predicted type (Figure~\ref{fig:05q}), and the two checks with each cell's total ADT count held fixed (Figure~\ref{fig:05r}), which is the isotype question, with four questions for the Bachireddy lab. The new figures follow Figure~\ref{fig:05f}, so Figures 19--23 keep their numbers and every later figure moves up by four.
\end{itemize}

\subsection*{Revision 2026-10-10, evening (version 3.1)}

\begin{itemize}[leftmargin=1.4em]
\item \textbf{Your questions on the protein part of section 5} (asks A-20261010-14 to -24; devlog 2026-10-10\_protein-explainer-v3) are answered inside the report, not only in the Pavan update directory. Section \emph{What is PC8}: Figure~\ref{fig:05c} is relabelled (malignant / healthy counts, which types are left out, where the malignant cells come from, and the definition of the mean oriented PC8 score in its caption); Table~\ref{tab:05adt} has plain column headers, the isotype reference per column and PC8's other strong malignant-cell correlates (CD116, CD33, CD146); the new Figure~\ref{fig:05rho} shows how one $\rho$ is computed from a score vector and a protein vector and what its sign means; the new Figure~\ref{fig:05j} sets the within-patient value of every protein against its value inside cell types, the two checks each protein gets (consistent across patients; persists within cell types), so the text now says which check CD99, CD48, CD116, CD33 and CD146 pass. Version 3's Figure 20 is now Figure~\ref{fig:05e} and its Figure 21 is now Figure~\ref{fig:05f}; every later figure is two higher (e.g.\ old 22 $\to$ 24). The update directory's working names F5, F6 and F7 are Figures~\ref{fig:05rho}, \ref{fig:05j} and \ref{fig:05c} here. The short read (\texttt{report\_core.pdf}, still 12 pages) shows compact versions: page 7, Figure 6 bottom (the cell-type means); page 8, Table 1 (Table~\ref{tab:05adt}'s PC8 malignant-cell rows) and Figure 7 (the sign of $\rho$ and the two checks). This PDF, the figure files and the update's PNGs are gitignored, so a clone of the repository has only the short read.
\end{itemize}

\subsection*{Revision 2026-10-10 (version 3)}

\begin{itemize}[leftmargin=1.4em]
\item \textbf{New section 2b}, \emph{What did we do with allograft rejection and apoptosis?}: why these two panels, the leverage test (D26) as run, the rule fixed before the run, what it found, and what it means for re-weighting and for healthy-cell downsampling. Section~2 keeps IFN-$\alpha$ and points there; sections 3--9 keep their version-2 numbers.
\item \textbf{Your decisions of 2026-10-10 are recorded where they apply}: new full runs keep only patients with at least 100 malignant cells (section \emph{Remaining TODOs}, \emph{Proposed for the new data}, with what that does to the summaries; A-20261010-07); re-weighting targets the Stage~2 labels (sections 2 and 2b; A-20261010-08); the Git LFS lines are gone (Appendix~C; A-20261010-10). Devlog 2026-10-10\_decisions-new-run-inclusion-reweighting-lfs records all three.
\item \textbf{Fixes from two agent reviews of version 2}: one account of where a downsampling gain could come from (section~2 and \emph{Remaining TODOs} now agree, and the latter summarizes your earlier donor downsampling, D10 and D11); section~8 brought up to date with D26; claims trimmed where they outran their devlog gates; a prevalence mix-up in section~2; figure labels that collided; glossary entries.
\end{itemize}

\subsection*{Revision 2026-10-09 (version 2): where your reading notes are taken up}

Your reading notes on version 1 (asks A-20261009-01 to -16) are not answered here; this list says where version 2 changed so that the report answers them. Figure and table numbers have shifted: the sections on shaping the factors, on aggregating patients and on PC8 gained figures, so the figure numbers in your notes refer to version 1.

{\footnotesize
\setlength{\tabcolsep}{4pt}
\renewcommand{\arraystretch}{1.12}
\begin{longtable}{@{}>{\raggedright\arraybackslash}p{0.33\linewidth} >{\raggedright\arraybackslash}p{0.635\linewidth}@{}}
\toprule
\textbf{Note (ask)} & \textbf{Where to look in version 2} \\
\midrule
\endhead
01 What ``descriptive'' means & This page (``How to read'') and the paragraph above the Answers-at-a-glance table, which also names the two higher levels; Appendix~A, entry ``Claim levels''. \\
02 The leverage test; why it stays open; other panels & Glance-table term (a). Section \emph{Do the highest-prevalence patients shape the factors?}: box ``The leverage test, step by step''. The run you started (D26) settles the question for PCA and FA: section 2b, \emph{What did we do with allograft rejection and apoptosis?} (version 3). Other panels: section 2b, Next step (W2.8). \\
03 How re-weighting works; its percentages; Stage~1 or Stage~2 & Glance-table term (b). Same section: box ``How patient-equal re-weighting works, and what its percentages mean'' and the list of Stage~1 and Stage~2 options. You chose the Stage~2 labels on 2026-10-10 (sections 2 and 2b). \\
04 ``Aggregate'' in numbers & Glance row for section \emph{Should we aggregate patients?} and term (c); that section's short answer and box ``The quantities, exactly''. \\
05 The L1 entry rule & Section \emph{IFN-$\alpha$: why does FA40 reach only P03 and P04?}, box ``The L1 entry rule, and why one $C$ is a different penalty in every fold'', which now ends with one fold in numbers, and the fold-by-fold margins after it. The same fold worked through with numbers (P02 held out: the bar, PC8's and FA7's gaps and margins, and what $C \times n$ would let FA7 in): section \emph{Learning curves}, paragraph ``One fold worked through''. \\
06 ``Neither is identifiable alone'' & Glance-table term (d). Section \emph{What is PC8}: box ``What \emph{not identifiable} means'' and a new figure. \\
07 What the 1,810 factor cards are & Glance-table term (e) and the row on the interactive app. Section \emph{What is PC8}: a new table of where the cards come from (by source, method and $K$), and the first paragraph of ``What we did'', which splits them by how the stored probe lets them be read. \\
08 Factor $\times$ cell type and surface protein for all factors; $\rho$ $+0.251$ and $+0.222$ & New analysis D27 (all 1,801 non-degenerate cards). Section \emph{What is PC8}: box ``How to read $\rho$'', a figure that takes the two numbers apart, and a figure over all cards. \\
09 Transductive vs inductive on IFN-$\alpha$ & Section \emph{Option 2} (unchanged). \\
10 Why $C$ is retuned & Section \emph{Learning curves}: the first paragraph of ``What we did''. \\
11 The 100-malignant-cell rule; P01, P06, P07 & Where the 100 comes from: the glance-table row on your 2026-10-09 proposals and section \emph{Remaining TODOs}, \emph{Options 1.1 and 3}, where the cut first appeared. Same section, subsection \emph{Proposed for the new data}: the inclusion threshold (W8.5), with a table of what 10, 20 or 100 malignant cells keep; you chose 100 on 2026-10-10 (version 3 adds what that does to the median and P20 lift). Section \emph{The cohort and the yardstick} points there (Next step). Right after that table, a list of the figures that already show P01, P06 and P07 (including P07's IFN-$\alpha$ lift); the integration plan (\texttt{workstreams/\allowbreak W8\_new\_cohort/\allowbreak integration\_plan.md}, \S4.1) has the same pointers with devlog ids. \\
12 Healthy-to-malignant ratio; healthy-cell downsampling; the paper & Section \emph{Remaining TODOs}, subsection \emph{Proposed for the new data}: the plan in brief and box ``Why fewer healthy cells could help, and why it might not''. Paragraphs on downsampling in sections \emph{Do the highest-prevalence patients shape the factors?}, 2b and \emph{Should we aggregate patients?}; your earlier donor downsampling (D10, D11) is summarized in \emph{Remaining TODOs} (version 3) and in \emph{Should we aggregate patients?}. The full draft pre-registration is \texttt{workstreams/\allowbreak W8\_new\_cohort/\allowbreak healthy\_downsampling\_plan.md}. As you decided, the run waits for the new cohort. The paper context lives in \texttt{research/paper/} (Appendix~B; task W0.13); the gradual repository update you asked for is task W0.14 (\texttt{workstreams/\allowbreak W0\_context\_provenance/\allowbreak README.md}) and is not covered here. \\
13 The two shaping tests and their metrics & Section \emph{Do the highest-prevalence patients shape the factors?}: boxes ``Score energy, intuitively'' and ``The leverage test, step by step''; section 2b, box ``The decision rule, fixed before the run''. \\
14 Figure 3 (score energy, 1.11), Figure 4 (axis loss); which patients new runs use & Same section: both figures redrawn (1.11 labelled; axis loss in percent with a second scale), plus a cell-type-adjusted figure (D28). New runs leave out patients without malignant cells (recorded in section \emph{Remaining TODOs}, subsection \emph{Proposed for the new data}). The allograft and apoptosis leverage run (D26) refits variants of the stored bases and ran as built, on all 13 patients, so it compares directly with the IFN-$\alpha$ test (section 2b, caveats). \\
15 Start the next steps & D28 done (shaping section), D27 done (section \emph{What is PC8}), D26 done after you ran it (section 2b, version 3). \\
16 The aggregation plots; normal vs malignant & Section \emph{Should we aggregate patients?}: three new figures that work one factor through the test, separately for normal and malignant cells, then show the three call types. \\
\bottomrule
\end{longtable}
\addtocounter{table}{-1}
}

\tableofcontents

% ---------------------------------------------------------------- sections
% One line per section: comment it out to drop the section.
% Section 00: answers at a glance (assembler). One row per question; the Section column resolves
% through \glanceref (report.tex) and shows a dash when that section is commented out.
% Delete a row here if you drop its section and do not want the row in the shared version.
\makeatletter
\providecommand{\glanceref}[1]{\@ifundefined{r@#1}{--}{\ref{#1}}}
\makeatother
% A sub-question row: the question cell is indented as a block.
\providecommand{\subq}{\leftskip=1em\relax}

\section*{Answers at a glance}
\phantomsection
\addcontentsline{toc}{section}{Answers at a glance}
\label{sec:glance}

One row per question from the 2026-10-08 team report, in the order of the sections that answer them. Lift is held-out AUPRC minus the patient's malignant prevalence, in percentage points (pp). Terms marked with a superscript letter are defined just below, before the table.

\textbf{What ``descriptive'' means in the ``How sure'' column.} The study's claim rules (THEME \S6, with the claim levels of the devlog rules) have three levels. \emph{Descriptive}, the level of every result in this table (rows marked \emph{Plan}, \emph{Status} or \emph{Built and checked} report work, not results, and the literature row names its own limit): an effect size measured on these data (seven supported patients, mostly a Stage~1 basis fit transductively on all cells), with no calibrated test, no patient-level interval and no outside cohort behind it, so it says what happened in these patients and bases and claims nothing beyond them. \emph{Inferential} would add a calibrated null, patient-level uncertainty and nested selection; none exists yet (section \emph{The three scientific questions}, Q2). \emph{Publication} is a claim whose reviewer alternatives are all addressed, whose blocking tasks are done and whose evidence comes from committed code. After the word ``Descriptive'' the column says what the evidence covers and what it does not.

{\footnotesize
\textbf{Terms in the table.}\par\nobreak
\begin{itemize}[leftmargin=1.6em,itemsep=1pt]
\item[a] \emph{Leverage test}: an in-silico perturbation. Refit Stage~1 without a set of patients, measure how far the malignant axis moves (axis loss: the share of the axis's variance the refit basis can no longer reproduce; matched $|r|$: correlation with the best single refit factor), and compare with 20 random drops of the same number of cells and of malignant cells (composition-matched controls). The decision rule, written for the allograft and apoptosis run after the IFN-$\alpha$ results were known (so post hoc on IFN-$\alpha$), calls it leverage beyond composition when the axis moves more than all 20 controls in at least 4 of the 6 method $\times$ $K$ cells, and asks, before calling it the patients' own leverage, that it also exceed that of every contrast patient set in at least 4. Section \emph{Do the highest-prevalence patients shape the factors?}, box ``The leverage test, step by step''; section 2b, box ``The decision rule, fixed before the run''.
\item[b] \emph{Fit-weight percentages}: each patient's share of the total weight in the Stage~1 fit. Unweighted, that is its share of the cells; patient-equal weighting gives each of the 13 patients $1/13 = 7.7\%$. Same section, box ``How patient-equal re-weighting works''.
\item[c] \emph{Overlap}: the shared area of two patients' histograms of one factor's scores (1 = identical, 0 = disjoint), computed separately for normal and malignant cells. Section \emph{Should we aggregate patients?}, worked example on one factor.
\item[d] \emph{Identifiable}: the data pin the factor's direction down uniquely, so a refit on slightly different cells returns the same factor. A small eigengap (PCA: neighbouring components explain nearly the same variance, so they can rotate into each other) or a determinacy below 0.707 (FA: the scores are uncertain estimates) means they do not, although the subspace spanned by the related factors can still be stable. Section \emph{What is PC8}, box ``What \emph{not identifiable} means''.
\item[e] \emph{Factor card}: a one-page record per factor of one stored Stage~1 basis (panel $\times$ method $\times$ $K$): gene poles, orientation, per-patient lift, matches in other bases, identifiability, cell-type and surface-protein associations. Section \emph{What is PC8}, table of where the 1,810 come from.
\item[f] \emph{Robust association}: a factor--protein correlation whose median within-patient $|\rho|$ over the seven supported patients is at least 0.20, with at least 5 of 7 patients of the same sign (fixed before computing). Section \emph{What is PC8}, box ``How to read $\rho$''.
\end{itemize}
\par}


{\footnotesize
\setlength{\tabcolsep}{4pt}
\renewcommand{\arraystretch}{1.15}
\begin{longtable}{>{\raggedright\arraybackslash}p{3.45cm} >{\raggedright\arraybackslash}p{7.55cm} >{\raggedright\arraybackslash}p{3.35cm} >{\centering\arraybackslash}p{0.75cm}}
\toprule
\textbf{Question (team wording, short)} & \textbf{Short answer} & \textbf{How sure} & \textbf{Sec.} \\
\midrule
\endfirsthead
\toprule
\textbf{Question (team wording, short)} & \textbf{Short answer} & \textbf{How sure} & \textbf{Sec.} \\
\midrule
\endhead
\midrule
\multicolumn{4}{r}{\emph{continued on the next page}} \\
\endfoot
\bottomrule
\endlastfoot
% --- section 2
Do P02, P03, P04 disproportionately shape the factors? &
Not more than other patients do, on the three panels tested. They hold about their 8.2\% cell share of factor score variance on all 41 panels (median ratio 0.96; 1.04 after adjusting for cell type). In the leverage test\textsuperscript{a}, refitting Stage~1 without them barely moves the IFN-$\alpha$ malignant axis (matched $|r| \geq 0.988$; above all 20 composition-matched controls in 3 of 6 method $\times$ $K$ cells, the rule needs 4); on allograft rejection and apoptosis see the next row. They dominate the Stage~2 labels (68.3\% of malignant cells). &
Descriptive; PCA/FA, no FactoSig; rule post hoc on IFN-$\alpha$; transductive bases &
\glanceref{glance:02} \\
\addlinespace
\subq Option 1.2: re-weight so each patient counts equally &
Patient-equal weights give every patient 1/13 of the Stage~1 fit, so P02--P04 go from 8.2\% to 23.1\% of the fit weight\textsuperscript{b} (P09 from 28.5\% to 7.7\%); Stage~2 labels are untouched. Piloted on IFN-$\alpha$: PC8 rotates (matched $|r|$ 0.90); no clear loss under tuned L1 (median per-patient change PCA40 $+1.35$, FA40 $+0.07$\pp), but FA40 loses a median 1.13\pp{} under L2. You chose the target on 2026-10-10: the Stage~2 labels (per-patient, per-class sample weights in the logistic loss), in the runs with the new patient data (W2.3). &
Descriptive; one panel, transductive pilot &
\glanceref{glance:02} \\
\addlinespace
\subq Why is P01 the most over-represented and most leveraged patient? &
Its score-energy excess is mostly matched by its cell-type mix (median ratio 1.24, 1.03 after adjustment; ncMono, pDC2 and cDC cells carry more energy per cell); this is arithmetic, not cause. Its leverage, the largest of any set on the primary axis in 17 of 18 method $\times$ $K$ cells of the three tested panels (IFN-$\alpha$ PCA40 15$\times$ its controls; allograft 21--141$\times$, apoptosis 48--128$\times$), is not explained. The IFN-$\alpha$ 1.11 of P02+P03+P04 is not cell-type mix (1.11 after adjustment). &
Descriptive; predicted cell types; reference-dependent for P01 (1.03 or 1.08) &
\glanceref{glance:02} \\
\addlinespace
% --- section 2b
What did we do with allograft rejection and apoptosis? (2026-10-10) &
Ran the leverage test\textsuperscript{a} on both (D26, your screen run; PCA and FA, $K$ 10/20/40). Leaving P02+P03+P04 out moves the malignant axis more than all 20 composition-matched drops in 5 of 6 (allograft) and 6 of 6 (apoptosis) cells, 1.3--4.1$\times$ and 2.0--4.0$\times$ their mean, mostly through P03; the axis stays matched ($|r| \geq 0.945$ and $\geq 0.994$). Leaving P01 out moves it 21--141$\times$ and 48--128$\times$, and the group never beats every contrast set on the primary axis (0 of 6 cells on each panel; 1 of 6 on the apoptosis probe axis), so by the rule fixed before the run H-S1-05 is not supported. For re-weighting this points to the Stage~2 labels (your choice); for downsampling, to P01. &
Descriptive; PCA/FA on two panels, no FactoSig; rule fixed before the run; transductive 13-patient bases &
\glanceref{glance:02b} \\
\addlinespace
% --- section 3
Should we aggregate patients? (per-patient factor distributions) &
Yes, on the current bases. ``Aggregate'' means that, averaged over factors with variance weights, the median pair of patients shares at least 75\% of its normal-cell and at least 45\% of its malignant-cell factor-score histogram (overlap\textsuperscript{c}), no patient's normal cells share less than 60\% with the rest, and the malignant overlap is less than 0.15 below that of label-shuffled patients. 45 of 47 representations pass, 2 are location shifts, none is a subtype candidate. The test measures where each patient's cells land, not how much each patient contributes to a factor. P03's only private axis is in broad gene bags, in none of the 41 biological representations tested (one per panel). &
Descriptive; transductive basis; uncalibrated gates, permutation $p$ at its floor &
\glanceref{glance:03} \\
\addlinespace
% --- section 4
IFN-$\alpha$: why does PCA40 reach four patients but FA40 only P03 and P04? &
A Stage~2 L1 entry threshold, not missing FA structure: a factor enters the L1 probe only if its standardized class-mean gap exceeds $4/(C\,n_{\text{train}})$, and FA7 misses that bar by 1.4--3.3\% in P02, P06, P07. With $C$ tuned per fold, FA40 has positive lift in 7/7 patients. &
Descriptive; one panel; the rule reproduces 2,772/2,772 audited atlas folds &
\glanceref{glance:04} \\
\addlinespace
\subq Hypothesis (a): FA keeps allocating factors to the highest-prevalence patient &
Retired. FA40 does not favour the highest-prevalence patients more than PCA40: P02--P04 hold 8.4\% of FA40's and 9.1\% of PCA40's score energy against 8.2\% of cells, and P03 tops none of either basis's factors. P02 tops most factors in both (27/36 FA40, 31/40 PCA40), through excess spread on low-variance trailing factors, so it does not separate FA from PCA. &
Descriptive; one panel &
\glanceref{glance:04} \\
\addlinespace
% --- section 5
PC8: two opposite-sign gene groups. What is it? &
One oriented axis: malignant cells score toward TXNIP, IFITM2, TENT5A (the team's ``negative'' group). FA7 is related, not identical (loading cosine 0.69). Neither is identifiable\textsuperscript{d} alone: the data pin down a small subspace around each (every Stage~1 refit still reproduces the stored axis, $R^2 \geq 0.977$), but not which single direction in it a fit calls ``PC8'' or ``FA7'' (PC8's variance is within 5.5\% of PC7's; FA7's scores have determinacy 0.649, below 0.707). So name the axis by its two poles. &
Descriptive; one panel; stored L1 probe &
\glanceref{glance:05} \\
\addlinespace
\subq Can we view factors in the interactive app? &
Yes, as a separate page: the read-only factor explorer shows all 1,810 factor cards\textsuperscript{e} and is linked from every atlas cell. The 1,810 are every factor of 56 of the atlas's 736 stored DR bases (one basis per panel, chosen by patient-median lift, plus 9 IFN-$\alpha$ and focus-panel bases), not of every Stage~0 $\times$ Stage~1 combination. &
Built and checked &
\glanceref{glance:05} \\
\addlinespace
\subq Correlate factors with cell types and surface protein &
Done for all 1,801 non-degenerate cards, now also within patients and within predicted cell types. PC8 (oriented toward malignant) tracks CD34 in normal cells in every patient (median $\rho$ $+0.245$, 7/7; pooled $+0.251$) but not inside cell types ($+0.019$): a cell-identity association. Its malignant-cell CD99 correlate (pooled $+0.222$) is $+0.159$ within patients (6/7) and fails the robust rule\textsuperscript{f}. Across cards, normal-cell correlates hold within patients but mostly track cell type; pooled malignant correlates are mostly between-patient. &
Descriptive; a screen over correlated factors; CLR ADT without dsb; patient nested in batch; recipient origin &
\glanceref{glance:05} \\
\addlinespace
% --- section 6
Option 2: exclude the held-out patient from Stage~1 &
Done on IFN-$\alpha$: lift stays positive in 7/7 patients under tuned $C$ or L2; median per-patient change $-1.29$ to $-0.05$\pp{} over 18 settings, so transductive numbers are slightly optimistic per patient. Pooled AP can fall more (FA40, tuned L2: 0.1325 to 0.0794). &
Descriptive; one panel, PCA/FA at K 10/20/40; no test &
\glanceref{glance:06} \\
\addlinespace
% --- section 7
AUPRC vs number of cells or patients (learning curve) &
Drawn. $C$ is retuned for every training set, because a fixed $C$ is a stronger penalty with fewer cells (at fixed $C$ the IFN-$\alpha$ curve is 0 up to 25,000 training cells). With $C$ retuned, one supported training patient (the six low-malignancy patients are always in training) already gives positive lift (median 4.45--8.79\pp); going from two to six adds 0.30--2.34\pp, with CIs that include 0 for three of four representations. &
Descriptive; frozen transductive basis; nothing extrapolated &
\glanceref{glance:07} \\
\addlinespace
\subq Useful for the grant? &
As feasibility, yes: transfer works with one supported training patient, on a label-free basis fit on all patients. As evidence that sequencing more patients raises AUPRC, not yet: seven patients are too few to measure scaling. &
Descriptive &
\glanceref{glance:07} \\
\addlinespace
% --- section 8
Q1. Shared or patient-specific? &
Tested directly on IFN-$\alpha$ only: one axis (PC8, or the related FA7) carries held-out lift in five of seven patients, including low-prevalence P07 and P09. Across the atlas, lift above 5\pp{} concentrates in P02, P03, P04 (64--80\% of the 640 biological DR cells vs at most 23\% for the others), so shared transfer and prevalence are not yet separated. &
Descriptive; transductive; 7 patients &
\glanceref{glance:08} \\
\addlinespace
Q2. How strong is the evidence? &
Effect sizes are well defined (lift, $\eta^2$, overlap), but there is no calibrated null, no patient-level interval on the atlas and no nested selection yet. &
Not yet inferential &
\glanceref{glance:08} \\
\addlinespace
Q3. What does DR add over direct genes? &
It depends on the estimand: the best DR cell (picked by pooled AP) beats direct genes in 38/41 programs on pooled AP, 16/41 on median lift and 17/41 on P20 lift; large gains in a few (allograft, apoptosis, IFN-$\alpha$). &
Descriptive; best DR chosen in-sample (not nested) &
\glanceref{glance:08} \\
\addlinespace
% --- section 9
New MRD patients from Yu-Hsiang &
The email and a pre-registered integration plan are drafted; the email is not sent. These patients are the external test for transfer and for the learning curve. &
Status &
\glanceref{glance:09} \\
\addlinespace
\subq Your 2026-10-09 proposals: healthy-cell downsampling; which patients new runs use &
Not run, by your decision: the downsampling design is pre-registered (draft, not locked) to run with the new patients. New full runs leave out patients without malignant cells and, by your decision of 2026-10-10, keep only patients with at least 100 malignant cells (below it: left out of every step, our reading, to confirm at the plan lock). On the current cohort that keeps P02, P03, P04 and P09 (any cut from 46 to 102 does), three of them the highest-prevalence patients, so medians and P20 lift over them are not comparable with today's seven-patient summaries (on the four learning-curve representations the median lift is 1.8--2.7 times higher). The 100 started as an illustrative cut for the team's Option 1.1; the study's support rule stays at 10. &
Plan &
\glanceref{glance:09} \\
\addlinespace
\subq Options 1.1 and 3: focus on P02, P03, P04 &
No new pipeline needed yet: training on two of them plus the unsupported patients and scoring the third differs from full LOPO by a median $+0.13$\pp{} over 12 folds, but single folds range from $-8.18$ to $+10.09$\pp{} and follow the held-out patient (P02 lower in all four representations, P04 higher in all four). &
Descriptive; 12 folds, 4 representations &
\glanceref{glance:09} \\
\addlinespace
\subq Days since transplant and before relapse &
Plotted. With seven patients, timing cannot be separated from patient: median lift tracks prevalence ($\rho$ 0.89) more closely than days since transplant ($\rho$ 0.57) or before relapse ($\rho$ 0.18), but at $n = 7$ these correlations cannot be told apart. &
Descriptive; $n = 7$ &
\glanceref{glance:09} \\
\addlinespace
\subq Ankit: factors learned jointly with the labels &
Deferred until the new patients arrive. Baseline to beat (inductive IFN-$\alpha$ PCA40, tuned L1): pooled AP 0.1164, median lift 10.55\pp, P20 lift 1.29\pp. &
Plan &
\glanceref{glance:09} \\
\addlinespace
\subq Review single-cell ML literature &
Done: two reviews (159 sources) with 32 recommended actions, each routed to a task. &
Agent-run reviews; their numbers count only once reproduced &
\glanceref{glance:09} \\
\end{longtable}
}
% The table above has no caption, but longtable still steps the table counter: undo it so the first
% captioned table is Table 1.
\addtocounter{table}{-1}

\subsection*{What changed in our understanding since the team report}

\begin{itemize}
\item \textbf{On the three panels tested, the imbalance specific to P02--P04 sits mainly in Stage~2.} P02, P03 and P04 hold about their cell share of factor score variance on 41 panels. On IFN-$\alpha$ their leverage on the discriminative axis is mostly what removing that many malignant cells does; on allograft rejection and apoptosis it goes beyond that, mostly through P03, but never beyond what leaving out other patients does (D26, rule fixed before the run). The patient with the most leverage is P01, a low-prevalence patient with 19.7\% of the cells (15--24$\times$ its composition-matched control under PCA); its excess score energy is mostly matched by its cell-type mix (1.24, 1.03 after adjustment; arithmetic on predicted labels, not cause), but its leverage is not explained. Re-weighting Stage~1 is therefore no longer the obvious fix; you chose to target the Stage~2 labels (2026-10-10). \src{2026-10-08\_d16-factor-patient-matrix; 2026-10-08\_d17-composition-matched-controls; 2026-10-09\_d28-celltype-score-shares; 2026-10-10\_d26-leverage-allograft-apoptosis-done}
\item \textbf{A zero in the atlas is often a threshold, not an absence of signal.} The IFN-$\alpha$ PCA/FA gap is the L1 entry rule $4/(C\,n_{\text{train}})$, which reproduces every audited atlas fold. 528 of 783 atlas cells select the lowest $C$ on the grid, and 136 folds in 49 cells sit at the knife-edge of entry, so method comparisons at a fixed $C$ need tuned-$C$ checks. \src{2026-10-08\_d19-regularization-contract}
\item \textbf{PC8's ``negative'' gene group is the malignant pole.} Once the axis is oriented by its probe coefficient, malignant cells score toward TXNIP, IFITM2 and TENT5A. The direction recurs across bases (PC8 has $R^2$ 0.920 on the FA40 scores) but is not identifiable as ``factor 8'': leaving P01 out of Stage~1 turns PC8's best single match to $|r|$ 0.58 while 98.8\% of it stays in the refit basis. So we name it by its poles. \src{2026-10-08\_d18-factor-cards; 2026-10-08\_d17-basis-variant-pilot}
\item \textbf{A factor's surface-protein correlate has to be read inside patients and inside cell types.} Over all 1,801 non-degenerate cards, normal-cell correlates hold within patients (median within/pooled ratio 0.89) but only 729 of 20,274 robust pairs keep half their size inside predicted cell types; pooled malignant correlates are mostly between-patient (ratio 0.24). PC8's CD34 correlate is a cell-identity axis, and its CD99 correlate fails the within-patient rule. \src{2026-10-09\_d27-factor-modality-assoc}
\item \textbf{Transductive optimism is small per patient on IFN-$\alpha$.} Fitting Stage~1 without the held-out patient costs a median 0.05--1.29\pp{} per patient; pooled AP can fall more (FA40, tuned L2: 0.1325 to 0.0794). The axis survives without the held-out patient, without P02--P04 and without malignant cells. This is one panel; the other representations are not yet tested. \src{2026-10-08\_d17-basis-variant-pilot}
\item \textbf{The full patient-count curve replaces the reduced-design numbers in the 2026-10-08 team update.} The median PCA40 gain from two to six training patients is 1.67\pp{} [95\% CI $-0.44$, 2.89], not 2.16 [0.20, 2.48]; only allograft's CI now excludes 0. ``More patients help'' stays out of the claims. \src{2026-10-08\_d15b-full-design; 2026-10-08\_d15b-patient-count-curve}
\end{itemize}
                                            % Answers at a glance (unnumbered)
\labelnextsection{glance:01}\section{The cohort and the yardstick}
\label{sec:cohort}

This primer fixes the cohort facts, the effect size used everywhere (lift), and what ``transductive'' means here. Every term is also in Appendix~A.

\begin{answerbox}
The cohort is 60,118 MRD CITE-seq cells from 13 patients, 1,569 of them malignant (2.61\%). Seven patients are \emph{supported} (at least 10 malignant cells and both classes): P01, P02, P03, P04, P06, P07, P09. P03, P02 and P04 are the highest-\emph{prevalence} patients (32.6\%, 14.4\%, 9.0\%). They hold 68.3\% of the malignant cells but only 8.2\% of all cells, whereas P09 and P01 alone are 48.2\% of the cells the label-blind Stage~1 is fit on. Results are reported as \emph{lift}: held-out AUPRC minus the patient's prevalence $\pi$, in percentage points (pp). We subtract $\pi$ because a constant score has AP exactly $\pi$. These numbers are counts and definitions, not estimates.
\end{answerbox}

\subsection*{What the cohort looks like}

The cells are the labelled (malignant or normal) CITE-seq cells from the MRD time point. After library-size normalization and log1p, genes detected in at least 1\% of cells were kept: 60,118 cells $\times$ 15,560 genes. All malignant cells are recipient-derived; the normal class is 58,499 donor-derived and 50 recipient-derived cells.

\begin{figure}[H]
  \centering
  \includegraphics[width=\linewidth]{figures/fig01a_cohort.pdf}
  \caption{\textbf{The cohort, patient by patient.} Dark: the seven supported patients, whose held-out lift is reported. Light gray: patients that stay in every training fold but are never summarized as held-out patients. Top: P09 and P01 dominate the cell count. Middle (log scale): malignant cells against the 10-cell support rule. Bottom: prevalence, which is also each patient's no-skill AP; note how far the bracketed patients sit above the rest.}
  \label{fig:01a}
\end{figure}

Figure~\ref{fig:01a} separates three rankings that the team discussion ran together. \textbf{By prevalence}, P03, P02 and P04 lead by a wide margin; the next patient is P09 at 2.2\%. \textbf{By malignant count}, the order is P03 (758), P09 (375), P02 (211), P04 (102), so ``the patients with the most malignant cells'' misdescribes P02 and P04. \textbf{By size}, P09 (17,126 cells) and P01 (11,860) are the largest patients, while P02, P03 and P04 are among the smallest (1,462, 2,327 and 1,128 cells).

This matters when we ask who ``shapes the factors''. Stage~1 is label-blind and weights every cell equally, so it sees cell shares, which P09 (28.5\%) and P01 (19.7\%) dominate. Stage~2 is supervised, so it sees malignant-cell shares: P02, P03 and P04 supply 1,071 of the 1,569 positives, and P03 alone 48.3\%. One consequence: weighting every patient equally in Stage~1 would \emph{raise} P02+P03+P04 from 8.2\% to 23.1\% of the fit (3 of 13 patients).

\subsection*{The yardstick: AUPRC, the no-skill floor, and lift}

Average precision (AP, our AUPRC) walks down a patient's cells ranked by the probe's score and asks how pure the list is each time it reaches another malignant cell. A ranking that knows nothing keeps that purity at the base rate, the prevalence $\pi$. So the no-skill floor differs by patient, from 0.38\% (P01) to 32.6\% (P03), and raw AP compared across patients mostly compares their prevalence.

\begin{explainbox}{Constant-score no-skill AP $=\pi$, exactly}
scikit-learn computes $\mathrm{AP}=\sum_n (R_n-R_{n-1})P_n$ over the distinct score thresholds. If every cell gets the same score, all cells form one tie group, and the only step goes from recall 0 to recall 1 at precision $N_+/(N_++N_-)=\pi$. Hence $\mathrm{AP}_{\text{const}}=\pi$ exactly; a random score only fluctuates around $\pi$. ``All healthy'' and ``all malignant'' both have AP $=\pi$. They differ only once a threshold is applied (P06: accuracy 0.9865 versus 0.0135).
\end{explainbox}

\begin{figure}[H]
  \centering
  \includegraphics[width=\linewidth]{figures/fig01b_yardstick.pdf}
  \caption{\textbf{Lift is the vertical gap between the probe's held-out AP (blue) and the patient's no-skill AP $=\pi$ (gray).} IFN-$\alpha$ PCA40, L1 at $C=10^{-4}$, transductive Stage~1, LOPO. P03 has the highest dot but not the longest gap. In P01 and P09 the dots coincide because the probe is all-zero. P07, the median patient here (11.08\pp), rests on 22 malignant cells (\emph{Remaining TODOs}, the inclusion decision).}
  \label{fig:01b}
\end{figure}

For a held-out patient $h$,
\[
  \lift_h=\mathrm{AP}_h-\pi_h \ \text{(pp)},\qquad
  \lift_h^{\mathrm{norm}}=\frac{\mathrm{AP}_h-\pi_h}{1-\pi_h}.
\]
Lift counts the AP points gained over a constant score, so 0 means no better than one. Its ceiling is $1-\pi_h$, which is 67.4\pp{} for P03 and 99.6\pp{} for P01. Normalized lift is the share of that headroom used; it is secondary. Per-patient lifts are summarized with every patient weighted equally, as the \textbf{median} and \textbf{P20} (20th percentile) over the seven supported patients. They are always quoted together with the cell-weighted \textbf{pooled LOPO AP} (AP over all held-out cells concatenated), because the three can disagree.

Figure~\ref{fig:01b} applies this to one stored probe (IFN-$\alpha$ PCA40, L1, $C=10^{-4}$). P03 has the highest AP (46.0\%), but its lift of 13.40\pp{} is below P04's (16.27\pp) and P02's (15.94\pp). Of P03's 46.0 AP points, 32.6 are prevalence. Normalized lift reorders the patients again (P03 0.199, P02 0.186, P04 0.179), because P03 has the least headroom. In P01 and P09 the L1 probe kept no factor at all, so the score is constant, AP $=\pi$, and lift is 0 \emph{by construction}. This is a regularization outcome, not evidence of no signal; the regularization section explains it.

\subsection*{The pipeline, and what ``transductive'' means}

\textbf{Stage~0} restricts the genes to a prior panel and z-scores them within it. The panel is either one of 41 curated programs (for example, Hallmark IFN-$\alpha$ response, 92 genes present) or one of 6 broad control bags (full\_34, core\_only and the HVG anchors). \textbf{Stage~1} reduces the panel to $K\in\{5,10,20,40\}$ factors with PCA, FA, FactoSig or FactoSig-promax (seed 42), or keeps the genes as they are (the \emph{direct-gene} representation). \textbf{Stage~2} fits a regularized logistic probe (L1, L2 or elastic net; balanced class weights; features standardized on the training cells). It is scored by \emph{LOPO}: each patient is held out in turn, the probe is trained on the other 12, and the held-out patient's cells are ranked. Each stored representation uses one of 65 settings ($C$ from $10^{-4}$ to $10^{4}$), chosen by pooled LOPO AP on the same folds, so the choice is not nested. Under balanced L1, a standardized factor enters only if its training class-mean gap exceeds $4/(C\,n_{\text{train}})$; that is where all-zero probes come from.

\begin{explainbox}{Transductive versus inductive Stage~1}
Stage~1 never sees labels, but the stored bases were fit on all 60,118 cells, \emph{including} those of the patient that Stage~2 later holds out. That patient's expression therefore helped define the axes. Such a basis is \emph{transductive}, and LOPO on it is label transfer on a shared atlas, not deployment to an unseen patient. An \emph{inductive} basis is refit without the held-out patient, whose cells are then projected onto it. This has been done for one panel only (IFN-$\alpha$; PCA and FA at $K=10,20,40$). There, lift stayed positive in 7/7 supported patients under tuned $C$ or L2. The per-patient median change was between $-1.29$ and $-0.05$\pp{} across 18 summaries, always negative: a small but consistent transductive optimism.
\end{explainbox}

\begin{explainbox}{How to read a 0 in the interactive heatmap}
A lift of exactly 0 is usually a threshold outcome, not a verdict. All 757 all-zero supported (cell, fold) pairs have $|\lift| \leq 5.6\times10^{-17}$, and 46 of the 783 atlas cells are all-zero in one to six supported folds (87 in all seven), so a count of patients above a lift cutoff partly counts all-zero folds. The cell view's per-fold strip shows each LOPO fold's nonzero-coefficient count at the cell's selected penalty and $C$ and flags all-zero supported folds (IFN-$\alpha$ FA40 is active in 3 of 13 folds, PCA40 in 11): open it before reading a patient pattern.

Treat \emph{knife-edge} cells as unresolved. A fold is knife-edge when its largest \emph{entry margin} (class-mean gap divided by the bar $4/(C\,n_{\text{train}})$) is within 10\% of 1, so a small change in $C$ flips the probe between all-zero and active. Among the 396 atlas cells with an L1 or elastic-net probe on a reduced basis, 136 of 2,772 supported folds in 49 cells are knife-edge (134 of them at the lowest $C$ on the grid), and crossing the bar changes lift by $\geq 5$\pp{} in 41 of 131 evaluable folds (up to 39.06\pp: hallmark hypoxia FA $K = 10$, P04 held out). The grid floor is common: 528 of 783 cells select it. The atlas picks $C$ by pooled LOPO AUPRC, which pools seven fold models and gives an all-zero fold one tied score, so it need not penalize all-zero folds: the fixed-$C$ L1 arms have the highest pooled AP of all IFN-$\alpha$ arms (PCA40 0.1630, FA40 0.1959) despite 2 and 5 all-zero folds. \src{2026-10-08\_d19-regularization-contract; 2026-10-08\_w7-atlas-strip-factor-explorer}
\end{explainbox}

\begin{explainbox}{Denominators used in this report}
\begin{itemize}[leftmargin=1.2em]
\item \textbf{783 atlas cells}: 47 panels $\times$ 17 Stage~1 configurations (16 DR = 4 methods $\times$ 4 $K$, plus direct genes), less 16 grid points where one of three small panels has fewer genes than $K$.
\item \textbf{640 biological DR cells}: the 681 biological atlas cells less their 41 direct-gene cells. Each panel has up to 16 DR cells; the three small panels have 8--12.
\item \textbf{643 representations} in the lift--prevalence correlations: the 640 biological DR cells plus the direct-gene probes of those three small panels.
\item \textbf{320 bases} in the factor-by-patient matrix: 41 programs $\times$ \{PCA, FA\} $\times$ $K \in \{5, 10, 20, 40\}$, less 8 grid points where a panel has fewer genes than $K$.
\end{itemize}
\end{explainbox}

\begin{caveatbox}
\begin{itemize}
  \item The normal class is 99.9\% donor-derived (58,499 of 58,549 cells), so malignant versus normal is almost recipient versus donor; genotype-linked genes could separate the classes without malignant biology. How the malignant labels were assigned may also differ by patient; a label-provenance audit is planned.
  \item Subtracting $\pi$ removes the floor, not the association. Across 643 representations, the median within-representation Spearman $\rho$ is 0.821 between lift and prevalence (0.964 for raw AP). These $\rho$ values are descriptive, not tests.
  \item Three supported patients rest on few positives: P07 (22 malignant cells), P06 (44) and P01 (45).
  \item Lift is a cell-level ranking quantity, not an improvement in recurrence risk.
\end{itemize}
\end{caveatbox}

\begin{teambox}
``Every result is scored as lift: held-out AUPRC minus the patient's malignant prevalence, which is exactly the AP a constant score gets. Prevalence runs from 0.38\% to 32.6\% across our seven scorable patients. P03, P02 and P04, the highest-prevalence patients, carry 68\% of the malignant cells but only 8\% of all cells.''

\textbf{Likely question:} ``Why not compare AUPRC directly, or use AUROC?''\\
\textbf{Answer:} Raw AP mostly ranks prevalence (median $\rho$ 0.964). Lift removes the floor but not the whole association ($\rho$ 0.821). AUROC has a 0.5 floor for everyone, yet it also ranks with prevalence (median $\rho$ 0.500). And because it ignores prevalence, it can be high while the top of the list is mostly normal cells. AUROC is planned as a companion column, not a replacement.

\textbf{Likely question:} ``Isn't malignant-vs-normal here just recipient-vs-donor?''\\
\textbf{Answer:} We cannot rule it out yet: 99.9\% of normal cells (58,499 of 58,549) are donor-derived. The planned label audit checks genotype-linked genes (HLA, ERAP2, chrY, XIST), sex match and the 50 recipient-normal cells; until then we say malignant-vs-normal transfer, not malignancy biology.
\end{teambox}

\subsection*{Next step}
\begin{itemize}
  \item Support- and prevalence-matched lift, to test whether the patient pattern survives equal prevalence (W1.5).
  \item Patient-level uncertainty and the AUROC companion column (W4.1).
  \item New MRD patients (W8.1). Runs on the new data leave out patients without malignant cells (your decision of 2026-10-09) and keep only patients with at least 100 malignant cells (your decision of 2026-10-10); the stored bases and the leverage refits of section~\ref{glance:02b} keep all 13 patients. \emph{Remaining TODOs}, subsection \emph{Proposed for the new data}, tabulates what the threshold keeps and what it does to median and P20 lift. The support rule for this report's summaries stays at 10 malignant cells.
\end{itemize}

\subsection*{How this was computed}
{\small\raggedright The counts come from \texttt{analysis/\allowbreak scorecards/\allowbreak stage2\_diagnostic/\allowbreak meeting\_followup\_v1/\allowbreak mrd\_patient\_cell\_composition.csv} (EXP-relative; D6, panel D6A). They match \texttt{scripts/\allowbreak mrd\_stage0\_2/\allowbreak research/\allowbreak study\_config.json}, and the shares are ratios of these counts. The worked example is the stored sharedness-LOPO result for IFN-$\alpha$ PCA40, L1, $C=10^{-4}$, on the transductive Stage~1 (\texttt{meeting\_followup\_v1/\allowbreak interferon\_alpha\_pca40\_fa40\_\allowbreak patient\_comparison.csv}). AP is \texttt{sklearn.metrics.\allowbreak average\_precision\_score}. C policy: the stored global selection by pooled LOPO AP, not nested; training $n$ unchanged. The Stage~2 runner is \texttt{scripts/\allowbreak mrd\_stage0\_2/\allowbreak stage2\_supervised/\allowbreak run\_stage2\_mrd\_multiobjective\_\allowbreak scorecard.py}. The figures come from \texttt{scripts/\allowbreak mrd\_stage0\_2/\allowbreak research/\allowbreak reports/\allowbreak 2026-10-08\_questions\_report/\allowbreak figs\_01\_cohort.py}. The inductive numbers are from D17 (\texttt{basis\_variant\_pilot\_ifna\_v1}). \src{D6, D8, D17; devlog 2026-08-31\_d6-d9-meeting-followup, 2026-10-08\_team-report-devlog-pivot, 2026-05-25\_set-a-old34-broad-screen, 2026-06-05\_set-a-all-biological-fullgrid, 2026-10-08\_d17-basis-variant-pilot}\par}
                      % 1. Cohort and yardstick (primer)
\labelnextsection{glance:02}\section{Do the highest-prevalence patients shape the factors?}
\label{sec:shape}

\begin{quote}
``We discussed quite a bit on the observation that patients with more malignant cells (P02, P03, P04) disproportionally shape the factors.''\\[2pt]
TODO: ``Re-weight cells per patient in Stage 1 so each patient contributes equally to DR \ldots{} then re-examine PCA and FA factors to see if the imbalance problem persists.''
\end{quote}

{\small\color{inkmuted}\emph{Reading guide.} What a leverage test is: the box ``The leverage test, step by step''; what it found on allograft rejection and apoptosis, the run you started (D26): section~\ref{glance:02b}. The metrics: the boxes ``Score energy, intuitively'' and ``The leverage test''. Figure~\ref{fig:02a} keeps cells, labels and score energy in separate panels; Figure~\ref{fig:02b} has its axis in percent and a reading of one row. How re-weighting works and what its percentages mean: the box ``How patient-equal re-weighting works''; your choice of target (2026-10-10): the options list. What the section implies for healthy-cell downsampling: the block after the re-weighting options.\par}

\begin{answerbox}
\textbf{Not more than other patients do, on the three panels tested (PCA and FA).}
\par\smallskip\noindent \emph{IFN-$\alpha$:} refitting Stage~1 without P02+P03+P04 leaves the malignant-vs-normal axis almost unchanged (under 1\% of its variance lost; matched $|r| \geq 0.988$) and not beyond random drops of the same number of cells and malignant cells (above all 20 in 3 of 6 method $\times$ $K$ cells; the rule needs 4, and on this panel it is applied post hoc).
\par\smallskip\noindent \emph{Allograft rejection and apoptosis} (D26, rule fixed before the run; section~\ref{glance:02b}): leaving them out moves the axis beyond those random drops (5 of 6 and 6 of 6 cells, mostly P03), but never more than leaving out other patients does; P01 moves it most. By the rule, the hypothesis that they shape the factors more than other patients do is not supported.
\par\smallskip\noindent \emph{Score space:} they hold about their 8.2\% cell share of factor variance on all 41 panels (median ratio 0.96; 1.04 after adjusting for cell type). \emph{Where they dominate:} Stage~2, with 68.3\% of the malignant labels. \emph{Not tested:} FactoSig bases (the atlas winners on allograft and apoptosis) and the other 38 panels.
\par\smallskip\noindent \emph{Decision it informed:} what re-weighting should target (W2.3). On the three panels tested (PCA and FA), the imbalance specific to these patients sits mainly in the Stage~2 labels: in the Stage~1 geometry it goes beyond composition on two panels (mostly P03) but, on the primary axis, never beyond other patients, and the largest Stage~1 leverage is P01's. You chose the Stage~2 labels on 2026-10-10, for the runs with the new patient data.
\end{answerbox}

\subsection*{What we did}

P03, P02 and P04 are the \emph{highest-prevalence} patients (32.6\%, 14.4\% and 9.0\% malignant), not the ones with the most malignant cells (P09 has more than P02 or P04); by total cells they are among the smallest. The stored Stage~1 bases are \emph{transductive}: fit once, without labels, on all 60,118 cells of the 13 patients, including any patient later held out and the four patients with no malignant cell at all. ``Shaping the factors'' can mean two things, and we measured both:
\begin{enumerate}[label=(\arabic*)]
\item \textbf{Where the variance sits.} On the fixed stored basis, how much of the factor-score variance sits in each patient's cells, compared with its share of the cells (``score energy'', next box; all 320 stored PCA/FA bases). D28 adds a cell-type adjustment: is a patient's excess just its cell-type mix?
\item \textbf{Leverage.} An in-silico perturbation: remove the patients from the Stage~1 fit, refit, and measure how far the malignant axis moves, against random removals of the same size and malignant composition (box ``The leverage test, step by step''; IFN-$\alpha$ here, allograft rejection and apoptosis in section~\ref{glance:02b}; PCA and FA at $K$ = 10, 20, 40).
\end{enumerate}
Only (2) tests whether a patient \emph{made} a factor; (1) is a cheap screen over all panels.

\begin{explainbox}{Score energy, intuitively}
\emph{Picture.} Each cell is a point in the $K$-dimensional space of its factor scores. Its \emph{score energy} is its squared distance from the cohort centre, $e_i = \sum_f (z_{if} - \bar z_f)^2$, on the stored (raw, not standardized) scores, so a high-variance factor counts more. Summed over all cells, the energy is $n$ times the total variance of the $K$ factor scores: for PCA exactly the quantity the factors are chosen to maximise.

\emph{Share and ratio.} A patient's \emph{energy share} is $\sum_{i \in p} e_i / \sum_i e_i$: the fraction of the captured variance that sits in its cells. Its \emph{ratio} is energy share $\div$ cell share, which equals the patient's mean energy per cell divided by the cohort's. A ratio of 1 means its cells are as spread out as an average cell; 1.11 means 11\% more squared distance per cell.

\emph{Worked example} (Figure~\ref{fig:02a}b). On IFN-$\alpha$ PCA40, P02+P03+P04 are 8.2\% of the cells and hold 9.1\% of the energy: $9.1 / 8.2 = 1.11$. On FA40 they hold 8.4\%: 1.02.

\emph{What it is not.} A share is an allocation on the fixed basis, not proof that those cells pulled the factors (that is the leverage test). It also ignores labels: the malignant-cell share in Figure~\ref{fig:02a}a is a different quantity, the Stage~2 label imbalance. With every factor standardized first (each counts equally), the 320-base median ratio of P02+P03+P04 is 1.00 instead of 0.96. Energy is summed over all of a patient's cells, malignant ones included; restricted to normal cells, the group's IFN-$\alpha$ ratio is still 1.11 (PCA40; FA40 1.03) and its 320-base median 0.99, so the excess is not carried by their malignant cells. \src{2026-10-09\_d28-celltype-score-shares}
\end{explainbox}

\subsection*{What we found}

\textbf{The labels are concentrated; the cells and the score energy are not} (Figure~\ref{fig:02a}).
Stage~1 sees cells, and the cell-heavy patients are P09 (28.5\%) and P01 (19.7\%). P02+P03+P04 are 4,917 of 60,118 cells (8.2\%) but 1,071 of 1,569 malignant cells (68.3\%; P03 alone 48.3\%). The four patients with no malignant cells (P08, P10, P11, P12) hold 17.7\% of the cells in the fit. On IFN-$\alpha$ at $K = 40$ the group's ratio is 1.11 (PCA) and 1.02 (FA); per patient, the most over-represented on this panel are P02 (1.43 PCA40, 1.29 FA40) and P07 (1.31, 1.44). PCA40's 1.11 is the third highest of the 320 bases, so IFN-$\alpha$ is a favourable case for the hypothesis and still near parity. Over all 320 bases the group's ratio has median 0.96 (range 0.58--1.14; above 1 in 75/320); per patient the medians are P01 1.24 (above 1 in 279/320), P02 1.06, P03 0.91 and P04 0.89, and the group's ratio creeps up with $K$ (PCA medians 0.92, 0.94, 0.98, 1.02 at $K$ = 5, 10, 20, 40) because its per-factor excess sits on low-variance trailing factors. \src{2026-10-08\_d16-factor-patient-matrix}

\begin{figure}[H]
\centering
\includegraphics[width=\linewidth]{figures/fig02a_who_dominates.pdf}
\caption{Two different imbalances, kept apart. (a) Each patient's share of the 60,118 cells (light; also its share of the unweighted Stage~1 fit) and of the 1,569 malignant cells (dark; its share of the Stage~2 positive labels). Dashed line: the share every patient would get under patient-equal weighting (1/13); over the P02+P03+P04 column, what the three get together (3/13). (b) Score-energy ratio on IFN-$\alpha$ at $K = 40$ (box ``Score energy, intuitively''), PCA40 blue, FA40 orange; a stem above 1 means the patient's cells hold more factor variance than their head count. The P02+P03+P04 values, 1.11 and 1.02, are labelled. (c) The same combined ratio on every stored PCA/FA basis; the ringed dots are IFN-$\alpha$ $K = 40$. Look at P03 in (a): 3.9\% of the cells, 48.3\% of the labels; then at its ratio in (b): 1.00 and 0.91.}
\label{fig:02a}
\end{figure}

\begin{explainbox}{The leverage test, step by step}
An in-silico perturbation, run on IFN-$\alpha$ (D17, this section) and on allograft rejection and apoptosis (D26, section~\ref{glance:02b}), for PCA and FA at $K$ = 10, 20 and 40:
\begin{enumerate}[label=\arabic*.,leftmargin=1.4em]
\item \emph{Reference axis.} On the stored basis, the factor with the largest standardized malignant--normal mean gap: PC8 (PCA) and FA7 (FA) at every $K$ on IFN-$\alpha$ (PC8 and FA8 for allograft rejection, PC6 and FA6 for apoptosis). Under balanced L1 it is the factor a sparse probe takes first (largest entry margin).
\item \emph{Perturb.} Refit Stage~1 without the patient set (genes re-standardized on the retained cells; same solver and seed).
\item \emph{Measure.} Over the retained cells, regress the original reference-axis scores on all $K$ refit factors (least squares with intercept). $R^2$ is the share of the axis's variance the refit basis still reproduces, and \textbf{axis loss} $= 1 - R^2$. Scale: 0 means fully recoverable; $10^{-3}$ means 0.1\% of the axis's variance is lost (best linear reconstruction correlates $R = \sqrt{1 - 10^{-3}} = 0.9995$ with the original); $10^{-2}$ means 1\% ($R = 0.995$). \textbf{Matched $|r|$} is the correlation with the single best-matching refit factor (one-to-one matching). Small axis loss with a low matched $|r|$ means the axis survives but is spread over several refit factors (rotated).
\item \emph{Compare.} Any refit on fewer cells moves the basis a little, so steps 2--3 are repeated for random drops. \emph{Size-matched} controls (3 draws) drop as many random cells. \emph{Composition-matched} controls (20 draws) drop as many cells \emph{and} as many malignant cells, at random from the whole cohort. Why both: the reference axis is the malignant-vs-normal direction, so removing malignant cells moves it whoever they belong to (dropping all 1,569 malignant cells moves the PCA40 axis 17.1$\times$ as much as dropping 1,569 random cells). A random draw of 4,917 cells holds about 128 malignant cells; P02+P03+P04 hold 1,071. Only the composition-matched draws separate ``these patients' cells'' from ``this many malignant cells''.
\item \emph{Summarise.} Ratio = observed axis loss $\div$ the mean of the 20 composition-matched controls; rank $p = (1 + \#\{\text{controls} \geq \text{observed}\})/21$, whose floor 0.048 means ``above all 20''. Descriptive and uncorrected.
\end{enumerate}
\emph{Reading one row} (Figure~\ref{fig:02b}a, P02+P03+P04 at PCA40): axis loss 0.0015, i.e.\ 0.15\% of PC8's variance (matched $|r|$ 0.995). Its 20 composition-matched controls lose 0.00054--0.00101 (mean 0.00079), so it is above all 20: rank $p$ 0.048, ratio 1.93. Against the size-matched mean (0.00030) it would have looked like 5.2$\times$.

\emph{Decision rule}, fixed for the allograft/apoptosis run (D26) before it ran: leverage beyond composition means above all 20 controls in at least 4 of the 6 method $\times$ $K$ cells; support for ``the highest-prevalence patients shape the factors'' also needs their ratio to exceed that of every contrast set (P01, P06, P09, P01+P09) in at least 4 cells. Both conditions were written after the IFN-$\alpha$ results below were known, so on IFN-$\alpha$ the rule is a post-hoc reading, not a pre-registered test; for allograft and apoptosis it is pre-specified.
\end{explainbox}

\textbf{On IFN-$\alpha$, leaving them out barely moves the axis; most of what moves is malignant composition} (Figure~\ref{fig:02b}).
Without P02+P03+P04 the IFN-$\alpha$ axis is essentially unchanged: PCA40 axis loss 0.0015 (matched $|r|$ 0.9953), FA40 0.00062 (0.9958); across method and $K$ the matched $|r|$ is 0.9885--0.9969 and the largest loss is 0.0069 (PCA10, 0.7\% of the axis). Size-matched controls had made this look like a 3.1--5.2$\times$ excess. Against composition-matched controls the FA excess disappears (1.00, 0.88, 0.36$\times$ at $K$ = 10, 20, 40; rank $p$ 0.48, 0.86, 0.19) and the PCA excess shrinks to 1.42--1.94$\times$ (rank $p$ 0.048 each). That is 3 of 6 cells above all controls, below the rule's 4, and in 0 of 6 cells is the group's ratio larger than every contrast set's. Panel~(c) shows why the second condition is needed: clearing composition-matched controls is common for any real patient (P01 and P09 in 6/6 cells, P02 and P06 in 5/6, P03 in 4/6, P04 in 2/6). The group clears fewer cells than P02 alone because its controls drop 1,071 random malignant cells, about two thirds of them from P02--P04 themselves, a stricter yardstick than a single patient's. A plausible reason: one patient's cells are a coherent block (one sample, one processing batch, one cell-type mix), while the controls are scattered across the cohort. P01 is the outlier (PCA 15.0--23.8$\times$, FA 3.74--21.0$\times$ its controls): without it the PCA40 axis loss is 0.0118 and PC8's best single match falls to $|r|$ 0.58, still in the subspace but rotated. \src{2026-10-08\_d17-basis-variant-pilot; 2026-10-08\_d17-composition-matched-controls}

\begin{figure}[H]
\centering
\includegraphics[width=\linewidth]{figures/fig02b_axis_leverage.pdf}
\caption{Leave-out refits of the IFN-$\alpha$ basis. (a, b) $K = 40$. Each row is one set removed from the Stage~1 fit (its cells and malignant cells on the left). Coloured dot: axis loss of PC8 (a) or FA7 (b) after refitting without that set. Bottom axis: axis loss as a percentage of the reference axis's variance that the refit factors cannot reproduce (log); top axis: the same value as $R = \sqrt{1 - \text{loss}}$. Gray ticks above the dot: the 3 size-matched random drops; dark dots below it: the 20 composition-matched random drops. $|r|$: matched correlation for the rows discussed in the text. (c) All $K$: axis loss divided by the mean of its 20 composition-matched controls (triangle, square, circle = $K$ 10, 20, 40); filled = above all 20 (rank $p$ 0.048), open = inside the control range; right column: cells above all 20, out of 6. Look at the P02+P03+P04 row: in (a) the blue dot sits just right of its control cloud (0.15\% vs 0.05--0.10\%); in (c) PCA is filled and FA is open, 3/6. For FA40, 2--10 of each set's 20 composition-matched drops fall in a second mode at 1.1--1.8\% (drawn as a separate segment), so FA40 dots there are within control noise.}
\label{fig:02b}
\end{figure}

\textbf{What lies behind P01? Its cell-type mix matches its score-energy excess, not its leverage} (Figure~\ref{fig:02c}).
D28 gives every cell the average energy of its predicted cell type (13 types) and adds these up per patient: the energy its cell-type mix alone predicts. The \emph{adjusted} ratio is observed $\div$ expected, an indirect standardization like a standardized mortality ratio; 1 means the mix fully accounts for the patient's ratio. Over the 320 bases, P01 goes from 1.24 (observed) to 1.21 (expected from its mix) and 1.03 (adjusted); the mix accounts for a median 90\% of its log departure from 1. The mechanism is in panel~(b): ncMono cells carry 3.25$\times$ the cohort's energy per cell, pDC2 1.26$\times$ and cDC 1.18$\times$, while HSC/MPP2 carry 0.65$\times$; P01's cells are 20.5\% pDC2, 14.2\% cDC and 8.4\% ncMono. P02+P03+P04 have a low-energy mix (16.3\% HSC/MPP2 at 0.65$\times$, 11.4\% MEP-EP and 10.6\% NP1 at 0.88$\times$ and 0.87$\times$, almost no ncMono or pDC2; their largest type, CLP at 29.1\%, sits at the cohort average): 0.96 observed, 0.93 expected, 1.04 adjusted (above 1 in 230/320). The IFN-$\alpha$ 1.11 is not the mix: PCA40 1.11 observed, 1.00 expected, 1.11 adjusted; FA40 1.02, 0.95, 1.08. Nor is P02's excess (1.06, 0.96, 1.11 over the 320 bases), or P07's on IFN-$\alpha$ (PCA40 1.31, 0.85, 1.54). The adjustment depends on the reference for P01, which holds 93\% of the cohort's ncMono cells: with the other patients' cells as the reference, P01's adjusted ratio is 1.08 and the mix accounts for 67\%.

``Accounts for'' is arithmetic here: the mix predicts most of P01's \emph{score energy}, which does not show that cell type causes it (the labels are transferred calls, and batch, sampling time and patient biology stay in the comparison). And it says nothing about P01's \emph{leverage}: on IFN-$\alpha$ PCA40, P01's ratio is 0.98 (adjusted 0.92), yet removing it moves PC8 15$\times$ as much as its composition-matched controls, the largest ratio of any set. One candidate, untested: leaving P01 out removes 93\% of the cohort's ncMono cells, the highest-energy cell type, almost a whole cell population. A leave-out control matched on cell type would test it. \src{2026-10-09\_d28-celltype-score-shares; 2026-10-08\_d17-composition-matched-controls}

\begin{figure}[H]
\centering
\includegraphics[width=\linewidth]{figures/fig02c_celltype_adjusted.pdf}
\caption{Is a patient's score-energy ratio its cell-type mix? (a) Open: the observed ratio (the quantity of Figure~\ref{fig:02a}b), here as the median (bar: interquartile range) over the 320 stored PCA/FA bases. Filled: the same ratio after the cell-type adjustment (observed $\div$ what the patient's predicted cell-type mix alone predicts). Right: the P02+P03+P04 values on IFN-$\alpha$ PCA40 (blue) and FA40 (orange). (b) Energy per cell of each predicted cell type relative to the cohort's (log; median and IQR over the 320 bases), with each type's percentage of P01's and of P02+P03+P04's cells underneath (bold: at least 10\%). Look at P01 in (a): its ratio falls from 1.24 to 1.03 once its pDC2, cDC and ncMono cells are accounted for; P02's rises from 1.06 to 1.11. Raw-score energy, cohort cell-type averages, all cells.}
\label{fig:02c}
\end{figure}

\textbf{Allograft rejection and apoptosis: beyond composition, mostly P03, but not beyond the other patients} (section~\ref{glance:02b}). The same test on the two panels where DR gains most over direct genes, with the rule fixed before the run (D26): without P02+P03+P04 the primary axis moves more than all 20 composition-matched drops in 5 of 6 cells (allograft) and 6 of 6 (apoptosis), at 1.31--4.15$\times$ and 2.01--3.98$\times$ their mean, mostly through P03; but the group's ratio beats every contrast set's in 0 of 6 primary-axis cells on each panel (1 of 6 on the apoptosis probe axis), and P01 moves the same axes 21--141$\times$ and 48--128$\times$ its controls. Section~\ref{glance:02b} has the design, the figures, the absolute sizes and what the result means for re-weighting and downsampling. \src{2026-10-10\_d26-leverage-allograft-apoptosis-done}

% Keep this heading with its explainer box: a new page only when fewer than twelve lines are left (the box
% does not start in a short remainder, which left the heading alone at a page foot). Same macro as report.tex;
% \providecommand keeps the standalone compile working.
\makeatletter
\providecommand{\softneedspace}[1]{\par\begingroup\dimen@ii\pagegoal\advance\dimen@ii-\pagetotal
  \ifdim #1>\dimen@ii\ifdim\dimen@ii>\z@\newpage\fi\fi\endgroup}
\makeatother
\softneedspace{12\baselineskip}
\textbf{Re-weighting (TODO option 1.2): what it changes, and the pilot.}

\begin{explainbox}{How patient-equal re-weighting works, and what its percentages mean}
\emph{Unweighted (the stored bases).} Every cell has weight 1, so a patient's influence on the Stage~1 fit is its cell share (Figure~\ref{fig:02a}a, light bars): P09 28.5\%, P01 19.7\%, P02+P03+P04 8.2\%.

\emph{Patient-equal.} Cell $i$ of patient $p$ gets weight $w_i = 1/(13\, n_p)$, so each patient's weights sum to $1/13 = 7.7\%$ (the dashed line). P02+P03+P04 go from 8.2\% to $3/13 = 23.1\%$ of the fit; P09 goes from 28.5\% to 7.7\%. \emph{These percentages are shares of the total fit weight}, not of cells or of labels.

\emph{Mechanics.} Genes are re-standardized with the weighted mean and SD. PCA takes the eigenvectors of the weighted covariance $\Sigma_w = \sum_i w_i (x_i - \mu_w)(x_i - \mu_w)^\top$. FA has no sample weights in scikit-learn, so it is fitted to pseudo-data whose covariance equals $\Sigma_w$ (the FA likelihood depends on the data only through the covariance). A second arm subsamples 1,128 cells per patient (the smallest patient's count; 14,664 cells; 5 seeds). All 60,118 cells are then projected, and Stage~2 runs unchanged on the standard LOPO folds ($C$ tuned per fold).

\emph{What it does and does not change.} It raises the highest-prevalence patients' influence on Stage~1 and leaves Stage~2 untouched: 68.3\% of the malignant labels still come from them.
\end{explainbox}

\emph{Pilot, IFN-$\alpha$ only (still transductive).} Patient-equal weighting rotates PC8 (matched $|r|$ 0.90 at $K = 40$) more than FA7 (0.98), with the axis still in the subspace ($R^2$ 0.990, 0.994). Under L1 with $C$ tuned per fold there is no clear loss (pooled AP / median / P20 lift): PCA40 0.1213 / 9.25 / 2.69\pp{} vs 0.1153 / 8.91 / 1.42\pp{} for the unweighted transductive refit with the same code (the stored basis gives 0.1001 / 8.71 / 1.37\pp; section \emph{Option 2} explains the difference), median per-patient change $+1.35$\pp, no patient lower; FA40 0.0863 / 6.85 / 1.20 vs 0.0894 / 7.03 / 1.25\pp, median paired change $+0.07$\pp{} but all three summaries slightly lower. Under tuned (and fixed) L2, FA40 loses a median 1.13\pp{} per patient, with 5 of 7 patients lower. \src{2026-10-08\_d17-basis-variant-pilot; 2026-10-08\_team-report-devlog-pivot}

\emph{Targeting Stage~1 geometry or Stage~2 labels: the options} (W2.3). \textbf{You chose the Stage~2 labels on 2026-10-10}, for the runs with the new patient data; the Stage~1 options stay listed for reference.
\begin{itemize}
\item \emph{Stage~1 geometry} (who shapes the factors): weights in the DR fit (patient-equal, as piloted, or weights by cell type), equal-$n$ subsampling, healthy-cell downsampling before the fit, or an inductive or normal-only fit. Stage~2 stays as is. Read-out: the leverage metrics above, then held-out lift.
\item \emph{Stage~2 labels} (who the probe learns from): per-patient sample weights in the logistic loss (each patient's malignant cells sum to the same weight, likewise its normal cells), class balancing within each patient instead of pooled, capping or downsampling one patient's training cells, or choosing $C$ by patient-equal median lift instead of pooled AP. The basis stays as is. Read-out: per-patient, median and P20 lift.
\item Either changes the effective training data, so $C$ is retuned per condition, as in the pilot.
\item The evidence behind the choice: Stage~1 leverage of these patients is not beyond composition on IFN-$\alpha$, and on allograft and apoptosis it is beyond composition but not beyond that of other patients (D26, section~\ref{glance:02b}), while the Stage~2 labels are 68.3\% from three patients and lift ranks with prevalence (median within-representation Spearman $\rho = 0.821$ over 643 representations). \src{D6 in 2026-08-31\_d6-d9-meeting-followup; 2026-10-10\_d26-leverage-allograft-apoptosis-done}
\end{itemize}

\textbf{What this section means for healthy-cell downsampling} (your idea; the plan waits for the new cohort).
\begin{itemize}
\item \emph{In Stage~1, downsampling healthy cells is a re-weighting of the fit.} In the unweighted fit a patient's weight is its cell count. P09 holds 28.5\% of the cells and about the same share of the IFN-$\alpha$ $K = 40$ score energy (28.8\% PCA40, 27.5\% FA40): its cells are not unusually spread out, they are many. That is a statement about shares, not leverage: removing P09 moves the IFN-$\alpha$ axis more than all 20 of its composition-matched controls in 6 of 6 method $\times$ $K$ cells (Figure~\ref{fig:02b}c). Thinning P09's healthy cells in Stage~1 is the experiment that would show how much of the geometry its cell count sets. The patient with the most healthy cells per malignant cell, P01 (about 263), has the largest leverage of all (section~\ref{glance:02b}).
\item \emph{Energy per cell differs by cell type} (Figure~\ref{fig:02c}b: ncMono 3.25$\times$, HSC/MPP2 0.65$\times$). A cell-type-stratified downsampling keeps the Stage~1 energy budget only if it keeps the type proportions; changing them re-weights the fit toward or away from these types.
\item \emph{On IFN-$\alpha$ at $K = 40$, a Stage~1 gain should be small.} The malignant axis is carried by normal-cell variation: a basis fitted on normal cells only keeps PC8 (matched $|r|$ 0.9937) and FA7 (0.9937), so the axis is already in the subspace (on allograft PCA it is not: section~\ref{glance:02b}). Where a Stage~1 gain can show is at small $K$ and on panels where the malignant direction is not a leading factor: fewer healthy cells give the malignant cells more of the fit's variance and move that direction to an earlier factor.
\item \emph{In Stage~2 the ratio as such should already be neutralized} (an argument, not a result). With \texttt{class\_weight='balanced'} each class carries half of the loss however many healthy cells there are, so thinning healthy cells at random is the same objective at a smaller $C$, plus noise; and AP only ranks the held-out cells. With $C$ retuned, a Stage~2 gain would have to come from \emph{which} healthy cells are kept (for example thinning abundant types more than rare ones), not from the ratio itself. Your earlier donor downsampling (D10, fixed $C$) is consistent with this, since L2 held-out AP barely moved, but it did not retune $C$, so it is not a test (section \emph{Remaining TODOs}).
\item \emph{So, to explain a gain} (your main concern), the design separates the channels: downsample in Stage~1 only, in Stage~2 only, and in both, with $C$ retuned per condition and the held-out patient untouched (its prevalence $\pi$, the lift's baseline, does not move). Section \emph{Remaining TODOs}, box ``Why fewer healthy cells could help, and why it might not'', lists the four channels and the read-outs for each.
\end{itemize}
\src{2026-10-09\_report-reading-notes; 2026-10-09\_w8-healthy-downsampling-plan; 2026-10-08\_d17-basis-variant-pilot; 2026-10-08\_d17-composition-matched-controls; 2026-10-09\_d28-celltype-score-shares}

\begin{caveatbox}
\begin{itemize}
\item \textbf{One panel for leverage here} (IFN-$\alpha$; allograft rejection and apoptosis in section~\ref{glance:02b}; PCA and FA only, no FactoSig); score shares cover 41 panels. Score share is not leverage, and both use transductive bases.
\item \textbf{No significance}: 20 controls (rank-$p$ floor 0.048), 48 uncorrected contrasts per panel, 7 supported patients. On IFN-$\alpha$ the decision rule is a post-hoc reading (it was written for D26 after these results). A leave-out refit is an in-silico perturbation: it shows how much a patient's cells move an axis, not why. FA40 controls are bimodal, so FA40 ratios below 1 mean ``within control range'', not ``less leverage''.
\item \textbf{The P02+P03+P04 control is not free of them}: they supply 68.3\% of malignant cells, so the matched draws remove many of theirs too. The test isolates what is specific to their cells beyond the label count.
\item \textbf{The cell-type adjustment uses predicted labels} (label-transfer calls) and removes one confound; batch, sampling time and patient biology stay in the adjusted ratio. For P01 the reference matters (adjusted 1.03 or 1.08).
\item \textbf{P01's leverage is not explained}, and P02--P04 still dominate Stage~2 labels, selection and lift summaries.
\end{itemize}
\end{caveatbox}

\begin{teambox}
``We tested whether P02, P03 and P04 shape the factors in two ways. Their cells hold about their share of the factor variance on all 41 panels (1.04$\times$ after adjusting for cell type). On IFN-$\alpha$, refitting Stage~1 without them changes the malignant axis by under 1\% of its variance and not clearly more than removing the same number of random cells and malignant cells. On allograft rejection and apoptosis it moves 1.3 to 4 times as much as such random removals, mostly because of P03, but still by under 3\% of the axis's variance, and less than removing other patients, P01 above all. Where they clearly dominate is Stage~2: 68\% of the malignant labels.''

\textbf{Likely question:} ``So should we re-weight Stage~1 so every patient counts equally?''\\
\textbf{Answer:} Not as a fix for this: it would raise P02+P03+P04 from 8.2\% to 23.1\% of the fit (the pilot rotates PC8 to $|r|$ 0.90; no clear loss under tuned L1, but FA40 loses a median 1.13\pp{} under L2). If the worry is that three patients drive the results, the data place that mainly in Stage~2, where they supply 68\% of the labels, so re-weighting will target the Stage~2 labels in the runs with the new patients. The patient that moves the primary axis most is P01 (19.7\% of cells, 45 malignant), in 17 of the 18 method $\times$ $K$ cells of the three panels; its excess score energy is mostly matched by its cell-type mix, but its leverage is still unexplained.
\end{teambox}

\subsection*{Next step}
\begin{itemize}
\item \textbf{Done:} the allograft/apoptosis leverage test (D26, W2.4), which you ran in screen W2\_d26; verdict for PCA and FA: not supported. Section~\ref{glance:02b}, with its own next steps (W2.8 all panels, W2.9 P01, W2.2 FactoSig).
\item \textbf{Done:} the cell-type stratification of the score shares (D28, W1.9; Figure~\ref{fig:02c}).
\item Proposed: locating P02's and P07's excess that cell type does not explain, by factor and cell type (W1.10).
\item Re-weighting (W2.3): target decided on 2026-10-10 (the Stage~2 labels). The arm is designed with the runs on the new patient data, with $C$ retuned per condition; it runs inside the healthy-cell downsampling design as its Stage-2-only sample-weight scheme. Those runs leave out patients with no malignant cells and keep only patients with at least 100 malignant cells (W8.5; section \emph{Remaining TODOs}, subsection \emph{Proposed for the new data}). \src{2026-10-10\_decisions-new-run-inclusion-reweighting-lfs}
\end{itemize}

\subsection*{How this was computed}
{\small\raggedright
60,118 MRD CITE-seq cells (13 patients, 1,569 malignant); stored transductive Stage~1 bases, seed 42. \emph{Score shares (D16):} 41 biological panels $\times$ \{PCA, FA\} $\times$ $K \in \{5,10,20,40\}$, less 8 grid points where a panel has fewer genes than $K$: 320 bases with 5,870 factors (5,838 non-degenerate); energy on raw cohort-centred scores; no Stage~2 model, so no $C$ policy. \emph{Cell-type adjustment (D28):} the same 320 bases; 13 predicted cell types (\texttt{predicted.annotation}); expected energy = each cell's cell-type average over the cohort; adjusted ratio = observed / expected; nothing refit, no $C$ policy. \emph{Leverage (D17):} IFN-$\alpha$ (92 covered genes), PCA and FA at $K$ = 10, 20, 40, refit with the stored sklearn randomized solvers and genes standardized on the fit cells; reference = stored FA scores and, for PCA, the projection onto the stored loadings; size-matched controls seeds 1--3, composition-matched seeds 101--120; alignment only, no $C$ policy. \emph{D26} (allograft rejection, apoptosis): section~\ref{glance:02b}. \emph{Balanced pilot:} patient-equal weights over 13 patients; Stage~2 L1/L2 with $C$ tuned per outer fold by inner LOPO on the standard LOPO folds (training $n$ unchanged). Code: \texttt{scripts/\allowbreak mrd\_stage0\_2/\allowbreak workstreams/\allowbreak W1\_patient\_aggregation/\allowbreak run\_d16\_factor\_patient\_matrix.py}, \texttt{\ldots/\allowbreak W1\_patient\_aggregation/\allowbreak run\_d28\_celltype\_score\_shares.py}, \texttt{\ldots/\allowbreak W2\_stage1\_basis\_variants/\allowbreak run\_d17\_basis\_variant\_pilot.py}; figures and quoted ratios from this report's \texttt{figs\_02\_shape.py}. Artifacts (EXP-relative, under \texttt{analysis/\allowbreak scorecards/\allowbreak stage2\_diagnostic/}): \texttt{factor\_patient\_matrix\_v1/}, \texttt{factor\_patient\_matrix\_celltype\_v1/}, \texttt{basis\_variant\_pilot\_ifna\_v1/}. \src{2026-10-08\_d16-factor-patient-matrix; 2026-10-09\_d28-celltype-score-shares; 2026-10-08\_d17-basis-variant-pilot; 2026-10-08\_d17-composition-matched-controls; D6 in 2026-08-31\_d6-d9-meeting-followup; 2026-10-08\_team-report-devlog-pivot}\par}
 % 2. Do P02/P03/P04 shape the factors? (D16, D17, D28)
\labelnextsectionb{glance:02b}\section{What did we do with allograft rejection and apoptosis?}
\label{sec:d26}

\begin{quote}
``What did we do with allograft and apoptosis? Please update all that as a new section in the report.'' (2026-10-10)\\[2pt]
``Why would it be open until allograft and apoptosis are tested? How about other stage 0 panels? Test what? Should we test them too?'' \ldots\ ``I like the next step, let's start them.'' (reading notes on version~1, A-20261009-02 and -15)
\end{quote}

\begin{answerbox}
\textbf{We ran the IFN-$\alpha$ leverage test on both panels; by the rule fixed before the run, the highest-prevalence patients do not shape these factors more than other patients do.}
\par\smallskip\noindent \emph{What we did} (D26, your screen \texttt{W2\_d26}, 10 October): refit the allograft-rejection and apoptosis Stage~1 bases without P02+P03+P04, without each of them, and without four contrast sets (P01, P06, P09, P01+P09), and compared how far the malignant--normal axis moves with 20 random drops of the same number of cells \emph{and} malignant cells. PCA and FA at $K$ = 10, 20, 40; 2,364 refits.
\par\smallskip\noindent \emph{What we found:} leaving P02+P03+P04 out moves the axis more than all 20 of those drops in 5 of 6 method $\times$ $K$ cells (allograft) and 6 of 6 (apoptosis), at 1.3--4.1$\times$ and 2.0--4.0$\times$ their mean, mostly through P03. The axis stays matched ($|r| \geq 0.945$ and $\geq 0.994$). Leaving P01 out moves the same axes 21--141$\times$ and 48--128$\times$ its own controls, and in no primary-axis cell on any of the three tested panels does the group's excess beat every contrast set's (0 of 6 each; 1 of 6 on the apoptosis probe axis). Hypothesis H-S1-05 is therefore not supported for PCA and FA.
\par\smallskip\noindent \emph{What it means:} the imbalance specific to P02--P04 sits mainly in the Stage~2 labels (68.3\% of the malignant cells), which fits your decision that re-weighting targets the Stage~2 labels. The largest Stage~1 leverage belongs to P01, the supported patient with the most healthy cells per malignant cell. \emph{Not tested:} FactoSig and FactoSig-promax (the atlas winners on both panels) and the other 38 panels. Descriptive.
\end{answerbox}

\subsection*{Why these two panels}

Allograft rejection and apoptosis are the two programs where DR gains most over direct genes: median lift 12.56 vs 0.23\pp{} (allograft) and 11.52 vs 0.005\pp{} (apoptosis). For allograft that gain sits in the highest-prevalence patients: selected DR minus direct genes is $+46.37$\pp{} in P03 and $+0.09$\pp{} in P01 (section \emph{The three scientific questions}). Their largest-gap factor also separates mostly those patients' cells: allograft FA8 at $K = 40$ has univariate held-out lift 39.04\pp{} in P03 and 22.68\pp{} in P04, but under 1\pp{} in P01, P06 and P07 (labels, not loadings). If P02--P04 shaped these factors, the DR advantage could be their fingerprint. On IFN-$\alpha$ the same test had found no leverage beyond malignant composition (section~\ref{glance:02}), so these two panels decided H-S1-05. \src{F2/F4, D1DD5; 2026-09-17\_team-feedback-f1-f7; 2026-09-10\_d1a-descriptive-deep-dives; 2026-10-08\_d17-composition-matched-controls}

\subsection*{What we did}

The design is D17's, unchanged (box ``The leverage test, step by step'' in section~\ref{glance:02}):
\begin{enumerate}[label=\arabic*.,leftmargin=1.4em]
\item \emph{Reference axis} on the stored transductive basis. Primary: the factor with the largest standardized malignant--normal gap over all cells (PC8 and FA8 for allograft rejection, PC6 and FA6 for apoptosis, at every $K$). Two secondary axes: the factor with the best median univariate held-out lift (a different factor from the primary in 8 of the 12 method $\times$ $K$ cells), and the atlas-selected Stage~2 probe refit on all cells, a direction made of 8--40 factors.
\item \emph{Perturb.} Refit Stage~1 without one set: P02, P03, P04 or P02+P03+P04, or one of the contrast sets P01, P06, P09 and P01+P09 (genes re-standardized on the retained cells; same solver and seed).
\item \emph{Measure} over the retained cells: \emph{axis loss} $= 1 - R^2$, the share of the reference axis's variance that all $K$ refit factors together can no longer reproduce, and \emph{matched $|r|$}, its correlation with the single best refit factor.
\item \emph{Compare} with 3 size-matched drops (as many random cells) and 20 composition-matched drops (as many random cells \emph{and} as many malignant cells, seeds 101--120). The ratio is the axis loss divided by the mean of the 20; rank $p$ 0.048 means above all 20.
\item \emph{Sanity floor}: a refit on all cells reproduces the stored basis (36 of 36 checks pass); a refit without all 1,569 malignant cells shows how much of the axis normal cells carry on their own.
\end{enumerate}
PCA and FA at $K$ = 10, 20 and 40; 2,364 refits in 8 minutes on 12 CPU workers. FactoSig and FactoSig-promax, the atlas winners on both panels, could not be refit: the runner is CPU-only and their refit needs the GPU code (W2.2). \src{2026-10-09\_d26-leverage-allograft-apoptosis; 2026-10-10\_d26-leverage-allograft-apoptosis-done}

\begin{explainbox}{The decision rule, fixed before the run}
Written into the run entry on 2026-10-09, before the screen started; ``cells'' are the 6 method $\times$ $K$ combinations of one panel and axis.
\begin{enumerate}[label=(\arabic*),leftmargin=1.6em]
\item \emph{Leverage beyond composition}: P02+P03+P04 lose more of the axis than all 20 composition-matched drops in at least 4 of 6 cells.
\item \emph{Support for H-S1-05} (``they shape the factors'') also needs their ratio to exceed that of every contrast set (P01, P06, P09, P01+P09) in at least 4 of 6 cells. Without this condition any real patient would ``shape'' the factors: on IFN-$\alpha$, P01 and P09 cleared their controls in 6 of 6 cells.
\item A verdict change needs the primary and the probe axis to agree.
\item Absolute size is reported regardless: matched $|r| \geq 0.98$ reads as ``beyond controls but small''.
\end{enumerate}
On IFN-$\alpha$ the same rule is a post-hoc reading (it was written after those results); here it is pre-specified.
\end{explainbox}

\subsection*{What we found}

\begin{table}[H]
\centering
\small
\caption{The decision rule per panel and axis (6 method $\times$ $K$ cells each). Ratio: axis loss of P02+P03+P04 divided by the mean of its 20 composition-matched drops. IFN-$\alpha$ is the D17 refit, primary axis only.}
\label{tab:d26rule}
\setlength{\tabcolsep}{4.5pt}
\begin{tabular}{l l c c c c}
\toprule
Panel & Axis & (1) above all 20 & (2) above every contrast & Ratio & P01's ratio \\
\midrule
IFN-$\alpha$ (D17) & primary & 3/6 & 0/6 & 0.36--1.94 & 3.7--23.8 \\
\addlinespace
Allograft rejection & primary & 5/6 & 0/6 & 1.31--4.15 & 21.2--140.6 \\
 & best held-out & 5/6 & 0/6 & 0.82--3.99 & 16.0--188.1 \\
 & probe & 5/6 & 0/6 & 1.22--3.24 & 7.5--35.1 \\
\addlinespace
Apoptosis & primary & 6/6 & 0/6 & 2.01--3.98 & 48.1--127.7 \\
 & best held-out & 6/6 & 0/6 & 1.76--3.98 & 30.7--126.2 \\
 & probe & 6/6 & 1/6 & 1.54--4.37 & 0.57--91.9 \\
\bottomrule
\end{tabular}
\end{table}

\textbf{Beyond composition, yes, and mostly P03} (Figure~\ref{fig:02b_d26a}a--c). Without P02+P03+P04 the primary axis moves more than all 20 composition-matched drops in 5 of 6 cells for allograft and 6 of 6 for apoptosis (IFN-$\alpha$: 3 of 6), at 1.31--4.15$\times$ and 2.01--3.98$\times$ the controls' mean; the held-out-lift and probe axes agree (Table~\ref{tab:d26rule}). P03 alone accounts for 74--93\% (allograft) and 72--88\% (apoptosis) of what the three together lose, against 45--69\% on IFN-$\alpha$, and clears its own controls in 5 of 6 and 6 of 6 cells; P02 alone in 4 of 6 and 6 of 6, P04 in 3 of 6 and 1 of 6. A union can show a smaller ratio than one of its members (allograft PCA40: P03 10.4$\times$, the three together 4.15$\times$) because its controls remove more malignant cells (1,071, against P03's 758), which makes its yardstick stricter; P02 and P04 do not cancel P03. The per-cell ratios behind these counts are printed by \texttt{figs\_02b\_d26.py} and tabulated in the run's findings file.

\textbf{Small in absolute terms} (Figure~\ref{fig:02b_d26a}d--f). Allograft PCA loses 0.9--2.8\% of its PC8's variance (the allograft panel's own largest-gap factor, not the IFN-$\alpha$ PC8 of section~\ref{glance:05}), with matched $|r|$ 0.945--0.970: by rule (4) that is not ``small''. Yet the 20 random drops of 4,917 cells, 1,071 of them malignant, also rotate PC8 to $|r|$ 0.90--0.97, so on matched $|r|$ the group sits inside its controls' range at $K$ 10 and 40 and just above it at $K$ 20 (0.9704 against a control maximum of 0.9696, i.e.\ it rotates PC8 slightly less than every control); and a refit without any malignant cells keeps PC8 at only $|r|$ 0.88--0.92 (IFN-$\alpha$ 0.99). The allograft PCA axis is partly made by the malignant cells themselves, whoever they belong to. Allograft FA loses 0.4--0.7\% (matched $|r|$ 0.988--0.991), apoptosis 0.07--0.18\% (0.994--0.995): ``beyond controls but small''.

\textbf{Not more than other patients: P01 dominates.} Clearing composition-matched drops is common to every real patient set, as on IFN-$\alpha$: P01 and P01+P09 do it in 6 of 6 cells on both panels, P09 and P06 in 4 of 6 (allograft) and 6 of 6 (apoptosis). On the primary axis P01 has the largest ratio of all eight sets in every allograft and apoptosis cell, and in 17 of the 18 cells of the three panels (the exception is IFN-$\alpha$ FA40, where P06 leads): 21.2--140.6$\times$ its controls on allograft and 48.1--127.7$\times$ on apoptosis, an axis loss of 1.1--9.0\% and 1.1--3.7\%. So the group's ratio beats every contrast set's in 0 of 6 primary-axis cells on each panel (1 of 6 on the apoptosis probe axis): condition (2) fails on both axes the rule names. The verdict does not hinge on P01: against P06 and P09 alone, the group's ratio is higher only in the three allograft PCA cells. Without P01 the apoptosis PCA axis stays in the refit span (1.6--2.6\% lost) but no single refit factor matches PC6 (matched $|r| \leq 0.043$): it is spread over several factors, as PC8 was on IFN-$\alpha$ (FA6 keeps 0.735--0.756; allograft PC8 0.667--0.874).

\textbf{The probe direction moves under any drop.} The dense $K = 40$ probes (up to 40 factors) lose a lot whoever is removed: allograft FA40 11.7\% without P02+P03+P04 and 9.6\% on average over its 20 random drops; apoptosis FA40 18.4\% against a control mean of 4.2\% (maximum 16.8\%).

\begin{figure}[H]
\centering
\includegraphics[width=\linewidth]{figures/fig02b_d26a_ratio_by_set.pdf}
\caption{Leave-out refits, primary axis, three panels (IFN-$\alpha$ from D17). (a--c) Each column is one patient set left out of the Stage~1 fit (02 = P02; 02+03+04 = P02+P03+P04); each marker is one method $\times$ $K$ (blue PCA, orange FA; triangle, square, circle = $K$ 10, 20, 40): its axis loss divided by the mean of its 20 composition-matched random drops (log). Filled: above all 20 (rank $p$ 0.048); open: inside their range. Above each column: cells above all 20, out of 6. Bold ranges: the P02+P03+P04 ratio (bottom) and P01's (shaded column, top). (d--f) Absolute size: matched $|r|$ of the primary axis after leaving P02+P03+P04 out (dot), the range over its 20 random drops (gray bar) and after leaving P01 out (diamond), on a log scale of $1 - |r|$ with $|r| = 1$ at the top; dashed: $|r| = 0.98$, the rule's ``small''. Look at the shaded P01 column against the P02+P03+P04 column in (a--c), and at the diamonds in (f).}
\label{fig:02b_d26a}
\end{figure}

\subsection*{What it means}

\textbf{For re-weighting.} You decided on 2026-10-10 that re-weighting should target the Stage~2 labels, for the runs with the new patient data (A-20261010-08). D26 is consistent with that choice: the Stage~1 leverage specific to P02--P04 is beyond composition on two of three panels but never beyond what removing other patients does, while the Stage~2 labels are 68.3\% theirs and lift ranks with prevalence (median within-representation Spearman $\rho = 0.821$ over 643 representations). As recorded, the re-weighting gives each training patient's malignant cells the same total weight in the logistic loss, and likewise its normal cells; it runs as the Stage-2-only sample-weight scheme of the healthy-cell downsampling design, with $C$ retuned per condition, since the effective training data change. No patient-reweighted Stage~1 refit is planned; the IFN-$\alpha$ pilot (section~\ref{glance:02}) stays the record of what Stage~1 re-weighting does. Re-weighting the labels leaves the basis as it is, including P01's leverage. \src{D6 in 2026-08-31\_d6-d9-meeting-followup; 2026-10-10\_d26-leverage-allograft-apoptosis-done; 2026-10-10\_decisions-new-run-inclusion-reweighting-lfs}

\textbf{For healthy-cell downsampling} (plan \texttt{healthy\_downsampling\_plan.md}, section \emph{Remaining TODOs}). Two points carry over, both predictions, not results:
\begin{itemize}[leftmargin=1.4em]
\item \emph{P01 is where Stage~1 thinning would bite.} P01 is the supported patient with the most healthy cells per malignant cell (about 263; 11,860 cells, 45 malignant) and the largest leverage on the primary axis of all three tested panels. The plan's patient-capped scheme thins P01 and P09 most, so in Stage~1 it should move the basis more than random thinning of the same number of cells. (Under your 100-cell rule P01 is not in new full runs at all; this applies to the current cohort's bases and to any new patient like it.) Whether that is P01's cell count, its ratio or its cell types (it holds 93\% of the cohort's ncMono cells, the highest-energy type) is untested (W2.9).
\item \emph{The Stage~1 channel is panel-specific.} On apoptosis and on IFN-$\alpha$ the axis survives without any malignant cells (matched $|r| \geq 0.99$, except IFN-$\alpha$ FA20 at 0.95), so changing the healthy-to-malignant mix of the fit set should change little there. On allograft it does not: PCA $|r|$ 0.88--0.92, FA 0.97--0.98, both under the rule's 0.98. Allograft is the panel where the plan's Stage~1 geometry channel has the most room, PCA most of all. These are axis movements, not held-out lift (D26 refits no Stage~2 probe).
\end{itemize}

\begin{caveatbox}
\begin{itemize}
\item \textbf{PCA and FA only}, on 3 of 41 panels. FactoSig and FactoSig-promax, the atlas winners on both panels, are untested; they are the only remaining route to a different H-S1-05 reading here (W2.2).
\item \textbf{Axis movement is not lift.} No Stage~2 probe was refit, so D26 says how far a patient set moves the factors, not how that changes held-out AUPRC.
\item \textbf{No significance}: 20 controls (rank-$p$ floor 0.048), 48 uncorrected contrasts per panel and axis, 7 supported patients. A leave-out refit is an in-silico perturbation: it shows how much a set's cells move an axis, not why.
\item \textbf{The group's controls are not free of the group}: they remove 1,071 random malignant cells, about two thirds of them from P02--P04 themselves. The test isolates what is specific to their cells beyond the label count.
\item \textbf{Which axis.} The primary axis is the largest-gap factor; the best held-out factor is a different one in 8 of 12 cells, and the dense probe direction moves under any drop. The verdict holds on all three axes, but the absolute sizes are those of the primary axis.
\item \textbf{Transductive, 13-patient bases.} The test refits variants of the stored bases, in which the four zero-malignant patients are 17.7\% of the fit. Under the new-run rule (no zero-malignant patients) the bases and their controls change; the test on that 9-patient fit set has not been run.
\end{itemize}
\end{caveatbox}

\begin{teambox}
``On allograft rejection and apoptosis, the two programs where DR helps most, removing P02, P03 and P04 from the factor fit moves the malignant axis 1.3 to 4 times as much as removing the same number of random cells and malignant cells, mostly because of P03, but the axis stays matched (under 3\% of its variance), and removing P01 moves it 21 to 141 times as much. So they do not shape the factors more than other patients do; where they dominate is the labels.''

\textbf{Likely question:} ``Then why does P01 move the axes so much?''\\
\textbf{Answer:} We don't know yet. It is the second-largest patient (19.7\% of the cells) and holds 93\% of the cohort's ncMono cells, the cell type with the most factor-score energy per cell; a leave-out control matched on cell type would test that (W2.9). Its malignant cells cannot be the reason: it has 45.
\end{teambox}

\subsection*{Next step}
\begin{itemize}
\item Proposed: the same test on all biological panels (W2.8; about 3 hours in screen; needs a campaigns row and your go-ahead). Its condition, leverage beyond composition, is met; condition (2) against the contrast sets is the informative read-out, since every real patient set clears composition-matched drops.
\item Proposed: a cell-type-matched leave-out control for P01 (W2.9), now the main open leverage question.
\item FactoSig and FactoSig-promax refits (W2.2, GPU), the only route to a different verdict on these panels.
\item Re-weighting (W2.3, closed by your decision): the Stage~2 sample-weight scheme runs with the new-patient runs (W8.6 runner, W8.7 run), with $C$ retuned per condition.
\end{itemize}

\subsection*{How this was computed}
{\small\raggedright
Stored transductive Stage~1 bases (60,118 cells, 13 patients, seed 42) of hallmark allograft rejection (147 genes) and hallmark apoptosis; refits with the stored scikit-learn solvers, genes standardized on the fit cells; size-matched controls seeds 1--3, composition-matched seeds 101--120 (the D17 seeds); all 13 patients kept, as in D17, so the panels compare directly. Alignment only: no Stage~2 probe refit or scored, so no $C$ policy; training $n$ unchanged. \par\smallskip Run, from the repository root:\\
{\footnotesize\hspace*{1em}\texttt{screen -S W2\_d26}\\
\hspace*{1em}\mbox{\texttt{bash scripts/mrd\_stage0\_2/workstreams/W2\_stage1\_basis\_variants/run\_d26\_leverage\_panels.sh}}\\}
(defaults \texttt{\mbox{-{}-ks 10,20,40} \mbox{-{}-resume} \mbox{-{}-n-jobs 12}}, single-thread BLAS); finished 2026-10-10.\par\smallskip Code: \texttt{\ldots/\allowbreak W2\_stage1\_basis\_variants/\allowbreak run\_d26\_leverage\_panels.py}; figures, the decision-rule counts (recomputed from \texttt{d26\_leverage\_contrast.csv} and asserted equal to \texttt{d26\_decision\_summary.csv}) and every quoted ratio from this report's \texttt{figs\_02b\_d26.py}. Artifacts (EXP-relative): \texttt{analysis/\allowbreak scorecards/\allowbreak stage2\_diagnostic/\allowbreak leverage\_allograft\_apoptosis\_v1/} (\texttt{d26\_leverage\_contrast.csv}, \texttt{d26\_decision\_summary.csv}, \texttt{d26\_reference\_axes.csv}, findings), with the run's own figures under \texttt{analysis/\allowbreak figures/\allowbreak stage2\_diagnostic/\allowbreak leverage\_allograft\_apoptosis\_v1/}; IFN-$\alpha$ rows from \texttt{basis\_variant\_pilot\_ifna\_v1/\allowbreak d17\_leverage\_contrast.csv}. \src{2026-10-09\_d26-leverage-allograft-apoptosis (run entry, decision rule); 2026-10-10\_d26-leverage-allograft-apoptosis-done (results, verdict); 2026-10-10\_core-report-v3 (the ranges this section derives); 2026-10-08\_d17-composition-matched-controls; 2026-10-09\_d28-celltype-score-shares}\par}
\plainsectionnumbers % 2b. Allograft rejection and apoptosis: the leverage test (D26)
\labelnextsection{glance:03}\section{Should we aggregate patients on one Stage 1 basis?}
\label{sec:03_aggregate}

\begin{quote}
``Data driven statistical test for whether to aggregate patients: 1. Learn factors on all patients, project each patient's cells onto those factors. 2. Check the per-patient factor distributions. 3. If distributions overlap significantly: patients share variation; if distinct: treat as separate subtypes.''
\end{quote}

\begin{answerbox}
Yes, on the current Stage~1 bases. \emph{Aggregate} has a fixed, quantitative meaning here: pooling is supported when, averaged over factors with variance weights, the median pair of patients shares at least 75\% of its normal-cell factor-score histogram and at least 45\% of its malignant-cell histogram, no patient's normal cells share less than 60\% with all the other patients' normal cells, and the malignant overlap sits less than 0.15 below that of label-shuffled patients. Of 47 representations (41 biological panels, 6 broad gene-bag controls), 45 meet all four gates, 2 are location shifts and none is a subtype candidate. On biological panels the median pair shares 0.777--0.890 of its normal-cell histogram (median 0.843) and 0.538--0.657 of its malignant-cell histogram (median 0.610). The only patient-private axis is P03's; it appears in three broad gene bags and in none of the 41 biological representations tested (one per panel). The test asks whether patients' cells \emph{land in the same score range}, not how much each patient \emph{contributes} to a factor. Confidence: descriptive, transductive basis, seven supported patients, uncalibrated gates.
\end{answerbox}

\subsection*{What we did}

\textbf{What the test asks, and what it does not.} Hold the Stage~1 factors fixed. For each patient, look at where that patient's cells fall on each factor, and ask whether that score distribution overlaps the other patients' distributions. If it does for every patient, one basis describes everybody and pooling the patients is reasonable. If one patient's cells sit in a range nobody else occupies, that patient behaves like a separate subtype on this basis. The test does \emph{not} measure how much each patient contributes to a factor: a patient can overlap the others completely whether or not its cells steered the loadings. Contribution is the subject of the section on whether the highest-prevalence patients shape the factors (score-energy shares and leave-out refits).

We ran the three steps as written, without refitting anything. Step~1 is how the pipeline already works. Every Stage~1 basis is \emph{transductive}: it was fit on all 60,118 cells of all 13 patients, including any patient later held out in LOPO. We took the stored scores of the 47 representations shown in the interactive atlas: one per Stage~0 panel, each at the method and $K$ with the highest patient-median lift.

\textbf{The method on one factor} (Figure~\ref{fig:03w1_one_factor}). Normal and malignant cells are never mixed. Each label is its own \emph{stratum}, with its own patients, bin grid, pairs and null: normal cells come from the 13 patients with $\geq 200$ normal cells (58,549 cells, 78 patient pairs) and malignant cells from the 7 supported patients, each with $\geq 20$ malignant cells (1,557 cells, 21 pairs). The split matters because P03's cells are 32.6\% malignant and P09's 2.2\%. On pooled cells, any factor that tracks malignancy would make P03 look different through prevalence alone. Within a stratum, for each factor:
\begin{enumerate}[leftmargin=1.6em]
\item bin all of the stratum's scores into 32 equal-width bins between the 1st and 99th percentile (the end bins take the tails), and turn each patient's counts into proportions (panels a and b);
\item for every pair of patients, take the area their two histograms share, the overlap coefficient OVL (panel c), and summarise the factor by the median OVL over all pairs (panel e);
\item for every patient, take the OVL between that patient's histogram and the pooled histogram of all the other patients' cells, the \emph{patient-vs-rest} overlap (panel d). A single patient that sits apart shows up here even when the median pair looks fine.
\end{enumerate}

\begin{figure}[H]
\centering
\includegraphics[width=\linewidth]{figures/fig03w1_one_factor.pdf}
\caption{\textbf{Worked example: how the test reads one factor} (IFN-$\alpha$ PCA $K = 40$, PC8, the axis the Stage~2 probe uses). \emph{(a, b)} Each row is one patient's PC8 histogram, as the share of that patient's cells in each of the 32 bins (cells in brackets; bold = highest-prevalence). Normal cells (a) and malignant cells (b) are binned separately, each on the 1st--99th percentile of its own stratum, so the two x-ranges differ. The number on the right is that patient's patient-vs-rest overlap on PC8 alone. Small patients give spiky histograms (one cell is 4.5\% of P07's 22). \emph{(c)} One pair drawn on the same bins: the shaded area is $\sum_b \min(p_b, q_b)$, their OVL. P02 and P06 are the median of the 21 malignant pairs, so 0.634 is PC8's median-pair overlap. \emph{(d)} P03 against the pooled histogram of the other six patients' 799 malignant cells: 0.857. \emph{(e)} All 21 malignant pairs; the boxed cell is the median. The same steps on normal cells give 78 pairs with median 0.909.}
\label{fig:03w1_one_factor}
\end{figure}

\textbf{From factors to one call} (Figure~\ref{fig:03w2_factors_to_call}). A representation has $K$ factors, so each stratum gives $K$ median-pair overlaps. They are averaged with weights equal to each factor's share of score variance among the stratum's cells: this is the \emph{weighted overlap}. Each patient's per-factor patient-vs-rest overlaps are averaged with the same weights, which gives one number per patient. On IFN-$\alpha$ PCA $K = 40$, PC8 carries 0.026 of the normal-cell weight and 0.022 of the malignant-cell weight, and no factor carries more than 0.100, so the summary describes the whole basis rather than the factor Stage~2 uses.

\begin{figure}[H]
\centering
\includegraphics[width=\linewidth]{figures/fig03w2_factors_to_call.pdf}
\caption{\textbf{From 40 factors to one call, for the worked example.} Top row normal cells, bottom row malignant cells; the x-axis is the same overlap scale everywhere. \emph{(a)} Each blue dot is one factor at its median-pair overlap (Figure~\ref{fig:03w1_one_factor}e computed for every factor), with area proportional to its share of score variance; the ringed dot is PC8. The solid black line is their variance-weighted mean, the \emph{weighted overlap} (normal 0.890, malignant 0.637). The dotted gray line is the same quantity after shuffling patient labels (mean of 25 shuffles); the \emph{gap} is the distance between the two. The dashed line is the aggregate gate. \emph{(b)} Each dot is one patient: its patient-vs-rest overlap per factor, averaged with the same weights. Only normal cells have a per-patient gate. The line under the panels shows the four checks this representation passes.}
\label{fig:03w2_factors_to_call}
\end{figure}

\begin{explainbox}{The quantities, exactly}
\emph{Overlap} of two histograms $p$ and $q$ on the same bins: $\mathrm{OVL}(p,q)=\sum_b \min(p_b,q_b)=1-\mathrm{TV}(p,q)$. It is 1 for identical histograms and 0 when no bin is shared. If A puts 50\% of its cells in each of two bins and B puts 20\% and 80\%, then $\mathrm{OVL}=\min(0.5,0.2)+\min(0.5,0.8)=0.7$. Two normal distributions one SD apart overlap by about 0.62.
\emph{Weighted overlap} of a representation in one stratum: $\sum_j w_j m_j$, where $m_j$ is the median OVL over all patient pairs on factor $j$, and $w_j$ is factor $j$'s score variance over the stratum's cells divided by the sum over all factors.
\emph{Patient-vs-rest overlap} of patient $i$: $\sum_j w_j\,\mathrm{OVL}(p_{ij}, r_{ij})$, where $r_{ij}$ is the histogram of all other patients' cells pooled.
\emph{Gap}: the mean weighted overlap over 25 label shuffles minus the observed weighted overlap.
\emph{$\eta^2$} of a factor: the share of its score variance explained by patient identity, i.e.\ the $R^2$ of a one-way ANOVA on patient.
\end{explainbox}

\textbf{Why not a KS test on cells.} A two-sample KS test asks whether two distributions are exactly equal. With 58,549 normal cells it rejects any shift, however small, and it treats the cells of one patient as independent. Aggregation does not need equality. It needs patients to occupy the same range, possibly shifted, rather than separate ranges. So we measure overlap and compare it with a null.

\textbf{The permutation null, and why its p-value cannot decide.} Within a stratum, we shuffle patient labels across cells, keeping each patient's cell count and the bins fixed, and recompute the weighted overlap (25 shuffles). This is the overlap that interchangeable patients of these sizes would show (the dotted line in Figure~\ref{fig:03w2_factors_to_call}). Its median over biological panels is 0.942 for normal cells and 0.659 for malignant cells; the malignant value is far below 1 because patients with 22--45 cells give noisy histograms. The p-value is $p=(1+\#\{\text{shuffles}\leq\text{observed}\})/26$. The observed overlap fell below all 25 shuffles in all 47 normal strata and in 43 of 47 malignant strata, giving $p=1/26=0.038$, the floor of this design. That only says patients are not interchangeable, which is expected given their different cell-type mix, genotype and batch. It cannot rank panels or decide aggregation.

\textbf{The three calls.} The call therefore rests on magnitude gates, set by judgment before the run and not yet calibrated (Figure~\ref{fig:03w3_call_types}). The subtype gates are checked first:
\begin{itemize}
\item \emph{Separate-subtype candidate} if either holds: some patient's normal patient-vs-rest $< 0.40$; or malignant weighted overlap $< 0.35$ together with a malignant gap $\geq 0.15$.
\item \emph{Aggregate} if all four hold: normal weighted overlap $\geq 0.75$; every patient's normal patient-vs-rest $\geq 0.60$; malignant weighted overlap $\geq 0.45$; malignant gap $< 0.15$.
\item \emph{Location shift} otherwise: the supports overlap, but one patient or the malignant cloud is displaced enough to miss an aggregate gate.
\end{itemize}
The malignant stratum has no per-patient gate; its patient-vs-rest values track cell count (Figure~\ref{fig:03b_patient_vs_rest}). A weighted average can hide one patient-private factor, so any factor with malignant $\eta^2 \geq 0.75$ and variance share $\geq 0.15$ is also flagged.

\begin{figure}[H]
\centering
\includegraphics[width=\linewidth]{figures/fig03w3_call_types.pdf}
\caption{\textbf{The three calls, drawn as schematic shapes (not data).} In each panel three blue curves stand for the other patients and the black curve for one patient. The overlap printed above each panel is computed on these synthetic curves, to show which side of the 0.60 and 0.40 per-patient gates each picture falls on. The rules under each panel are the exact gates from the decision file; the last line counts the 47 screened representations.}
\label{fig:03w3_call_types}
\end{figure}

\subsection*{What we found}

\begin{figure}[H]
\centering
\includegraphics[width=\linewidth]{figures/fig03a_decision_map.pdf}
\caption{\textbf{The aggregation call for every representation.} Each coloured point is one Stage~0 panel at its atlas-selected Stage~1 method and $K$ (colour = method; open squares = the six broad gene-bag controls), placed at its two weighted overlaps: normal cells on $x$, malignant cells on $y$ (the two black lines of Figure~\ref{fig:03w2_factors_to_call}a). Gray dots: the same representations' label-shuffle nulls, where interchangeable patients of these sizes would sit. Blue region: both weighted-overlap gates pass. Two gates are not drawn: the per-patient normal gate (Figure~\ref{fig:03b_patient_vs_rest}) and the malignant gap (how far a point's height sits below the height of its own gray dot; it must stay under 0.15). Read the distance to the red subtype band, which is large everywhere; the distance to the gray cluster only says that patients are not interchangeable.}
\label{fig:03a_decision_map}
\end{figure}

\textbf{Pooling passes almost everywhere} (Figure~\ref{fig:03a_decision_map}). The calls are 45 aggregate, 2 location shifts and 0 subtype candidates; 40 of the 41 biological panels are aggregate. Every point sits below its own null: on biological panels the gap is 0.053--0.163 for normal cells (median 0.098) and 0.008--0.111 for malignant cells (median 0.060). Patients are distinguishable but share most of their score support, and neither location shift is disjoint. Allograft rejection (FactoSig-promax $K = 40$) has high weighted overlaps (0.850 normal, 0.619 malignant), but P02's normal cells overlap the rest at 0.575, under the 0.60 gate. The HVG 10k bag has a malignant gap of 0.182.

\textbf{No patient sits apart, including the highest-prevalence ones} (Figure~\ref{fig:03b_patient_vs_rest}). In normal cells, patients' medians over the 41 biological panels run from 0.820 (P02) to 0.920 (P09); of 533 patient $\times$ panel values, one is below 0.60 (allograft rejection, P02, 0.575) and none below 0.40. In malignant cells, P02, P03 and P04 overlap the rest at medians of 0.788, 0.806 and 0.747. The seven malignant medians rank exactly in the order of malignant cell counts, from P07 (22 cells, 0.531; lowest in 38 of 41 panels) to P03 (758 cells, 0.806). Small-sample noise predicts this ordering, so the low overlaps of P07, P06 and P01 cannot be read as distinct biology, or as its absence, until a support-matched subsampling is done (D20).

\begin{figure}[H]
\centering
\includegraphics[width=\linewidth]{figures/fig03b_patient_vs_rest.pdf}
\caption{\textbf{Each patient against all the others, over the 41 biological panels.} Each gray dot is one patient's patient-vs-rest overlap on one panel's selected representation, averaged over that representation's factors with variance weights: the per-patient dots of Figure~\ref{fig:03w2_factors_to_call}b, now for all 41 panels (41 dots per patient). Black bar: the patient's median over the 41. Patients are ordered by their cell count in the stratum. Normal cells: aggregation needs every patient $\geq 0.60$ on a panel; one value of 533 falls under it and none under the 0.40 subtype gate. Malignant cells: no gate applies, and the medians rise with cell count.}
\label{fig:03b_patient_vs_rest}
\end{figure}

\begin{figure}[H]
\centering
\includegraphics[width=\linewidth]{figures/fig03c_factor_spectrum.pdf}
\caption{\textbf{Shared, shifted and private single factors (malignant cells).} The layout of Figure~\ref{fig:03w1_one_factor}b: each row is one patient's histogram on that factor's 32 bins (cells in brackets; bold = highest-prevalence), and the number on the right is that patient's overlap with all other patients pooled, on this factor alone. \emph{Left:} IFN-$\alpha$ PC8 is shared; P07's 0.483 rests on 22 cells. \emph{Middle:} the biological factor closest to the flag, on which P03 is shifted but still overlaps. \emph{Right:} a broad-bag factor on which P03 sits apart; the other patients' vs-rest overlaps are low too because P03's 758 cells, far to the right, make up 49--64\% of each one's rest pool, and only P03's own histogram is disjoint. Scores keep their raw, arbitrary sign.}
\label{fig:03c_factor_spectrum}
\end{figure}

\textbf{Only broad gene bags hide a patient-private factor} (Figure~\ref{fig:03c_factor_spectrum}). Of the 1,500 factors in the malignant stratum, 3 pass the single-factor flag. All three isolate P03 and come from broad gene bags; none of the 1,260 biological-panel factors meets the flag.

\begin{center}
\small
\setlength{\tabcolsep}{4pt}
\begin{tabular}{lrrr}
\toprule
Representation (malignant cells) & P03-vs-rest overlap & Variance share & $\eta^2$ \\
\midrule
HVG 10k bag, FactoSig $K = 40$, factor 21 & 0.007 & 0.373 & 0.913 \\
HVG 3k bag, FA $K = 40$, factor 27 & 0.013 & 0.196 & 0.907 \\
All gene sets (full\_34), FactoSig $K = 40$, factor 30 & 0.009 & 0.299 & 0.898 \\
\bottomrule
\end{tabular}
\end{center}

The closest biological factor is p53 pathway FactoSig $K = 40$ factor~9 ($\eta^2$ 0.569, variance share 0.104), on which P03 is shifted but still shares 0.263 of its histogram with the rest. The HVG 3k and full\_34 bags pass the aggregate gates despite their P03-private factors, which is why the per-factor flag must accompany the weighted average.

\textbf{What this means for downsampling healthy cells.} Pooling patients and thinning their healthy cells rest on the same property: each patient's normal cells occupy the score range that the other patients' normal cells occupy (532 of the 533 normal patient $\times$ panel values above are at least 0.60). Within single cell types the patients sit closer still: on IFN-$\alpha$ PCA $K = 40$, across eight predicted cell types (8--12 patients each), the gap to the shuffled null is 0.024--0.047, against 0.053 over all normal cells, and no patient's within-type patient-vs-rest overlap falls below 0.725. That is what makes downsampling healthy cells within cell types plausible: thinning P09's 16,751 normal cells (29\% of the normal stratum) type by type should not empty a region of healthy score space that only P09 covers. It does not show that downsampling helps, or why. The overlap was measured on bases fit with all cells, while downsampling changes the fit itself, the Stage~2 class ratio and the training $n$ (so $C$ must be retuned per condition). The earlier donor downsampling, stratified by patient and cell type (D10, D11), held $C$ fixed; its sparse probes mostly collapsed to all-zero, and under L2 held-out AP barely moved. The draft pre-registration (not yet locked), \nolinkurl{scripts/mrd_stage0_2/workstreams/W8_new_cohort/healthy_downsampling_plan.md}, uses this overlap both as the reason stratified thinning is plausible and as a check that the retained healthy cells still overlap the full healthy set; it runs when the team's new data arrive. \src{D14 focus\_celltype\_overlap.csv; 2026-08-21\_donor-sensitivity-d10-d11 (D10, D11)}

\begin{caveatbox}
\begin{itemize}
\item \textbf{Transductive basis.} Each patient helped fit the axes it is scored on, and an inductive basis could separate it more. The test must be rerun on inductive bases.
\item \textbf{Overlap is neither predictive sharing nor loading influence.} Whether the same factors drive lift is a separate, predictive question; whether P02/P03/P04 shape the axes needs refits without them. \src{D3; D17}
\item \textbf{The summary weights all factors by variance.} It is not about the factor Stage~2 uses: PC8 carries 0.022--0.026 of the IFN-$\alpha$ PCA $K = 40$ weight. A basis can pass while its one discriminative factor separates patients; the single-factor flag is the guard against that.
\item \textbf{Small malignant samples.} The overlaps for P07 (22 cells), P06 (44) and P01 (45) are dominated by noise. No support-matched subsampling has been done yet.
\item \textbf{Promax weighting.} Degenerate late factors dominate the variance weights; in allograft rejection $K = 40$, factor 40 carries 0.71 of the malignant weight. The 13 promax calls mostly reflect these factors.
\item \textbf{Not significance.} The gates are judgment calls, $p$ is at its floor, and each panel was tested at one lift-selected representation.
\item \textbf{The P03 axis is not a subtype.} Genotype, label provenance, batch and composition have not been excluded.
\end{itemize}
\end{caveatbox}

\begin{teambox}
``On the current bases, patients' factor-score distributions overlap enough to pool them. 45 of 47 representations pass the aggregate gates and none looks like a separate subtype. The one patient-private axis, P03's, appears only in broad gene bags, not in any biological program.''

\textbf{Likely question:} ``So every patient contributes to every factor? Doesn't that contradict the worry that P02/P03/P04 dominate the factors?''\\
\textbf{Answer:} No to both. This test asks whether their cells land in the same score range as everyone else's, and they do (malignant patient-vs-rest medians 0.788, 0.806, 0.747). Whether they steer the axes is answered by refitting without them. On IFN-$\alpha$, the axis that best matches PC8/FA7 after dropping all three still correlates with it at $|r|$ 0.9885--0.9969 (PCA and FA, $K$ = 10, 20, 40). \src{D17}

\textbf{Likely question:} ``Where do the gates come from?''\\
\textbf{Answer:} They were set before the run, by judgment, and are not calibrated. With $p$ at its floor the call rests on them; calibrating them with matched subsampling is the next step.
\end{teambox}

\subsection*{Next step}
\begin{itemize}
\item Harden the test (D20, W1.2). Down-weight scale-outlier promax factors, screen all PCA/FA representations, use $\geq 200$ shuffles, calibrate the gates and match cell counts by subsampling.
\item Rerun the test on the inductive leave-patient-out bases (W1.3, after W2.6).
\item Audit P03's broad-bag axis for genotype, label provenance, batch and composition (D25, W1.4).
\item Healthy-cell downsampling (W8, with the new cohort): reuse this overlap to check that the retained healthy cells still cover the full healthy range, per the draft pre-registration.
\end{itemize}

\subsection*{How this was computed}
{\small\raggedright
\emph{Data:} stored transductive scores \nolinkurl{stage1_dr/panel_id=*/method=*/k=*/seed=42/scores.npy} (no refit) for the 47 D1-selected representations; \texttt{CN.label} and patient from the cohort cache. \emph{Estimand:} within label, the variance-weighted mean over factors of the median pairwise 32-bin histogram overlap; patient-vs-rest overlap; per-factor $\eta^2$. \emph{Null:} 25 within-label shuffles of patient labels (counts and bins fixed). \emph{C policy:} not applicable (no Stage~2 fit). \emph{Code:} \nolinkurl{scripts/mrd_stage0_2/stage2_supervised/patient_factor_overlap.py}, via \nolinkurl{scripts/mrd_stage0_2/notebooks/stage0_2/stage2_old34_patient_factor_overlap.ipynb}. Figures: \nolinkurl{figs_03_aggregate.py} (this report's directory) reads the D14 tables and the gates in \nolinkurl{overlap_decision_rule.json}. For the worked example (Figures~\ref{fig:03w1_one_factor}--\ref{fig:03w2_factors_to_call}) and Figure~\ref{fig:03c_factor_spectrum} it recomputes single-factor histograms with the D14 functions (\texttt{evaluate\_stratum} and its helpers), on the same cells and bins, and checks every recomputed summary against the D14 tables. Figure~\ref{fig:03w3_call_types} is a schematic. The within-cell-type numbers are from \nolinkurl{focus_celltype_overlap.csv}. \emph{Artifacts} (relative to the experiment dir): \nolinkurl{analysis/scorecards/stage2_diagnostic/patient_factor_overlap_v1/}. \src{2026-10-02\_d14-patient-factor-overlap (D14); 2026-10-08\_d17-composition-matched-controls (D17); 2026-08-21\_donor-sensitivity-d10-d11 (D10, D11)}\par}
           % 3. Should we aggregate patients? (D14)
\labelnextsection{glance:04}\section{\texorpdfstring{IFN-$\alpha$}{IFN-alpha}: why does FA40 reach only P03 and P04?}
\label{sec:ifna-pca-vs-fa}

\begin{quote}
We looked at the odd observation wrt Interferon alpha response: PCA 40 factor achieve lift in four patients whereas FA(K=40) only capture P03 and P04.
(a) Data science angle: Hypothesizing PCA's orthogonality constraint forces later components to capture variance beyond P03 whereas FA keep allocating factors toward the highest-prevalence patient. Need to inspect the factor-by-patient matrix for FA to test this.
\end{quote}

\begin{answerbox}
The gap is a Stage~2 regularization threshold, not missing FA structure. At the atlas setting (L1, $C = 10^{-4}$) a factor can enter a fold's probe only if its class-mean gap beats $4/(C\,n_{\mathrm{train}})$; PCA40's best factor (PC8) clears that bar by 7.3--8.5\% in P02, P06 and P07, while FA40's best factor (FA7) misses it by 1.4--3.3\%, so those FA folds keep no factor and score lift 0 by construction. A 5\% higher $C$ gives FA40 18.06\pp{} in P02 and 6.97\pp{} in P07, and with $C$ tuned per fold FA40 has positive lift in all 7 supported patients under L1, L2 and elastic net, so hypothesis (a) is retired. Confidence: descriptive, one panel, 7 supported patients, transductive Stage~1 (the threshold account also holds on inductive bases); the rule reproduces 14/14 IFN-$\alpha$ folds and 2,772/2,772 audited atlas folds.
\end{answerbox}

\subsection*{What we did}

Everything below rests on one property of L1 probes.

\begin{explainbox}{The L1 entry rule, and why one $C$ is a different penalty in every fold}
\textbf{Intuition.} An L1 probe starts with every coefficient at zero. Moving one off zero costs a fixed price per unit; the initial benefit is how far apart the classes sit on that factor times how much loss the probe sums ($C$ times the training cells). A factor enters only if the benefit rate beats the price; otherwise its coefficient is exactly zero, not small.

\textbf{Exact.} Stage~2 minimizes $C\sum_{i} s_i\,\ell(y_i,\, b + w^{\top} z_i) + \lVert w\rVert_1$ with balanced weights $s_i = n/(2n_{c(i)})$, a free intercept and training-fold standardized scores. At $w = 0$ the optimal intercept gives every cell probability $1/2$, so the gradient for factor $j$ is $C\,(n/4)\,\mathrm{gap}_j$, with $\mathrm{gap}_j = \bar z_j^{\,\mathrm{normal}} - \bar z_j^{\,\mathrm{malignant}}$ in training-SD units. The subgradient (KKT) condition keeps $w_j = 0$ unless $|\mathrm{gap}_j| > \tau = 4/(C\,n_{\mathrm{train}})$. The \textbf{entry margin} is $m_j = |\mathrm{gap}_j|/\tau$; if every $m_j \leq 1$ the probe is all-zero, and a fold is \textbf{knife-edge} when its largest margin is within 10\% of 1. (Elastic net with L1 ratio $\alpha$: $\tau = 4\alpha/(C\,n_{\mathrm{train}})$.)

\textbf{Consequence.} Balanced weights remove the class sizes but not the total, so $\tau$ depends on $n_{\mathrm{train}}$, the cells outside the held-out patient. At $C = 10^{-4}$, $\tau$ is 0.678\,SD with P04 held out ($n_{\mathrm{train}} = 58{,}990$) and 0.930\,SD with P09 held out (42,992): the same $C$ demands a 37\% larger gap in the P09 fold. The atlas grid steps by $\times 4.64$ (13 values per penalty, $10^{-4}$ to $10^{4}$) and $10^{-4}$ is its floor, so the next $C$ the atlas could choose was $4.64\times10^{-4}$, far above the 1.5--3.5\% increase FA7 needs. The atlas still chose the floor because, of its 65 settings, that one had the highest pooled LOPO AP, which P03 dominates (see \emph{Reading the atlas heatmap} below).

\textbf{One fold in numbers.} With P02 held out, $n_{\mathrm{train}} = 58{,}656$, so $\tau = 4/(10^{-4} \times 58{,}656) = 0.682$\,SD. PC8's gap gives margin 1.085 and it enters (the PCA40 probe's held-out lift is 15.94\pp); FA7's gives 0.967, so the FA40 probe keeps no factor and scores lift 0 by construction. The section on learning curves (paragraph ``One fold worked through'') takes this fold step by step and shows what $C \times n$ would let FA7 in.
\end{explainbox}

The observation is one pair of atlas cells: the hallmark interferon-alpha response panel (92 genes) reduced by PCA or FA to $K = 40$ factors, both with the atlas-selected probe, L1 at $C = 10^{-4}$. \emph{Lift} is held-out AUPRC minus the held-out patient's malignant prevalence $\pi$ under leave-one-patient-out (LOPO). PCA40 has positive lift in P02 15.94, P03 13.40, P04 16.27, P07 11.08 and P06 0.06\pp{} (the team's ``four'' are those above 5\pp), FA40 in P03 14.37 and P04 19.87\pp. The other folds are \emph{all-zero}: the probe kept no factor, every held-out cell gets one score, AUPRC $= \pi$ (constant-score no-skill AP) and lift is 0 by construction. We tested the L1 optimality condition fold by fold, then moved $C$ or tuned it per outer fold by an inner LOPO over the six other supported patients, on the stored \emph{transductive} basis (fit without labels on all cells, the held-out patient's included) and on \emph{inductive} bases (fit without the held-out patient).


\subsection*{What we found}

\textbf{FA40 has the direction.} FA7 is FA40's largest-margin factor in every fold, as PC8 is PCA40's. On its own FA7 gives held-out lift 18.06\pp{} in P02 and 6.97\pp{} in P07 (PC8: 15.94 and 11.08\pp). They are related, not identical, directions: $r(\mathrm{PC8}, \mathrm{FA7}) = -0.685$ over all cells ($-0.575$ to $-0.722$ within patients), and PC8 has $R^2$ 0.920 on the 40 FA scores. Among their top-loading genes both put malignant cells toward IFITM2 and TENT5A and away from IFI27, PSME2 and PSMA3 (orientation from the L1 coefficient sign at $C = 10^{-4}$). The factor-by-patient matrix the team asked for (Figure~\ref{fig:04c}) shows no FA allocation toward the highest-prevalence patients. P02, P03 and P04 hold 8.4\% of FA40's score energy (summed squared deviations of factor scores from the cohort mean) against 8.2\% of cells (PCA40: 9.1\%), and P03 \emph{tops} (has the largest over-representation ratio of the 13 patients on) none of either basis's factors. P02 tops most factors in both bases (27 of 36 in FA40, 31 of 40 in PCA40), through excess spread on low-variance trailing factors, so it does not separate FA from PCA either.

\begin{figure}[H]
  \centering
  \includegraphics[width=\linewidth]{figures/fig04c_factor_patient_matrix.pdf}
  \caption{The factor-by-patient matrix for IFN-$\alpha$ PCA40 (left) and FA40 (right), transductive bases. Colour: the over-representation ratio, a patient's share of that factor's score energy divided by its share of cells (1$\times$ = in proportion; log scale, clipped at 0.35$\times$ and 2.8$\times$). Dot: the patient that tops the factor; numbers above the columns count the factors each patient tops. Look at the P02 column, red on the trailing factors in both bases, and at P03 and P04, near 1$\times$ throughout; the outlined rows are PC8 and FA7. Rows are factor indices, so row $j$ in one basis is not row $j$ in the other.}
  \label{fig:04c}
\end{figure}

\textbf{The rule explains every fold (Figure~\ref{fig:04a}).} Predicted and stored zero/nonzero outcomes agree in 14 of 14 folds. PC8's gap is 3--12\% larger than FA7's in every fold (0.704--0.770 vs 0.658--0.705\,SD), and in P02, P06 and P07 that small difference straddles the bar: PC8 margins 1.085, 1.074, 1.073; FA7 margins 0.967, 0.986, 0.981. P01 and P09 fall below the bar for both methods because they are the two largest patients (11,860 and 17,126 cells): holding either out shrinks $n_{\mathrm{train}}$ and raises $\tau$.

\begin{figure}[H]
  \centering
  \includegraphics[width=\linewidth]{figures/fig04a_entry_margin.pdf}
  \caption{Entry margins at the atlas setting, one pair of points per held-out supported patient. \textbf{(a)} Margin of PC8 (blue) and FA7 (orange); above 1 the factor enters, below it the probe is all-zero; gray band = knife-edge zone. Look at P02, P06, P07: PC8 just above, FA7 just below. \textbf{(b)} The same folds in training-SD units: gray bar $= 4/(C\,n_{\mathrm{train}})$, dots = class-mean gaps. The gaps are nearly constant; the bar is highest when P01 or P09 is held out.}
  \label{fig:04a}
\end{figure}

\textbf{Lift is a step in $C$, and tuning $C$ closes the gap (Figure~\ref{fig:04b}).} Once one factor enters, the held-out ranking is that factor's ranking, so lift jumps from 0 to the factor's full single-factor lift, whatever the coefficient's size. FA7 would enter P02 at $1.035\times$, P06 at $1.015\times$ and P07 at $1.020\times$ the atlas $C$; at $1.05\times$ FA40 covers the same five patients as PCA40, whose lifts do not change up to $1.25\times$ (except P01, 0.23\pp{} at $1.25\times$). Conversely, at $0.9\times$ PCA40 keeps only P03. With $C$ chosen by inner LOPO on the training patients only, both methods have positive lift in 7 of 7 supported patients under L1, L2 and elastic net (L1 ratio 0.5), and neither is ahead overall. PCA40 vs FA40, pooled AP / median lift / P20 lift: tuned L1 0.1001 / 8.71 / 1.37\pp{} vs 0.0894 / 7.03 / 1.25\pp; tuned L2 0.0871 / 5.60 / 1.17\pp{} vs 0.1325 / 6.38 / 1.85\pp; tuned elastic net 0.1209 / 10.16 / 1.30\pp{} vs 0.0891 / 7.40 / 1.21\pp. P09, all-zero for both at the atlas $C$, gets 5.63 (PCA40) and 7.03\pp{} (FA40) under tuned L1. One PCA40 advantage survives tuning, in P07 (22 malignant cells): 8.71 vs 4.97\pp{} (L1), 10.16 vs 7.40 (elastic net), 5.60 vs 5.07 (L2); it is one patient and its size depends on the penalty.

\begin{figure}[H]
  \centering
  \includegraphics[width=\linewidth]{figures/fig04b_policy_lift.pdf}
  \caption{Per-patient held-out lift of PCA40 (blue) and FA40 (orange) under four $C$ policies; filled = a factor entered, hollow = all-zero probe. Look at the hollow orange circles for P02, P06 and P07 in the first panel: they fill at $1.05\times$ the atlas $C$, and every patient is filled once $C$ is tuned per fold.}
  \label{fig:04b}
\end{figure}

\textbf{The Stage~1 basis does not change this.} On bases fit without the held-out patient, on normal cells only, with patient-balanced weights, or without P02, P03 and P04, tuned L1 and tuned L2 give positive lift in 7 of 7 patients for both methods (tuned L1 medians, PCA40 / FA40: inductive 10.55 / 6.69\pp, normal-only 6.98 / 7.57, balanced 9.25 / 6.85, without P02--P04 6.31 / 15.08). At the fixed atlas $C$ coverage moves with the basis in both directions: on the patient-balanced basis both cover 5 patients and FA40 has the higher median (14.28 vs 8.29\pp), while on the inductive basis PCA40 loses P07 because PC8's margin falls from 1.073 (stored basis) to 0.928.

\textbf{Why hypothesis (a) fails: FA is not ``PCA without orthogonality''.} Both methods give uncorrelated scores, the space Stage~2 sees: the largest off-diagonal score correlation is 0.000 for PCA40 and 0.003 for FA40 (FactoSig-promax, an oblique rotation, reaches 0.924), because unrotated maximum-likelihood FA makes $W^{\top}\Psi^{-1}W$ diagonal ($W$ loadings, $\Psi$ per-gene noise variances). What FA drops is orthogonality of the loadings in gene space (largest off-diagonal entry of $W^{\top}W$: 0.403, vs below $10^{-6}$ for PCA40); what it adds is a noise variance per gene, so it models shared rather than total variance. So L1 has no collinear FA columns to choose between. Stage~1 never sees labels, so it can favour a patient only through that patient's expression; it sometimes does (the P03-private axes in three broad gene bags, in the section on aggregating patients), but on IFN-$\alpha$ it does not for either method. The threshold accounts for the gap fold by fold.

\textbf{Reading the atlas heatmap.} A lift of exactly 0 there is this threshold, not an absence of signal; the cell view's per-fold strip flags all-zero folds (IFN-$\alpha$ FA40 is active in 3 of 13 folds, PCA40 in 11). The atlas picks $C$ by pooled LOPO AUPRC, which pools seven fold models and gives an all-zero fold one tied score, so it need not penalize all-zero folds: the fixed-$C$ L1 arms have the highest pooled AP of all IFN-$\alpha$ arms (PCA40 0.1630, FA40 0.1959) despite 2 and 5 all-zero folds. Across the atlas, 528 of 783 cells select the grid floor, and 136 L1 or elastic-net folds in 49 cells are knife-edge, where crossing the bar changes lift by $\geq 5$\pp{} in 41 of 131 evaluable folds (the primer section's note on reading a 0 has the details).

\begin{caveatbox}
\begin{itemize}
  \item Descriptive: one panel, $K = 40$, seven supported patients (P07 has 22 malignant cells, P06 44, P01 45); no test statistic on any PCA-versus-FA difference.
  \item Transductive Stage~1 for the main numbers (LOPO is label transfer, not inductive deployment); the inductive check covers this panel only.
  \item The atlas-$C$ arms (stored, $1.05\times$, fixed L2) are not nested: the atlas $C$ was chosen by pooled LOPO AUPRC including each held-out patient. Tuned arms are nested in $C$ only (method and $K$ were fixed after looking); tuned L2 lands on the grid floor in 6 of 7 FA40 and 5 of 7 PCA40 folds.
  \item PC8 and FA7 are not one axis, and neither is individually identifiable (PC8 relative eigengap 0.055; FA7 determinacy 0.649). Nothing here ranks PCA against FA.
\end{itemize}
\end{caveatbox}

\begin{teambox}
``On IFN-$\alpha$ the PCA-versus-FA difference was a regularization threshold, not a difference in what the factors capture: FA's best factor sat 1.4 to 3.3\% below the L1 entry bar in P02, P06 and P07; a 5\% higher $C$ brings those patients in, and with $C$ tuned per fold FA reaches all seven supported patients.''

\textbf{Likely question:} ``So which is better here, PCA or FA?''\\
\textbf{Answer:} Neither, on this evidence. With $C$ tuned per fold PCA40 is ahead under L1 (median 8.71 vs 7.03\pp) and elastic net (10.16 vs 7.40\pp), FA40 under L2 (6.38 vs 5.60\pp, also on P20 lift and pooled AP). The one per-patient gap that survives tuning (P07, 22 malignant cells) changes with the penalty and cannot be identified from one patient.

\textbf{Likely question:} ``Why did the atlas pick $C = 10^{-4}$ at all?''\\
\textbf{Answer:} It picks $C$ by pooled LOPO AP, which P03 dominates; the sparse FA40 probe has the highest pooled AP of any IFN-$\alpha$ arm (0.1959) despite five all-zero folds. 528 of 783 atlas cells sit at the grid floor, so the selection is censored there.
\end{teambox}

\subsection*{Next step}
\begin{itemize}
  \item Nested selection over method, $K$ and $C$ for the whole atlas, with pooled AP, median and P20 lift together and the 49 knife-edge cells flagged (W3.3, D23).
  \item Re-score the 49 knife-edge cells with per-fold tuned $C$ under L1, L2 and elastic net (proposed W3.8); add knife-edge badges to the per-fold strip (W7.1 follow-up).
  \item Name PC8 and FA7 from both poles under the tuned-$C$ probes (W5.1).
\end{itemize}

\subsection*{How this was computed}
{\small\raggedright
Data: stored transductive Stage~1 scores (seed 42) for the hallmark interferon-alpha response panel, PCA and FA at $K = 40$ (FA40 has 36 non-degenerate factors); 60,118 cells from 13 patients, LOPO over all 13, summaries over the 7 supported patients. Estimands: held-out lift per supported patient; entry margin $= |\mathrm{gap}|/(4/(C\,n_{\mathrm{train}}))$ on training-fold standardized scores. C policies: the stored atlas L1 probe at $C = 10^{-4}$ (cuML; the sklearn-objective refit reproduces its per-patient lift within 0.00014\pp); L1 refits at 0.90--1.25$\times$ that $C$; L2 at that $C$; L1, L2 and elastic net (L1 ratio 0.5) with $C$ tuned per outer fold by inner LOPO over the six supported training patients (median inner lift; ties to the sparser probe, then the smaller $C$; grid $1/\mathrm{logspace}(-4, 4, 17)$). Code: \texttt{scripts/\allowbreak mrd\_stage0\_2/\allowbreak workstreams/\allowbreak W3\_dr\_value\_regularization/\allowbreak run\_d19\_\allowbreak regularization\_\allowbreak contract.py} (module \texttt{regularization\_\allowbreak contract.py}) and \texttt{W2\_stage1\_\allowbreak basis\_variants/\allowbreak run\_d17\_\allowbreak basis\_variant\_pilot.py}; figures: \texttt{figs\_04\_threshold.py} in this report's directory. Artifacts (EXP-relative, under \texttt{analysis/\allowbreak scorecards/\allowbreak stage2\_diagnostic/}): \texttt{regularization\_\allowbreak contract\_v1/} (\texttt{d19\_anchor\_\allowbreak ifna\_l1\_entry\_check.csv}, \texttt{d19\_ifna\_\allowbreak pca40\_vs\_fa40\_by\_patient.csv}, \texttt{d19\_ifna\_\allowbreak arm\_summary.csv}); \texttt{basis\_variant\_\allowbreak pilot\_ifna\_v1/}; \texttt{feedback\_deep\_\allowbreak dive\_20260917/}; \texttt{factor\_patient\_\allowbreak matrix\_v1/} (Figure~\ref{fig:04c} from \texttt{factor\_patient\_\allowbreak matrix\_long.csv.gz}); \texttt{factor\_cards\_v1/}.
\src{2026-10-08\_d19-regularization-contract (D19); 2026-10-08\_d17-basis-variant-pilot (D17); 2026-09-17\_team-feedback-f1-f7 (F3, F6); 2026-10-02\_f9-ifna-pc8-loadings (F9); 2026-10-08\_d18-factor-cards (D18); 2026-10-08\_d16-factor-patient-matrix (D16); 2026-10-08\_w7-atlas-strip-factor-explorer; H-S2-08 supported (descriptive), H-S2-07 retired}
\par}
                         % 4. IFN-alpha PCA40 vs FA40 (D19, D17)
\labelnextsection{glance:05}\section{What is PC8, and how should we read a factor?}
\label{sec:pc8}

\begin{quote}
``Two groups of genes in the PC8's loading with opposite signs: Positive: IFI27, PSME2, PSMA3, CNP, LAP3, USP18. Negative: TXNIP, IFITM2, TENT5A, NCOA7, TRIM5, ELF1.'' \ldots\ ``(expand beyond IFN-alpha; can we view these in the interactive app?): correlate factors with other modalities (cell types, surface protein); review interpretability frameworks from other fields (e.g.\ psychology's factor rotation); no causal claim without intervention.''
\end{quote}

{\small\color{inkmuted}\emph{Reading guide (v2).} What ``not identifiable'' means: the box ``What \emph{not identifiable} means'' and Figure~\ref{fig:05d}. What the 1,810 factor cards are: Table~\ref{tab:05cards}. How to read $\rho = +0.251$ and $+0.222$: the box ``How to read $\rho$'' and Figure~\ref{fig:05e}. Factor $\times$ cell type and surface protein for every card, not only PC8: ``Every card'' and Figure~\ref{fig:05f}.\par}

\begin{answerbox}
PC8 is a single axis that contrasts two groups of IFN-$\alpha$ genes. Once PC8 is oriented by its L1 coefficient, the team's ``negative'' list (TXNIP, IFITM2, TENT5A, NCOA7, TRIM5, ELF1) is the pole that malignant cells score toward. FA7 is a related axis (Tucker $\phi$ 0.691, below the 0.85 ``fair'' similarity band). Either factor alone gives 4.81--19.87\pp{} of held-out lift in the same five supported patients (P02, P03, P04, P07, P09) and under 1\pp{} in P01 and P06.

\emph{Not identifiable} means the data pin down a small subspace (the span of PC8 and its near-tied neighbours, which every refit reproduces) but not which single direction inside it is ``factor 8'': PC8's variance is within 5.5\% of PC7's, and leaving P01 out of Stage~1 turns PC8 so far that its best single match falls to $|r|$ 0.58 while 98.8\% of it stays in the refit basis; FA7's scores are estimates (determinacy 0.649). So name the axis by its two poles, not ``factor 8''.

\emph{Surface protein} (now for all 1,801 non-degenerate cards): PC8 tracks CD34 in normal cells inside every supported patient (median $\rho$ $+0.245$, 7/7) but not inside predicted cell types ($+0.019$), so it is a cell-identity association; its malignant-cell CD99 correlate (pooled $+0.222$) is $+0.159$ within patients (6/7) and fails the rule fixed in advance.

Descriptive: one panel for the gene poles; transductive Stage~1; seven supported patients; ADT without dsb.
\end{answerbox}

\subsection*{What we did}
We built a \emph{factor card} for every factor of 56 stored Stage 1 representations (1,810 cards, nothing refit): orientation, gene poles, per-patient held-out lift, matches across methods and K, identifiability, and cell-type and ADT associations. Table~\ref{tab:05cards} shows where the 1,810 come from. They are every factor of 56 of the atlas's 736 stored DR bases (13,540 factors in all), not every Stage~0 $\times$ Stage~1 combination: one basis per panel, plus the IFN-$\alpha$ and focus-panel bases the 2026-10-08 questions needed. How a card can be read depends on its basis's stored probe, which splits the 1,810 into two groups that do not overlap. In 10 of the 56 representations the probe is sparse (in the median supported fold at most half the factors are non-zero): their 330 cards are 119 \emph{selected drivers} (non-zero in at least one supported fold) and 211 never selected. In the other 46 the probe is dense (31 L2, 14 elastic net, 1 L1): their 1,480 cards are partial effects among many correlated factors and cannot be read factor by factor; 46 of them are never selected either. So $119 + 211 + 1{,}480 = 1{,}810$, and the never-selected cards, $211 + 46 = 257$, include the 9 degenerate ones. For IFN-$\alpha$ PCA40 and FA40 the stored probe is L1 at $C=10^{-4}$.

\begin{table}[H]
\centering
\small
\caption{\textbf{Where the 1,810 factor cards come from.} One card per factor of one stored Stage~1 basis (panel $\times$ method $\times$ $K$, seed 42, transductive). A basis with $K$ factors gives $K$ cards; \emph{non-degenerate} excludes factors with zero score variance (IFN-$\alpha$ FA40 factors 37--40 and MYC targets v2 FA40 factors 36--40), so it is the sum of the effective $K$.}
\label{tab:05cards}
\setlength{\tabcolsep}{4pt}
\begin{tabularx}{\linewidth}{@{}Xrrr@{}}
\toprule
Source of the basis & Bases & Cards & Non-degenerate \\
\midrule
Atlas screen: for each of the 41 biological panels, the DR cell (method, $K$) with the highest patient-median lift (the bases of section \emph{Should we aggregate patients}) & 41 & 1,260 & 1,255 \\
Atlas screen: the same choice for the 6 broad gene-bag controls (HVGs 500, 1,000, 3,000, 10,000; all Stage~0 gene sets; core-only programs) & 6 & 240 & 240 \\
Focus panels: PCA40 and FA40 for allograft rejection, apoptosis and TNF-$\alpha$ via NF-kB & 6 & 240 & 240 \\
IFN-$\alpha$ FA40, the pair of the screen's IFN-$\alpha$ PCA40 & 1 & 40 & 36 \\
IFN-$\alpha$ FA10 and FA20, to follow FA7 across $K$ & 2 & 30 & 30 \\
\midrule
Total & 56 & 1,810 & 1,801 \\
\bottomrule
\multicolumn{4}{@{}p{\linewidth}@{}}{\footnotesize By method: FA 590 cards (19 bases), FactoSig 380 (11), FactoSig-promax 410 (13), PCA 430 (13). By $K$: 37 bases at $K = 40$, 14 at 20, 5 at 10. Orientation: 1,217 cards oriented by a consistent probe sign, 336 whose sign flips across folds, 257 never selected.} \\
\end{tabularx}
\end{table}

\textbf{Lift} is held-out AUPRC minus the patient's malignant prevalence $\pi$ (constant-score no-skill AP $=\pi$). \textbf{Univariate lift} scores cells with one factor alone, its sign fixed on the training folds. Stage 1 is \textbf{transductive} (fit on all cells, held-out patient included, without labels), so LOPO lift is label transfer, not inductive deployment.

\textbf{New in this version: factor $\times$ cell type and surface protein for every card (D27).} The cards held, for each factor, the pooled correlation with each surface protein over all cells of one label. D27 recomputes this for all 1,801 non-degenerate cards and adds the two views a pooled number cannot give: Spearman $\rho$ inside each supported patient (summarised as the median over the seven and the number that agree in sign), and inside each supported patient $\times$ predicted cell type ($\geq 30$ cells). The rule was fixed before computing: an association is \emph{robust} if the median within-patient $|\rho|$ is at least 0.20 and at least 5 of the 7 patients agree in sign. The 9 isotype antibodies (no specific target) are the noise reference. Nothing is refit, so there is no $C$ policy.

\begin{explainbox}{How to read a factor}
\small
\begin{enumerate}[leftmargin=*,itemsep=0pt,topsep=0pt]
\item \textbf{The sign is arbitrary, so orient by $\beta$.} The oriented loading is $\mathrm{sign}(\beta)\times$loading, so positive means toward malignant. PC8's $\beta$ is $-0.040$ to $-0.064$ in all five active supported folds.
\item \textbf{Numbers are labels.} Match by score correlation (Hungarian) and many-to-one $R^2$; one factor can split over two.
\item \textbf{L1 keeps one of several correlated directions.} A factor enters only if its standardized class-mean gap exceeds about $4/(C\cdot n_{\text{train}})$ (0.68--0.93\,SD across folds here); gap over threshold is the \emph{entry margin}. $\beta=0$ means below threshold, not no signal.
\item \textbf{FA loadings are generative} (how expression moves with the factor), not score weights.
\item \textbf{Check identifiability, and read ablation as attribution.}
  \begin{itemize}[leftmargin=*,itemsep=0pt,topsep=0pt]
  \item A relative eigengap of 10\% or less means the component can rotate.
  \item FA determinacy $\rho<0.707$ means equally valid scores may correlate only $2\rho^2-1$.
  \item Tucker $\phi$ (loading cosine) $\geq 0.95$ means the same factor; 0.85--0.94 is fair.
  \item Ablation is attribution inside the model, not causation.
  \end{itemize}
  The box ``What \emph{not identifiable} means'' below works these through for PC8 and FA7.
\item \textbf{Read a protein correlate inside patients and inside cell types} before naming a factor by it (box ``How to read $\rho$'').
\end{enumerate}
\end{explainbox}

\subsection*{What we found}
\textbf{The two lists are the two poles of one axis} (Figure~\ref{fig:05a}). The malignant pole runs from TXNIP ($+0.269$) to ELF1 ($+0.176$). The normal pole runs from IFI27 ($-0.315$) to USP18 ($-0.161$). Interferon-stimulated genes sit on both poles, so the axis contrasts two parts of the response rather than measuring ``more interferon''. FA7 agrees in sign on 19 of the 20 genes shown; IL7 is the exception. The two factors are related, not identical: PC8 has $R^2$ 0.920 on the FA40 scores, yet its score $r$ with FA7 is only $-0.685$, and their loading cosine (Tucker $\phi$) is 0.691, below the 0.85 ``fair'' band. By psychometric convention these are different factors that share a subspace: one subspace, cut two ways. PC8 recurs as factor 8 of PCA $K = 20$ ($r$ $+0.942$) and FA7 as factor 7 of FA $K = 10$ ($r$ $+0.932$); no $K = 5$ basis spans either.

\begin{figure}[H]
\centering
\includegraphics[width=\linewidth]{figures/fig05a_pc8_poles.pdf}
\caption{\textbf{PC8's poles, oriented so positive points toward malignant} (top 10 per pole; bold = genes from the team report). Bars: PC8 weights (orientation $-1$, so the malignant-leaning pole has negative raw PC8 loadings). Dots: FA7 loadings for the same genes (orientation $+1$). The scales differ, so compare signs and order. Orientation is from L1 at $C=10^{-4}$.}
\label{fig:05a}
\end{figure}

\begin{explainbox}{What \emph{not identifiable} means, and why we name the axis by its poles}
\emph{In plain words.} A factor is \emph{identifiable} when the data pin its direction down uniquely: refit the same method on slightly different cells and the same factor comes back. When they do not, the data still fix a \emph{subspace}, the span of a few neighbouring factors (for PC8, itself and the components whose variances nearly tie with it), but not where one factor points inside it. Every rotation inside that subspace fits about equally well, and which one a fit returns depends on noise, the seed or a few cells. ``Factor 8'' is then a name for one arbitrary cut through the subspace. ``Neither is identifiable alone'' means this holds for PC8 and for FA7, each for its own reason; the two are related cuts through overlapping subspaces (PC8 has $R^2$ 0.920 on the FA40 scores), not the same factor.

\emph{PCA: a small eigengap.} PCA orders components by the variance they explain. PC8 explains 1.40\% of the panel's variance, PC7 1.48\% and PC9 1.26\%: PC7 is 5.5\% above it and PC9 10.0\% below (Figure~\ref{fig:05d}a; relative eigengap 0.055). When two variances are nearly tied, a small change in the data can mix the two components; the Davis--Kahan theorem bounds how far a component turns by the size of the change divided by the gap, so a small gap allows a large turn. Past PC3 this is the rule in this basis: 37 of the 40 IFN-$\alpha$ components have a neighbour within 10\%.

\emph{FA: low determinacy.} FA scores are not observed; they are estimated from the genes. Determinacy $\rho$ is the correlation between the factor and its best estimate. FA7 has 0.649, and 30 of the 36 non-degenerate FA40 factors are below 0.707 (Figure~\ref{fig:05d}b). Below 0.707, two equally valid sets of scores can correlate as little as $2\rho^2-1$; for FA7 that is $-0.16$.

\emph{What it looks like when Stage~1 is refitted} (Figure~\ref{fig:05d}c; the refits of sections \emph{Do the highest-prevalence patients shape the factors?} and \emph{Option 2}). Two numbers per refit: the best single matching factor ($|r|$) and the share of the stored axis that the whole refit basis reproduces ($R^2$). Leaving out P02, P03, P04 or all three barely moves PC8 ($|r| \geq 0.995$). Leaving out P01 drops PC8's best single match to 0.58, although $R^2$ is 0.988; patient-equal weights give 0.90 ($R^2$ 0.990). FA7 turns when P09 is left out (0.76; $R^2$ 0.981). Across all these refits $R^2 \geq 0.977$: the subspace survives, the single factor need not.

\emph{What to do.} Name the axis by both poles (the genes on each end), with its orientation and the Stage~2 penalty; compare bases by subspace ($R^2$, principal angles), not by factor number; and read ``PC8'' as shorthand for the direction, not as a unique program.
\end{explainbox}

\begin{figure}[H]
\centering
\includegraphics[width=\linewidth]{figures/fig05d_identifiability.pdf}
\caption{\textbf{Why PC8 and FA7 are not identifiable on their own.} (a) Share of the 92-gene panel's variance explained by each of the first 16 IFN-$\alpha$ PCA40 components; PC8 filled, its neighbours PC7 and PC9 ringed; gray band: $\pm 10\%$ of PC8's value. (b) Factor-score determinacy of the 36 non-degenerate IFN-$\alpha$ FA40 factors; dashed line at 0.707 ($\rho^2 = 0.5$). (c) Stored transductive axis against Stage~1 refits (D17; $K = 40$; every refit uses fewer cells or different cell weights). Filled dot: $|r|$ with the single best-matching refit factor (one-to-one matching). Open dot: many-to-one $R^2$, the share of the stored axis that all refit factors together reproduce. Blue PC8, orange FA7. Look at the P01 row: the open blue dot stays near 1 while the filled one falls to 0.58.}
\label{fig:05d}
\end{figure}

\textbf{The axis carries signal beyond P03 and P04} (Figure~\ref{fig:05b}). Alone, PC8 gives 15.94, 13.40 and 16.27\pp{} in the highest-prevalence patients P02, P03 and P04. It gives 11.08\pp{} in P07 and 5.47\pp{} in P09. FA7 gives 18.06, 14.37, 19.87, 6.97 and 4.81\pp{} in the same five patients. The FA40 probe uses FA7 only in P03 and P04. In P02 and P07, FA7's entry margins are 0.967 and 0.981. With a 5\% larger $C$ it enters and gives those same lifts. $\eta^2$ is the share of score variance explained by patient within a label. For PC8 it is 0.010 in normal cells and 0.031 in malignant cells. A Stage 1 fit on normal cells alone recovers the axis (matched $|r|$ 0.9936). So normal cells already vary along PC8, and malignant cells sit toward one end.

\textbf{The axis is a combination, not its best gene.} A gene's $|$loading$|$ tracks its own held-out lift only weakly (Spearman $\rho$ 0.30, 92 genes). The best single gene, BST2, reaches only 4.06\pp{}, even with its sign picked on the held-out patient (an optimistic choice); PC8's median is 11.08\pp.

\begin{figure}[H]
\centering
\includegraphics[width=\linewidth]{figures/fig05b_axis_lift.pdf}
\caption{\textbf{Alone, each factor gives 4.81--19.87\pp{} of held-out lift in five supported patients (P02, P03, P04, P07, P09) and under 1\pp{} in P01 and P06.} Univariate lift uses the training-fold sign. The strip shows whether the stored L1 probe ($C=10^{-4}$) kept the factor in that patient's fold. Open squares mark all-zero probes, whose fitted lift is 0 by construction.}
\label{fig:05b}
\end{figure}

\textbf{Cell type and surface protein for PC8 and FA7} (Figure~\ref{fig:05c}, Table~\ref{tab:05adt}). In normal cells, cell type explains part of the score: $\eta^2$ 0.081 pooled and 0.105 within patients (median), against 0.0047 expected with no association and 0.010 for patient. The score is highest in EBM, MEP-MkP and HSC/MPP1. Within cell types, though, malignant cells score higher in 8 of 10. Held out, PC8 keeps its lift inside HSC/MPP2 alone: $+21.8$\pp{} in P02, $+12.7$ in P03, $+10.3$ in P07 and $+27.1$ in P09. An oracle that scores each held-out cell by its own patient's malignant fraction within the cell's predicted type (it uses held-out labels, so it is a ceiling, not a competitor) beats PC8 in P02 (57.33 vs 15.94\pp) and P03 (38.04 vs 13.40\pp) but not in P07 (8.17 vs 11.08\pp): in the highest-prevalence patients, which cell types are malignant carries more ranking signal than the axis. Table~\ref{tab:05adt} gives the surface-protein correlates pooled, within patient and within patient $\times$ cell type; the next box explains how to read them.

\begin{figure}[H]
\centering
\begin{minipage}[t]{0.47\linewidth}\vspace{0pt}
\includegraphics[width=\linewidth]{figures/fig05c_celltype.pdf}
\captionof{figure}{\textbf{Within most cell types, malignant cells sit further along the axis.} Mean oriented PC8 score (cohort z) per predicted type, pooled over patients. (n) is the number of malignant cells.}
\label{fig:05c}
\end{minipage}\hfill
\begin{minipage}[t]{0.505\linewidth}\vspace{0pt}
\captionof{table}{\textbf{Strongest within-patient ADT correlates of PC8 and FA7.} Spearman $\rho$ of the oriented score with CLR ADT, within label: pooled over all cells; median over the 7 supported patients (number with the same sign); median inside patient $\times$ cell type.}
\label{tab:05adt}
\footnotesize
\setlength{\tabcolsep}{2.6pt}
\begin{tabular}{@{}llrrr@{}}
\toprule
 & & pooled & within (k/7) & in type \\
\midrule
\multicolumn{5}{@{}l}{\emph{PC8, normal cells (58,549)}} \\
 & CD34 & $+0.251$ & $+0.245$ (7) & $+0.019$ \\
 & CD22 & $-0.221$ & $-0.246$ (6) & $-0.046$ \\
\multicolumn{5}{@{}l}{\emph{PC8, malignant cells (1,569)}} \\
 & CD99 & $+0.222$ & $+0.159$ (6) & $+0.150$ \\
 & CD48 & $-0.192$ & $-0.308$ (6) & $-0.011$ \\
\multicolumn{5}{@{}l}{\emph{FA7, normal cells}} \\
 & CD22 & $-0.354$ & $-0.312$ (7) & $-0.085$ \\
 & CD34 & $+0.260$ & $+0.298$ (7) & $-0.004$ \\
\multicolumn{5}{@{}l}{\emph{FA7, malignant cells}} \\
 & CD116 & $-0.108$ & $-0.316$ (6) & $-0.035$ \\
 & CLEC12A & $-0.158$ & $-0.291$ (7) & $-0.041$ \\
\midrule
\multicolumn{2}{@{}l}{Isotype ref., PC8} & \multicolumn{3}{r@{}}{0.109 normal / 0.227 malignant} \\
\multicolumn{2}{@{}l}{Isotype ref., FA7} & \multicolumn{3}{r@{}}{0.071 normal / 0.257 malignant} \\
\bottomrule
\end{tabular}

{\footnotesize Isotype ref.: the card's largest $|$median within-patient $\rho|$ over the 9 isotype antibodies.\par}
\end{minipage}
\end{figure}

\begin{explainbox}{How to read $\rho$: the PC8 $\times$ CD34 $+0.251$ and PC8 $\times$ CD99 $+0.222$ of report v1}
\emph{Which cells.} One label at a time, every patient together. $+0.251$ ranks the 58,549 normal cells of all 13 patients by PC8 score and by CD34 level and correlates the two rankings (Spearman). $+0.222$ does the same over the 1,569 malignant cells. Neither is a malignant-vs-normal contrast.

\emph{Which score and which protein.} PC8 is oriented (multiplied by $-1$), so $\rho>0$ means cells with more of the protein sit toward the malignant-leaning pole. The protein value is the CLR of raw ADT counts: $\log(1+\text{count})$ minus the cell's mean of that quantity over all 170 antibodies, with no dsb background correction. Spearman uses ranks, so only the per-cell centring matters.

\emph{How big is $+0.25$?} Pick two normal cells at random: the one with more CD34 has the higher PC8 score about 58\% of the time instead of 50\%, and about 6\% of the rank variance is shared ($\rho^2$). That is a population tendency, not a marker that identifies the axis in a single cell. With 58,549 cells the sampling SE is about 0.004, so a p-value carries no information; inside P07's 22 malignant cells it is about 0.23.

\emph{Noise reference.} Isotype antibodies bind nothing specific, yet the same calculation gives them up to 0.109 (normal) and 0.227 (malignant) on this card, from non-specific binding and the CLR centring. A value near 0.2 must be read against that.

\emph{Pooled vs within patient.} A pooled $\rho$ mixes how cells vary inside each patient with how the patients' averages differ, in score and in staining (patients are nested in processing batches, so the two cannot be told apart). Ranking inside one patient removes the between-patient part; we report the median over the seven supported patients and how many agree in sign. Ranking inside patient $\times$ predicted cell type also removes cell identity.

\emph{Why a pooled value holds or shrinks within patients.} In normal cells every patient contributes thousands of cells over the same cell types, and the patients' score distributions overlap (section \emph{Should we aggregate patients}). Pooled and within-patient $\rho$ then agree: median ratio 0.89 over the 23,133 card $\times$ marker pairs with $|$pooled$| \geq 0.20$, and only 2 reverse sign (with a within-patient $|\rho| \geq 0.10$). In malignant cells P03 holds 758 of the 1,569 cells and patients differ in both score and staining, so a pooled $\rho$ can come from patient averages alone. There the median ratio is 0.24 over 16,190 pairs, and 707 reverse sign. D27's pooled values reproduce the cards' (D18) values to $5\times10^{-11}$ (all 10 regression checks pass), so what changes between the two numbers is the stratification, not the computation.

\emph{Confounds that remain.} Patient and batch; recipient origin (every malignant cell is recipient-derived, so recipient genotype, e.g.\ HLA-binding antibodies, is not separated from malignancy); predicted cell type (a label transferred from a healthy reference); and correlated factors across bases, so counts of robust pairs overstate independent findings.
\end{explainbox}

\textbf{The two numbers taken apart} (Figure~\ref{fig:05e}). \emph{PC8 $\times$ CD34, normal cells}: every supported patient agrees ($+0.14$ to $+0.38$; median $+0.245$, 7/7), so this is a within-patient association, well above the card's isotype reference (0.109). Inside patient $\times$ predicted cell type it is $+0.019$ (71 strata): it is carried by which cell types sit at the malignant-leaning end (EBM highest, $+0.79$ in cohort z; ncMono lowest, $-0.73$), not by variation inside a cell type. In normal cells PC8's CD34 correlate is a cell-identity axis. \emph{PC8 $\times$ CD99, malignant cells}: $+0.159$ within patients, 6 of 7 positive (P07, 22 cells, $-0.12$), so it fails the 0.20 rule and sits inside this card's malignant isotype reference (0.227). What there is persists inside cell types ($+0.150$; 11 strata from P02, P03, P04 and P09). \emph{PC8 $\times$ CD48, malignant cells} shows the opposite case: $-0.308$ within patients (6/7) but only $-0.192$ pooled, because P03, with half the malignant cells, sits at $+0.03$. PC8's other strong malignant correlates (CD116 $-0.287$, CD33 $-0.272$, CD146 $+0.270$) share this pattern: P03 sits near 0 on all of them.

\begin{figure}[H]
\centering
\includegraphics[width=\linewidth]{figures/fig05e_rho_anchors.pdf}
\caption{\textbf{One pooled $\rho$, taken apart.} Spearman $\rho$ of the oriented PC8 score with CLR surface protein. Each panel title gives the verdict of the robust rule and, below it, both of its conditions with this pair's values: (b) has 6 of 7 patients of one sign, which passes, but fails on size. Top rows: inside each supported patient (number of cells of that label in brackets; large dark dot: P03, which holds 758 of the 1,569 malignant cells). Bottom rows: the median of the seven (filled circle; ``robust'' needs $|$median$| \geq 0.20$, dotted lines, and $\geq 5$ of 7 with the same sign), the pooled value over all cells of the label (diamond; the number in report v1) and the median inside supported patient $\times$ predicted cell type strata with $\geq 30$ cells (square). Gray band: the card's isotype reference for that label. Look at (a): every patient sits near $+0.25$, but the square is near 0. Then (c): the patients sit near $-0.3$, but P03 sits at 0 and drags the diamond toward it.}
\label{fig:05e}
\end{figure}

\textbf{Every card: factor $\times$ surface protein and cell type} (Figure~\ref{fig:05f}). \emph{Normal cells.} 20,274 of 289,961 card $\times$ marker pairs (7.0\%, on 843 of the 1,801 cards) meet the robust rule, against 2.3\% of isotype rows under the same rule. Pooled and within-patient $\rho$ agree (Figure~\ref{fig:05f}c). The strongest are myeloid, erythroid and progenitor markers (CD11a, CD18, CD33, CD200, CLEC12A, CD123, CD151, CD71; Figure~\ref{fig:05f}a), up to $+0.809$ within patients (HVGs (500) FA40 factor 3 $\times$ CD11a, 7/7). Inside patient $\times$ cell type most of them shrink: only 729 of the 20,274 keep at least half their size (median retention 0.15). In 1,665 of 1,801 cards, predicted cell type explains more of the score within patients than patient identity explains across patients (Figure~\ref{fig:05f}d). So in normal cells a factor's surface-protein correlates mostly say which cell types the factor separates. \emph{Malignant cells.} 8,474 pairs (2.9\%, on 954 cards) meet the rule, against 1.7\% of isotype rows. Pooled values are a poor guide here: some of the strongest within-patient associations are weak when pooled (CD133 on HVGs (3000) FA40 factor 16: $+0.598$ within, 7/7, $+0.181$ pooled; Figure~\ref{fig:05f}b). Of the 8,474, 1,797 keep at least half their size inside cell types (median retention 0.26), but the malignant cell-type strata come from four patients (P02, P03, P04, P09) only. Of the 119 selected drivers, 61 have at least one robust marker in normal cells and 75 in malignant cells. Every number here is a screen over correlated factors and 161 markers, not a test.

\begin{figure}[H]
\centering
\includegraphics[width=\linewidth]{figures/fig05f_modality_all.pdf}
\caption{\textbf{Factor $\times$ surface protein and cell type over all 1,801 non-degenerate cards (D27).} (a, b) For each of the ten markers with the strongest robust association on an oriented card, that card (panel $\cdot$ method $K$ factor; promax = FactoSig-promax) and its $|\rho|$: median within patient (dot, coloured by Stage~1 method: orange FA, aqua FactoSig or FactoSig-promax), pooled (diamond) and inside patient $\times$ cell type (square); the sign of the within-patient median is in brackets, the number of patients sharing it on the right. Gray band: the 99th percentile of the isotype rows ($|$median within$|$ 0.231 normal, 0.225 malignant). Broad gene-bag controls supply 3 of the 10 cards in (a) and 7 of the 10 in (b). (c) For every card $\times$ marker pair with $|$pooled $\rho| \geq 0.20$ (23,133 normal, 16,190 malignant), the within-patient median divided by the pooled value: 1 means it holds within patients, 0 that it is all between patients, below 0 that it reverses. (d, e) Per card: $\eta^2$ of predicted cell type within patients (median over supported patients) against $\eta^2$ of patient (from the cards); above the dashed diagonal, cell type explains more. Malignant within-patient $\eta^2$ uses P02, P03, P04 and P09 only.}
\label{fig:05f}
\end{figure}

\textbf{Beyond IFN-$\alpha$: which factors a probe uses.} 46 of the 56 representations have a dense selected probe that cannot be read factor by factor. Only 119 cards in 10 representations qualify as selected drivers; besides IFN-$\alpha$ they include TNF-$\alpha$ via NF-kB PCA40, inflammatory interferon FA40 and TGF-$\beta$ FA40. An L1 zero is not absence of signal. 101 of the 248 never-selected non-degenerate cards still reach $\geq 5$\pp{} in some patient; IFN-$\alpha$ PCA40 factor 4, for example, reaches 25.40\pp{} in P02.

% Keep the heading with the start of its list: a new page only when fewer than six lines are left. Unlike
% \needspace or a negative \penalty, this adds no break point, so it cannot end a page early
% (report.tex defines the same macro; \providecommand keeps the standalone compile working).
\makeatletter
\providecommand{\softneedspace}[1]{\par\begingroup\dimen@ii\pagegoal\advance\dimen@ii-\pagetotal
  \ifdim #1>\dimen@ii\ifdim\dimen@ii>\z@\newpage\fi\fi\endgroup}
\makeatother
\softneedspace{6\baselineskip}
\subsection*{Frameworks from other fields}
The factor-interpretability literature review suggests five practices worth adopting:
\begin{enumerate}[leftmargin=*,itemsep=0pt,topsep=0pt]
\item From psychometrics: name a factor only after reporting its eigengap (Davis \& Kahan 1970), determinacy (Guttman 1955; Grice 2001) and Tucker $\phi$ against the 0.85/0.95 bands (Lorenzo-Seva \& ten Berge 2006). The cards already report all three.
\item From neuroimaging decoders: read each probe's forward pattern $\mathrm{cov}(X,s)/\mathrm{var}(s)$, which is rotation-invariant, instead of its weights (Haufe et~al.\ 2014).
\item Report selection frequency across patient resamples and seeds (Meinshausen \& B\"uhlmann 2010; Kotliar et~al.\ 2019).
\item Test ``shared'' as measurement invariance with per-patient refits (Meredith 1993); iterate with domain knowledge by target rotation to sub-signatures fixed in advance (Browne 1972).
\item Rotation itself, the team's psychology example. Varimax can make single factors identifiable when factor scores are heavy-tailed (Rohe \& Zeng 2023); the review proposes varimax-rotated PCA40 and FA as Stage~1 variants. FactoSig-promax is already an oblique rotation, but its stored loadings are the pattern matrix, and pattern and structure can rank genes differently, so the promax gene poles (allograft, apoptosis) should report both.
\end{enumerate}

{\sloppy \textbf{In the app.} The read-only factor explorer (\texttt{scripts/\allowbreak mrd\_stage0\_2/\allowbreak apps/\allowbreak old34\_factor\_\allowbreak explorer.py}, port 8051) shows every card: the fold strip, lift, poles, ADT, cell type and matches. Each atlas cell links to it. It is a separate page because the atlas is a representation-by-patient heatmap, while cards are per factor and carry caveats that do not fit in a cell. The within-patient and within-cell-type ADT views of D27 are not in the app yet (proposed W7.4). Until then, any card's D27 summary is one row of \texttt{factor\_modality\_assoc\_v1/\allowbreak d27\_card\_summary.csv} (EXP-relative under \texttt{analysis/\allowbreak scorecards/\allowbreak stage2\_diagnostic/}; one row per non-degenerate card, normal and malignant columns side by side), and every robust pair is in \texttt{d27\_adt\_robust.csv} next to it.\par}

\begin{caveatbox}
\begin{itemize}[leftmargin=*,itemsep=0pt,topsep=0pt]
\item Transductive Stage 1; one panel for the gene poles, one seed, seven patients, no intervals (P07: 22 malignant cells).
\item Poles are not differential expression: TXNIP up or IFI27 down in malignant cells is not shown.
\item Identifiability is shown for PC8 and FA7 on IFN-$\alpha$; the refits in Figure~\ref{fig:05d}c are PCA and FA at $K = 40$ only.
\item Surface protein: CLR of raw counts, no dsb; isotype controls reach $|\rho|$ about 0.23 (99th percentile); patient is nested in batch; every malignant cell is recipient-origin. The robust rule is a screen over correlated factors and 161 markers, not a test, and the malignant cell-type strata come from four patients.
\item Predicted cell types are labels transferred from a healthy reference; for malignant cells they name the nearest healthy state.
\item No causal reading. Name the axis by both poles, with orientation and penalty (L1, $C=10^{-4}$).
\end{itemize}
\end{caveatbox}

\begin{teambox}
``PC8 contrasts two groups of interferon-stimulated genes. Malignant cells score toward TXNIP, IFITM2 and TENT5A and away from IFI27, PSME2 and PSMA3. FA7 is a related axis, and either one alone gives 4.8 to 19.9\pp{} of held-out lift in five of seven patients, not only P03 and P04.''

\textbf{Likely question:} ``So TXNIP is up in malignant cells?''\\
\textbf{Answer:} Not shown yet. A pole says which genes push the score, not each gene's own shift (see Next step).

\textbf{Likely question:} ``Isn't PC8 just cell type?''\\
\textbf{Answer:} In normal cells, partly: its CD34 correlate holds in every patient ($+0.245$) but vanishes inside predicted cell types ($+0.019$). Yet malignant cells score higher within 8 of 10 types, and PC8 keeps its lift inside HSC/MPP2 alone (P02 $+21.8$, P03 $+12.7$, P07 $+10.3$, P09 $+27.1$\pp).

\textbf{Likely question:} ``Is CD99 a surface marker of the malignant pole?''\\
\textbf{Answer:} Not on this evidence: within patients it is $+0.159$ (6 of 7), below the rule fixed in advance and inside the isotype noise level for malignant cells.
\end{teambox}

\subsection*{Next step}
\begin{itemize}[itemsep=0pt,topsep=0pt]
\item W5.2: per-patient oriented gene directions and Haufe patterns for PC8 and FA7.
\item W5.5: a census of program break-up across all 56 representations, reading the D27 card summary to flag sub-axes that are cell-identity axes.
\item W5.3: advanced by D27 (all cards). To close: fix the F7b column label (it holds raw counts, not log1p). dsb-normalized ADT needs your decision on the background source.
\item Proposed W5.7: held-out lift inside predicted cell types for the 119 selected drivers, plus composition-matched comparisons, reusing the D27 strata.
\item Proposed W7.4: show the within-patient and within-cell-type views in the factor explorer.
\item W9.3: CD99 is a literature lead to verify before any text uses it; it stays out of the claims.
\end{itemize}

\subsection*{How this was computed}
{\small\raggedright
\textbf{Data and method.} The factor-card bundle \texttt{analysis/\allowbreak scorecards/\allowbreak stage2\_diagnostic/\allowbreak factor\_cards\_v1/} (EXP-relative) and the F9 tables in \texttt{feedback\_deep\_dive\_20260917/}. FA7 loadings for PC8's genes: stored FA40 \texttt{loadings.npy}, oriented by the card sign. Orientation is $\mathrm{sign}(\beta)$ of the stored, globally selected LOPO probe (L1, $C=10^{-4}$, not retuned). Entry margins come from the regularization-contract analysis. Univariate lift uses the training-fold sign. The card count is read from \texttt{factor\_cards\_index.csv} and \texttt{factor\_cards\_shortlist.csv}; the atlas totals from \texttt{old34\_d1\_atlas\_representation\_table.csv}. Identifiability: PCA explained variance from the stored IFN-$\alpha$ PCA40 \texttt{diagnostics.json} (the relative eigengap reproduces the card's), determinacy from the card index, refits from \texttt{basis\_variant\_pilot\_ifna\_v1/\allowbreak d17\_alignment.csv} (\texttt{matched\_abs\_r\_fit}, \texttt{many\_to\_one\_r2\_fit}). \emph{D27:} 1,801 non-degenerate cards $\times$ 2 labels $\times$ 170 ADT features (612,340 rows); per-cell CLR of raw counts over all 170 features; Spearman within label, pooled, within each patient ($\geq 20$ cells; median over the 7 supported patients), within supported patient $\times$ predicted cell type ($\geq 30$ cells; 71 normal and 11 malignant strata); cell-type $\eta^2$ pooled and within patient; robust rule fixed before computing; nothing refit, no $C$ policy. Artifacts \texttt{analysis/\allowbreak scorecards/\allowbreak stage2\_diagnostic/\allowbreak factor\_modality\_assoc\_v1/} (\texttt{d27\_*.csv}, \texttt{d27\_adt\_assoc\_full.parquet}, findings with a ``How to read rho here'' section).

\textbf{Code.} \texttt{scripts/\allowbreak mrd\_stage0\_2/\allowbreak workstreams/\allowbreak W5\_factor\_\allowbreak interpretation/\allowbreak run\_d18\_\allowbreak factor\_cards.py}; \texttt{\ldots/\allowbreak W5\_factor\_\allowbreak interpretation/\allowbreak run\_d27\_\allowbreak factor\_\allowbreak modality\_\allowbreak assoc.py} with \texttt{factor\_\allowbreak modality\_\allowbreak assoc.py}; \texttt{scripts/\allowbreak mrd\_stage0\_2/\allowbreak notebooks/\allowbreak stage0\_2/\allowbreak stage2\_old34\_\allowbreak feedback\_\allowbreak deep\_dives\_\allowbreak 20260917.ipynb} (section 9); \texttt{figs\_05\_pc8.py} (this report; it prints every number above); review: \texttt{scripts/\allowbreak mrd\_stage0\_2/\allowbreak workstreams/\allowbreak W9\_literature/\allowbreak 2026-10-08\_factor\_\allowbreak interpretability\_\allowbreak review.md}.

\src{2026-10-08\_d18-factor-cards; 2026-10-09\_d27-factor-modality-assoc; 2026-10-02\_f9-ifna-pc8-loadings; 2026-10-08\_d19-regularization-contract; 2026-10-08\_d17-basis-variant-pilot; 2026-10-08\_d17-composition-matched-controls; 2026-10-08\_literature-reviews; 2026-10-08\_w7-atlas-strip-factor-explorer}\par}
                            % 5. What is PC8? (D18, F9)
\labelnextsection{glance:06}% Section 06: option 2, an inductive Stage 1 (D17 basis-variant pilot, IFN-$\alpha$).
% Figures: figs_06_inductive.py (fig06a_inductive_lift, fig06b_paired_change).
\section{Option 2: does lift survive when Stage 1 never sees the held-out patient?}

\begin{quote}
``Options to address: [\ldots] 2. Exclude held-out patient from Stage 1 entirely, then project their cells onto learned factors.''
\end{quote}

\begin{answerbox}
Yes, on the one panel tested (IFN-$\alpha$; PCA and FA at $K$ = 10, 20 and 40). With Stage 1 fit without the held-out patient and that patient's cells projected onto the frozen factors, held-out lift stays positive in all 7 supported patients under per-fold tuned C or L2 (126 of 126 patient $\times$ setting pairs). Against the transductive basis, the median per-patient change is between $-1.29$ and $-0.05$\pp{} across the 18 tuned/L2 settings: always negative, so transductive LOPO numbers are slightly optimistic, but by little next to median lifts of 2.5--10.6\pp. Pooled AP is more sensitive: FA40 under tuned L2 falls from 0.1325 to 0.0794. The discriminative axis (PC8, FA7) survives the refit, and a basis fit on normal cells only carries it too. Descriptive: one panel, seven supported patients, no significance test.
\end{answerbox}

\subsection*{What we did}

Every lift in the atlas, the interactive heatmap and the factor analyses so far comes from a Stage 1 basis fit on all 60,118 cells of all 13 patients. (Lift = held-out AUPRC minus the patient's malignant prevalence $\pi$, in pp; a constant score gets AP $= \pi$, so lift 0 means no skill.) Stage 1 never sees labels, but it does see the held-out patient's cells; option 2 asks how much of the transfer depends on that.

\begin{explainbox}{Transductive vs inductive: who sees whose cells}
Take one leave-one-patient-out (LOPO) fold, P03 held out (2,327 cells, 758 malignant).

\smallskip
\begin{tabularx}{\linewidth}{@{}p{2.3cm}XXl@{}}
\toprule
Arm & Stage 1 (PCA or FA, no labels) is fit on & Stage 2 probe is trained on & Scored on \\
\midrule
Transductive (atlas) & all 60,118 cells of all 13 patients, P03 included & labelled cells of the other 12 patients & P03 \\
Inductive & the 57,791 cells of the other 12 patients; P03 is projected afterwards & same & P03 \\
Normal-only & the 56,980 normal cells of the other 12 patients & same & P03 \\
\bottomrule
\end{tabularx}

\smallskip
Transductively, P03's labels are never used, but its cells help decide where the factors point: if P03 had a distinctive malignant population, Stage 1 could spend a factor on it and hand the probe a ready-made axis. The inductive arm closes that route by fitting the basis, gene means and SDs included, without P03 and freezing it, as for a new patient. We report $\lift_h^{\text{ind}} - \lift_h^{\text{trans}}$ per patient $h$; negative means the transductive number was optimistic.
\end{explainbox}

The pilot refit IFN-$\alpha$ (92 covered genes) with PCA and FA at $K$ = 10, 20 and 40: one inductive and one normal-only basis per supported held-out patient, plus a transductive refit with the same code as the reference. Stage 2 used the standard LOPO folds (training $n$ unchanged), with L1 and L2 at the stored atlas C ($10^{-4}$) and at a C tuned per outer fold by inner LOPO. ``The 18 tuned/L2 settings'' are tuned L1, fixed L2 and tuned L2 for each method $\times$ K; fixed-C L1 is treated last.

\subsection*{What we found}

\textbf{Lift survives in every patient.} Figure~\ref{fig:06a} shows each supported patient at $K$ = 40 under tuned L1. Summaries (pooled AP / median lift / P20 lift), inductive versus transductive refit: PCA40 0.1164 / 10.55 / 1.29\pp{} versus 0.1153 / 8.91 / 1.42\pp; FA40 0.0869 / 6.69 / 1.42 versus 0.0894 / 7.03 / 1.25. Under tuned L2 the picture is the same (PCA40 median change $-0.30$\pp, FA40 $-0.44$\pp; 7 of 7 positive), though FA40's pooled AP falls more (0.1325 to 0.0794). Read paired per-patient changes rather than differences of medians: PCA40 inductive has the higher median lift (10.55 vs 8.91\pp) only because P07, the median patient, gains 1.64\pp, while 5 of 7 patients lose (median change $-0.67$\pp). P01 and P06, near zero in the figure, stay positive in all three arms (0.19--0.52 and 0.05--0.37\pp).

\begin{figure}[H]
\centering
\includegraphics[width=\linewidth]{figures/fig06a_inductive_lift.pdf}
\caption{Held-out lift per supported patient at $K$ = 40, C tuned per fold, L1, on three Stage 1 bases: transductive refit (open circle), inductive (filled circle) and normal-only (diamond). Each gray line joins one patient's three values. Look at how little the line moves next to the lift itself, and that no filled marker reaches zero. Labels mark inductive changes of at least 1.5\pp. Tick labels give each patient's malignant prevalence; P02, P03 and P04 are the highest-prevalence patients.}
\label{fig:06a}
\end{figure}

\textbf{The optimism is small and consistent in sign.} Across all 18 tuned/L2 settings (Figure~\ref{fig:06b}, top) the median per-patient change runs from $-1.29$\pp{} (FA20, tuned L1) to $-0.05$\pp{} (FA40, fixed L2) and is negative in every setting. Single patients move more, in both directions: $-3.83$\pp{} (FA40 P03, tuned L1) to $+6.13$\pp{} (PCA10 P04, tuned L1). With seven patients and C picked from a discrete grid, single-patient differences are not identifiable effects.

\begin{figure}[H]
\centering
\includegraphics[width=\linewidth]{figures/fig06b_paired_change.pdf}
\caption{Change in held-out lift versus the transductive refit for each supported patient (gray dots) in each of the 18 tuned/L2 settings; bars are the median over the seven patients. Top: inductive bases. Bottom: normal-only bases. Look at where the bars sit relative to zero (just below it everywhere in the top row) and at the spread of the dots, which is single-patient noise.}
\label{fig:06b}
\end{figure}

\textbf{The axis survives without the held-out patient.} The reference axis is the factor with the largest standardized malignant-vs-normal gap (PC8, FA7 at every K; malignant cells score lower on PC8, higher on FA7). \emph{Matched $|r|$}: Hungarian matching of refit factors to it on $|$correlation$|$ of cell scores; \emph{many-to-one $R^2$}: the share of it explained by all K refit factors. At $K$ = 40 the inductive PC8 has matched $|r|$ 0.97--0.9995 in 6 of 7 bases, and FA7 0.97--0.9997 in 6 of 7. The exceptions keep the direction but spread it over neighbouring factors: P01-out PCA ($|r|$ 0.59, $R^2$ 0.986) and P09-out FA ($|r|$ 0.75, $R^2$ 0.981); in that P09 fold the tuned L1 probe gives the matched factor zero weight and still reaches 6.69\pp{} through its neighbours.

\textbf{Normal cells alone carry the axis.} A basis fit on the normal cells of all patients reproduces PC8 at $|r|$ 0.9937 and FA7 at 0.9937 ($K$ = 40). The per-fold normal-only bases, which also leave out the held-out patient, keep lift positive in 126 of 126 patient $\times$ setting pairs, with median changes from $-1.00$ to $+0.25$\pp{} (Figure~\ref{fig:06b}, bottom); at $K$ = 40 under tuned L1, PCA 0.0983 / 6.98 / 1.36\pp{} and FA 0.1126 / 7.57 / 1.17\pp. The direction that separates malignant from normal cells is one along which normal cells already vary; malignant cells differ in where they sit on it.

\textbf{At the fixed atlas C, inductive bases look worse; that is the L1 entry threshold.} Under balanced L1 at fixed C, a standardized factor can enter only if its training malignant-vs-normal mean gap exceeds $4/(C\,n_{\text{train}})$ SD (0.678\,SD at $C = 10^{-4}$ with P04 held out); the \emph{entry margin} is the gap divided by that threshold. All margins below 1 means an all-zero probe, AP $= \pi$, lift 0 by construction (not ``no signal''). At $C = 10^{-4}$ the inductive PCA40 median lift falls from 11.10 to 0.05\pp{} (positive in 4 of 7 patients instead of 5). One fold moves the median: P07's largest margin drops from 1.075 (transductive refit) to 0.928, so its lift goes from 11.10\pp{} to 0; with C tuned in that fold ($10^{-3}$) it is 10.55\pp. Normal-only bases have lower margins in 13 of 14 PCA40/FA40 folds and are all-zero at the fixed C in 5 of 7 (PCA40) and 7 of 7 (FA40) folds, yet positive in 7 of 7 once C is tuned. An inductive atlas rerun must retune C per fold.

\textbf{Related arm: patient-balanced Stage~1 (option 1.2).} Weighting each patient equally in the Stage 1 fit (all cells, so still transductive) rotates PC8 more than FA7 (matched $|r|$ 0.90 vs 0.98 at $K$ = 40), with no clear loss under tuned L1: PCA40 0.1213 / 9.25 / 2.69\pp{} (median change $+1.35$\pp, no patient lower), FA40 0.0863 / 6.85 / 1.20\pp{} (median paired change $+0.07$\pp, though all three summaries are slightly below the refit's 0.0894 / 7.03 / 1.25). Under tuned L2, FA40 loses a median 1.13\pp{} (5 of 7 patients lower). Single patients swing more ($-11.74$ to $+12.77$\pp{} across the 18 settings): read it as ``no clear loss on IFN-$\alpha$'', not a gain.

\begin{caveatbox}
\begin{itemize}
\item One panel (IFN-$\alpha$, 92 genes), PCA and FA only, $K$ = 10, 20, 40, seed 42. The 47 atlas-selected representations, FactoSig included, are untested; that needs a GPU or CPU screen job.
\item Seven supported patients, no significance test; single-patient changes (FA40 P03 $-3.83$\pp, PCA10 P04 $+6.13$\pp) are not identifiable effects.
\item C tuning is not fully inductive: the inner LOPO folds reuse the outer fold's inductive basis, whose fit included the inner held-out patients.
\item Same cohort and processing: no evidence yet for new data such as Yu-Hsiang's patients; cell-level lift is not recurrence-risk improvement.
\item The reference is the transductive refit (exact projection). Stored PCA scores use a slightly different convention, so stored-basis PCA40 lifts differ from it by up to 1.23\pp{} (tuned L1; 0.32\pp{} at the atlas's fixed-C L1); FA is identical.
\item Nothing here says whether P02, P03 and P04 shape the loadings; that is the leave-patients-out analysis.
\end{itemize}
\end{caveatbox}

\begin{teambox}
``On IFN-$\alpha$, refitting Stage 1 without the held-out patient keeps lift positive in all 7 patients and costs 0.05 to 1.29\pp{} at the median across 18 settings, so our LOPO numbers are slightly optimistic, not an artifact of the held-out patient's cells.''

\textbf{Likely question:} ``So should every atlas number be rerun inductively?''\\
\textbf{Answer:} For anything presented as transfer to new patients, yes, with C retuned per fold: at the fixed atlas C one small margin shift zeroes a fold (PCA40 P07, 11.10\pp{} to 0). Until that run covers all 47 representations, call atlas numbers label transfer on a transductive basis; on IFN-$\alpha$, with tuned C or L2, that overstates median per-patient lift by at most about 1.3\pp.
\end{teambox}

\subsection*{Next step}
\begin{itemize}
\item W2.6: inductive Stage~1 for the 47 atlas-selected representations, with fold-specific storage, per-fold tuned C and L2 alongside (GPU or CPU, in a user screen session).
\item W1.3: rerun the patient factor-overlap analysis (D14) on the inductive and normal-only bases.
\item W2.7: audit stored PCA scores against the exact projection across the atlas before projecting new patients onto stored bases.
\end{itemize}

\subsection*{How this was computed}
{\small\raggedright
Data: IFN-$\alpha$ hallmark panel (92 covered genes, z-scored), 60,118 MRD cells, 13 patients; 7 supported held-out patients. Stage 1: scikit-learn PCA and FactorAnalysis (randomized, seed 42), $K$ = 10/20/40, gene standardization on the fit cells, held-out cells projected with the fitted transform. Stage 2: balanced logistic probes on train-fold standardized scores (L1 proximal Newton, L2 L-BFGS); C fixed at $10^{-4}$ or tuned per outer fold by inner LOPO over the six supported training patients (median inner lift; 17-value grid $10^{-4}$--$10^{4}$). Estimand: per-patient lift, variant minus transductive refit, median over seven patients; pooled AP, median and P20 lift quoted together. Anchor: the stored-basis arm reproduces 84 of 84 stored per-patient atlas lifts within 0.1\pp{} with the same nonzero support. Code: \texttt{scripts/\allowbreak mrd\_stage0\_2/\allowbreak workstreams/\allowbreak W2\_stage1\_\allowbreak basis\_variants/\allowbreak run\_d17\_\allowbreak basis\_variant\_\allowbreak pilot.py} (with \texttt{stage1\_refit.py}); C policy from \texttt{scripts/\allowbreak mrd\_stage0\_2/\allowbreak workstreams/\allowbreak W3\_dr\_\allowbreak value\_\allowbreak regularization/\allowbreak regularization\_\allowbreak contract.py}; figures from \texttt{figs\_06\_inductive.py} in this report's directory. Artifacts (EXP-relative): \texttt{analysis/\allowbreak scorecards/\allowbreak stage2\_\allowbreak diagnostic/\allowbreak basis\_variant\_\allowbreak pilot\_ifna\_v1/} (\texttt{d17\_stage2\_\allowbreak by\_patient.csv}, \texttt{d17\_stage2\_\allowbreak summary.csv}, \texttt{d17\_alignment.csv}). \src{2026-10-08\_d17-basis-variant-pilot (D17); entry rule from 2026-10-08\_d19-regularization-contract (D19)}\par}
                       % 6. Option 2: inductive Stage 1 (D17)
\labelnextsection{glance:07}\section{Learning curves: does more data raise AUPRC?}
\label{sec:learning-curves}

\begin{quote}
``Plot AUPRC vs number of cells (learning curve / marginal benefit curve); useful for grant applications to justify sequencing more patients; randomly subsample to estimate the curve.''
\end{quote}

{\small\color{inkmuted}\emph{Reading guide (v2).} Why $C$ is retuned when the number of training cells or patients changes: the first paragraph of ``What we did'', drawn in Figure~\ref{fig:07a}b and in the gray lines of Figure~\ref{fig:07b}. The L1 entry rule worked through for one fold (P02 held out, PCA40 vs FA40): ``One fold worked through'', right after it.\par}

\begin{answerbox}
Yes, we can draw it, and we now have two versions, but neither yet justifies sequencing more patients. Adding \emph{cells} from the same 12 training patients (C fixed) raises median apoptosis lift from 7.65\pp{} at 500 training cells to 11.52\pp{} on the full fold. Adding \emph{patients} (C retuned for every training set, because a fixed C is a stronger penalty with fewer training cells; all 1,764 subset units) shows lift is already positive with one supported training patient (median 4.45--8.79\pp{}; the six low-malignancy patients are always in training), and going from two to six supported training patients adds a median 0.30--2.34\pp{} per held-out patient, with a 95\% CI that includes 0 for three of four representations. Descriptive only: one frozen transductive basis, seven supported patients, no extrapolation to new patients.
\end{answerbox}

\subsection*{What we did}

\textbf{First: why $C$ is retuned whenever the training cells or patients change.} In scikit-learn the data term of the logistic loss is a sum over training cells and the penalty is not, so the same $C$ is a stronger penalty when there are fewer cells. Under class-balanced L1 a standardized factor enters the probe only if its training class-mean gap exceeds about $4/(C\,n_{\text{train}})$\,SD (the L1 entry rule; the IFN-$\alpha$ PCA40 vs FA40 section derives it). A curve drawn at one fixed $C$ therefore partly measures the penalty, not the data.

\emph{The IFN-$\alpha$ example} (Figure~\ref{fig:07a}b): at the atlas $C = 10^{-4}$ the bar is 8\,SD with 5,000 training cells and 1.6\,SD with 25,000, and it falls below PC8's class gap (0.704--0.770\,SD) only near 52,000--57,000 cells. So the fixed-$C$ cell-depth curve of IFN-$\alpha$ PCA40 is exactly 0 at every subsample up to 25,000 cells (an all-zero probe scores AP $= \pi$ by construction) and jumps to 11.08\pp{} on the full fold: flat because of the penalty, not because the signal is missing. The same holds across patients. On the same training subsets, PCA40 at the fixed atlas $C$ climbs from 0 at one supported training patient to 11.08\pp{} at six, with all-zero probes in 100\% and then 29\% of the units; with $C$ retuned it is already 4.63\pp{} at one patient and ends at 8.71\pp{} (at six patients the atlas $C$, picked by pooled AP on these same held-out folds, scores higher). Most of the fixed-$C$ ``learning'' was the bar falling as $n$ grew.

\emph{One fold worked through} (the same arithmetic gives the PCA40 vs FA40 gap of section \emph{IFN-$\alpha$: why does FA40 reach only P03 and P04?}). With P02 held out, Stage~2 trains on $n_{\text{train}} = 58{,}656$ cells, so at $C = 10^{-4}$ the bar is $4/(10^{-4} \times 58{,}656) = 0.682$\,SD. PC8's class-mean gap is 0.740\,SD, an entry margin of $0.740/0.682 = 1.085$: it enters, and the PCA40 probe scores 15.94\pp{} in P02. FA7's gap is 0.659\,SD, a margin of 0.967: no FA40 factor enters, so FA40 scores 0 by construction, although FA7 used alone gives 18.06\pp{} there. Because the bar depends on $C \times n_{\text{train}}$, FA7 would enter with a $C$ about 3.5\% higher, or at the same $C$ with about 3.5\% more training cells: whether a factor is used depends on the penalty and on $n$, not on the factor alone.

\emph{Hence the rule we follow:} whenever the training set changes (subsets, subsampling, re-weighting, downsampling, new patients), $C$ is retuned for each training set by inner leave-one-patient-out over its training patients, and L2 is reported alongside. The patient-count curve below does this; the cell-depth curve, drawn earlier, keeps the atlas $C$ and is read with that caveat.

Two curves, two questions, both on the same four representations of three gene programs: IFN-$\alpha$ PCA40 and FA40 (atlas penalty L1, $C = 10^{-4}$), apoptosis FactoSig-promax40 (L2) and allograft-rejection FactoSig-promax40 (elastic net).
\begin{itemize}
\item \textbf{Cell depth.} Keep all 12 other patients in training, randomly subsample their cells (500 to 25,000, 3 draws, plus the full fold), and keep penalty and C at the atlas choice. Is the probe data-starved at a fixed patient mix?
\item \textbf{Patient count.} For each held-out supported patient, train on every subset of $k$ of the other six supported patients ($k = 1,\dots,6$; 63 subsets; the six unsupported patients always in), retuning penalty (L1 or L2) and C for every training set. $k = 6$ is the ordinary LOPO fold. This is the grant question: what does one more patient buy?
\end{itemize}

\begin{explainbox}{Four terms used in this section}
\emph{Lift} = the held-out patient's AUPRC (average precision, AP) minus its malignant prevalence $\pi$, in percentage points (pp). A constant score has AP $= \pi$ exactly (no skill), so lift 0 means no ranking skill.

\emph{Fixed-C entry threshold.} The bar $4/(C \cdot n_{\text{train}})$\,SD of the first paragraph. A probe with no factor over the bar is all-zero and scores lift 0 \emph{by construction}, not for lack of signal; retuning C per training set by inner leave-one-patient-out removes this (the patient-count curve does; the cell-depth curve does not).

\emph{Transductive basis.} Stage 1 was fit once on all 60,118 cells, held-out patient included; both curves refit only the Stage 2 probe, so they measure label transfer on a basis that has seen everyone. An \emph{inductive} curve would refit the basis on the $k$ training patients.

\emph{Paired gain.} Per held-out patient, lift in its LOPO fold ($k = 6$) minus its mean lift over the 15 two-patient subsets; median over the seven patients, 95\% CI by resampling patients (2,000 draws). It starts at $k = 2$ because at $k = 1$ inner LOPO is impossible and C is tuned by cell-level CV, a different protocol.
\end{explainbox}

\subsection*{What we found}

\textbf{Cell depth, fixed C (Figure~\ref{fig:07a}).} Apoptosis is at 7.65\pp{} median lift with 500 training cells (about 13 malignant), 9.28\pp{} at 1,000 and 11.52\pp{} on the full fold: nearly saturated by about 1,000 cells. Allograft (elastic net, L1 ratio 0.5) is 0 through 2,000 cells because its probe is all-zero in all 21 units there, the same entry-threshold collapse as IFN-$\alpha$ (its L1 part needs a class gap above $4 \times 0.5/(C\,n)$, about 2.2\,SD at 2,000 cells); it then reaches 1.67, 3.14, 8.81 and 12.56\pp{} at 5k, 10k, 25k and full. Both IFN-$\alpha$ L1 probes are exactly 0 at every subsample up to 25,000 cells; on the full fold PCA40 reaches 11.08\pp{} and FA40 stays at 0. The reason is the penalty: at $C = 10^{-4}$ the entry bar is 8\,SD at 5,000 cells and 1.6\,SD at 25,000, and it falls below PC8's class gap (0.704--0.770\,SD) only at about 52,000--57,000 cells, close to the full-fold size (median 57,791).

\begin{figure}[H]
\centering
\includegraphics[width=\linewidth]{figures/fig07a_cell_depth.pdf}
\caption{Cell-depth curve: fixed 12-patient training mix, fixed atlas C. (a) Median held-out lift over the seven supported patients (each patient's mean over 3 draws); malignant training cells in brackets. Look at apoptosis (aqua circles), already positive at 500 cells, against allograft (aqua diamonds), all-zero through 2,000 cells, and both IFN-$\alpha$ L1 curves, pinned at 0 until the full fold. (b) The L1 entry threshold $4/(C\,n)$ at $C = 10^{-4}$; the blue band is PC8's class-mean gap. The IFN-$\alpha$ probe can become non-zero only right of the gray crossing.}
\label{fig:07a}
\end{figure}

\textbf{Patient count, C retuned (Figure~\ref{fig:07b}, Table~\ref{tab:07}).} No representation collapses, and lift is positive with a single supported training patient for all four. From $k = 2$ to 6 the median curves rise by 1.58 to 2.83\pp{} (the difference of the two medians; the paired gain we quote, which compares each held-out patient with itself, is 0.30--2.34\pp) and are not monotone (PCA40 peaks at 9.30\pp{} at $k = 4$ and ends at 8.71\pp{}). Each median's 95\% CI is 7.7 to 37.6\pp{} wide, several times that rise. The lower tail (P20 lift, the 20th percentile over patients) stays between 1.18 and 3.29\pp{} at every $k$. Only allograft's paired-gain CI excludes 0.

\begin{figure}[H]
\centering
\includegraphics[width=\linewidth]{figures/fig07b_patient_count.pdf}
\caption{Patient-count curve (full design, 1,764 units). Coloured line (colour = representation): median over held-out patients of each patient's mean lift across all subsets of size $k$, C and penalty retuned per training set (open $k = 1$ point: cell-level CV tuning). Band: 95\% patient-bootstrap CI of that median. Gray: the same subsets at the fixed atlas C, which reproduces the cell-depth collapse (FA40 all-zero in every unit at $k = 1$--4 and in 90\% and 71\% at $k = 5$ and 6; PCA40's median climbs from 0 to 11.08\pp{} only because the bar falls as $n$ grows). With C retuned, FA40 has no all-zero unit and PCA40 only 5\% of its $k = 2$ units. Median malignant training cells at each $k$ are in the note of Table~\ref{tab:07}. Look at how flat the coloured lines are relative to the band.}
\label{fig:07b}
\end{figure}

\begin{table}[H]
\centering
\small
\caption{Patient-count curve (full design), tuned arm: median held-out lift (pp) over the seven held-out patients by number of supported training patients $k$; P20 lift range over $k$; paired gain from $k = 2$ to 6 (median, 95\% patient-bootstrap CI, patients with gain $> 0$). Pooled AP is not computed for these curves.}
\label{tab:07}
\setlength{\tabcolsep}{3.6pt}
\begin{tabular}{@{}lrrrrrrcl@{}}
\toprule
 & \multicolumn{6}{c}{median lift at $k$ =} & P20 & gain $k$ = 2 $\to$ 6 \\
\cmidrule(lr){2-7}
representation & 1 & 2 & 3 & 4 & 5 & 6 & range & median [95\% CI], $n_{>0}$ \\
\midrule
IFN-$\alpha$ PCA40 & 4.63 & 5.88 & 7.83 & 9.30 & 8.77 & 8.71 & 1.18--1.61 & 1.67 [$-0.44$, 2.89], 5/7 \\
IFN-$\alpha$ FA40 & 4.45 & 4.77 & 5.51 & 6.66 & 6.20 & 7.03 & 1.27--1.96 & 0.30 [$-0.53$, 4.18], 4/7 \\
Apoptosis FactoSig-promax40 & 5.34 & 8.68 & 8.27 & 9.96 & 10.31 & 10.26 & 1.26--2.25 & 1.14 [$-0.53$, 2.38], 5/7 \\
Allograft FactoSig-promax40 & 8.79 & 7.63 & 8.13 & 8.63 & 9.32 & 10.12 & 2.22--3.29 & 2.34 [0.14, 2.73], 6/7 \\
\midrule
\multicolumn{9}{@{}l}{\footnotesize Median malignant training cells at $k = 1,\dots,6$: 254, 497, 740, 982, 1,224, 1,467.} \\
\bottomrule
\end{tabular}
\end{table}

\textbf{The gain is uneven across patients (Figure~\ref{fig:07c}).} Per-patient gains run from $-2.80$\pp{} (P02, FA40) to $+13.32$\pp{} (P03, apoptosis: 28.54 to 41.86\pp{}); P01 and P06 stay within $-0.53$ to $+1.14$\pp{}. At $k = 6$, per-patient lift ranges from 0.06\pp{} (P06, PCA40) to 41.86\pp{} (P03, apoptosis): which patient is held out matters far more than how many are in training. Which patients are \emph{added} matters too (descriptive, since subsets also differ in their other members): at matched $k$, subsets containing P02 have 5.54\pp{} higher allograft lift, and subsets containing P09, a large low-prevalence patient, 2.70\pp{} lower apoptosis lift (medians over held-out patients).

\begin{figure}[H]
\centering
\includegraphics[width=\linewidth]{figures/fig07c_patient_gain.pdf}
\caption{Per-patient paired gain (tuned arm): lift in the LOPO fold minus the mean over the 15 two-patient subsets. Each lollipop is one held-out patient; the vertical line and band are the median gain over patients and its 95\% patient-bootstrap CI (numbers above each panel). Look at the bracketed highest-prevalence patients P02--P04: they hold the largest gain in every panel and the only IFN-$\alpha$ losses beyond 1\pp{}, while P01 and P06 barely move.}
\label{fig:07c}
\end{figure}

\textbf{No reliable scaling law.} By a rule fixed before looking, an inverse power law $a - b\,k^{-c}$ fitted on $k \le 4$ is reportable only if it converges, has no parameter at a bound, rises by $\geq 0.5$\pp{} and predicts $k = 5$ inside the observed CI. Only apoptosis passes (predicted 9.51\pp{}, observed 10.31 [1.96, 35.80]), and its asymptote, 9.80\pp{}, is below the observed $k = 5$ and 6 medians: a plateau, not growth. With CIs this wide, passing is a weak check.

\textbf{What would make this grant-grade.}
(1)~\emph{An inductive basis per subset}: refit Stage 1 on each training subset, project the held-out patient and redraw the curve on the same subsets. This is feasible; for IFN-$\alpha$ the $k = 6$ endpoint already has an inductive counterpart (the inductive pilot; read it as paired changes against the same-code transductive refit: PCA40 tuned-L1 median per-patient change $-0.67$\pp, 5 of 7 patients lower; section \emph{Option 2}).
(2)~\emph{New patients as external points}: score Yu-Hsiang's patients with models trained on this cohort and compare them with the curve.
(3)~\emph{More patients with enough malignant cells}: with seven supported patients $k$ stops at 6 and every CI rests on seven clusters; a scaling claim needs more held-out patients, not more cells per patient.
A sentence we can defend in a grant today:
\begin{quote}
``In seven post-transplant MRD patients with enough malignant cells, a knowledge-prior factor classifier whose malignant training examples came from a single other patient already ranked malignant cells above their prevalence in held-out patients (median lift 4.5--8.8 percentage points across four representations of three gene programs; label-free factors learned from all patients' cells). Increasing training from two to six such patients changed median lift by 0.3--2.3 points, with confidence intervals that include zero for three of the four. The cohort is too small to estimate how performance scales with more patients; the proposed sequencing would let us measure that curve on held-out patients.''
\end{quote}

\begin{caveatbox}
\begin{itemize}
\item That sequencing more patients raises AUPRC: the curve stops at six training patients and nothing is extrapolated.
\item An inductive learning curve: the basis saw every patient at every $k$; a per-subset basis could move the level or slope either way.
\item A patient-count effect from $k = 1$: tuning switches from cell-level CV to inner LOPO between $k = 1$ and 2.
\item A scaling law: one of four fits passes a weak validation, and it describes a plateau.
\item Patients separated from malignant cells: adding a patient adds its malignant cells (P03 carries 758, P09 375).
\item That the cell-depth curve measures data alone: at fixed C the penalty weakens as $n$ grows, for L2 and elastic net as well as L1; the cell-depth curve with C retuned per subsample would separate the two.
\item Recurrence risk: lift is a cell-level ranking metric.
\end{itemize}
Descriptive results on a transductive Stage 1 basis over seven supported patients; they do not establish inductive deployment, causal biology or recurrence risk.
\end{caveatbox}

\begin{teambox}
``With C retuned for each training set, the classifier already transfers to a held-out patient after training on one supported patient (the six low-malignancy patients are always in training), and going from two to six supported patients adds only a few pp with CIs that mostly include zero, so we cannot yet claim that more patients help.''

\textbf{Likely question:} ``Why retune C for every training set?''\\
\textbf{Answer:} scikit-learn's C does not scale with the number of cells, so the same C is a stronger penalty on a smaller training set. At a fixed C the IFN-$\alpha$ curve is 0 up to 25,000 cells because no factor clears the bar, not because the signal is missing.

\textbf{Likely question:} ``So the curve is useless for the grant?''\\
\textbf{Answer:} No. It shows feasibility (transfer works with very few patients) and that seven patients are too few to measure scaling, which is itself a case for sequencing more. It does not show that more patients raise accuracy; that needs the inductive per-subset curve and new patients as external points.
\end{teambox}

\subsection*{Next step}
\begin{itemize}
\item W6.6: rerun the patient-count curve on an inductive Stage~1 basis refit per training subset and compare it with this curve on the same subsets.
\item W6.5 (with W8.1): score Yu-Hsiang's new patients as external points against the curve.
\item W6.2 and W6.3: the cell-depth curve with C retuned per subsample, and a patients $\times$ malignant-cells-per-patient design; then the W6.4 grant note.
\end{itemize}

\subsection*{How this was computed}
{\small\raggedright
\emph{D15.} Stored transductive Stage 1 scores (K = 40, seed 42). Per held-out supported patient, label-stratified subsets of the other 12 patients' cells (3 draws) plus the full fold (42,992 to 58,990 cells); scaler and probe refit at the atlas penalty and C (C policy: global, fixed). Full-fold refits match the D1 atlas within 0.000164 AUPRC. Code \texttt{scripts/\allowbreak mrd\_stage0\_2/\allowbreak stage2\_supervised/\allowbreak cell\_count\_learning\_curve.py}; artifacts \texttt{analysis/\allowbreak scorecards/\allowbreak stage2\_diagnostic/\allowbreak cell\_count\_learning\_curve\_v1/}. \src{D15; devlog 2026-10-02\_d15-cell-count-learning-curve}

\emph{D15b, full design.} Same scores; 4 representations $\times$ 7 held-out patients $\times$ 63 subsets = 1,764 units. C policy: per training set, penalty in \{L1, L2\} and C tuned by inner LOPO through the D19 regularization contract (1,008 tuning runs; cell-level stratified 5-fold CV at $k = 1$); fixed-atlas and $n$-scaled C as reference arms. Per $(h, k)$ mean lift over subsets, then median and P20 over held-out patients; patient-cluster bootstrap, 2,000 draws. At $k = 6$, 84 of 84 reference checks reproduce D1, D15 and D19 (max $|\Delta\,\mathrm{AUPRC}| = 1.64 \times 10^{-4}$). CPU, screen \texttt{W6\_d15b}, 12 workers, 565.6\,s wall. Code \texttt{scripts/\allowbreak mrd\_stage0\_2/\allowbreak workstreams/\allowbreak W6\_learning\_curves/\allowbreak run\_d15b\_patient\_count\_curve.py}; artifacts \texttt{analysis/\allowbreak scorecards/\allowbreak stage2\_diagnostic/\allowbreak patient\_count\_learning\_curve\_v1/} (\texttt{d15b\_full\_*}); figure code \texttt{figs\_07\_curves.py} in this report. \src{D15b; devlog 2026-10-08\_d15b-full-design}

\emph{Worked fold (D19).} The P02 row of \texttt{analysis/\allowbreak scorecards/\allowbreak stage2\_diagnostic/\allowbreak regularization\_contract\_v1/\allowbreak d19\_anchor\_ifna\_l1\_entry\_check.csv} (training cells, entry bar, PC8 and FA7 class-mean gaps on training-fold standardized scores, margins, critical $C$ multiplier 1.035) and \texttt{d19\_ifna\_pca40\_vs\_fa40\_by\_patient.csv} (stored PCA40 L1 lift; FA40 L1 at $1.05\times C$, where FA7 is the only factor). \src{D19; devlog 2026-10-08\_d19-regularization-contract}\par
}
                        % 7. Learning curves (D15, D15b)
\labelnextsection{glance:08}\section{The three scientific questions (the paper spine)}
\label{sec:three-questions}

\begin{quote}
\emph{Current scientific questions:}
1.~Shared or patient-specific? Do the same factors drive lift across the high-cell patients, or different factors for each?
2.~How strong is the evidence? How much variation do cohort-shared versus patient-specific factors explain, and with what effect size and statistical significance?
3.~What does dimensionality reduction add? Compare with direct-gene classification to establish the method's value.
\end{quote}

This section pulls the earlier results into the three questions a paper would be built on. For each question it gives the current answer, the evidence that only this section shows, short pointers to the sections that hold the rest, and what a publication claim still needs.

\begin{answerbox}
\textbf{Q1.} Tested directly on one program only. On IFN-$\alpha$ the same axis carries held-out lift in five of seven supported patients, including two low-prevalence ones. Across the atlas, lift above 0 reaches every supported patient, but lift above 5\pp{} occurs in 64--80\% of the 640 biological DR cells for P02, P03 and P04 and in at most 23\% for the other four, so shared transfer and prevalence are not yet separated.
\textbf{Q2.} The effect size is well defined (lift above constant-score no-skill AP $=\pi$), but there is no inferential framework yet: no calibrated null, no patient-level intervals on the atlas, no nested selection.
\textbf{Q3.} It depends on the estimand and on how the winning DR cell is chosen. When the winner is picked by pooled AP, best DR beats direct genes on pooled AP in 38/41 programs, on median lift in 16/41 and on P20 lift in 17/41. When it is picked by median lift, it wins on median lift in 28/41 programs (30/47 with the control bags). Large gains are confined to a few programs, and broad control bags favour direct genes.
All of this is descriptive, on a transductive Stage~1 basis, over seven supported patients.
\end{answerbox}

\begin{explainbox}{Terms used in this section}
\begin{itemize}
\item \textbf{Lift}: held-out AUPRC minus that patient's malignant prevalence $\pi$, in percentage points (pp). A constant score has AP exactly $\pi$, so lift 0 means no skill. An L1 probe that keeps no factor (an \emph{all-zero} probe) is such a constant score.
\item \textbf{Median and P20 lift}: patient-equal summaries over the seven supported patients (the median and the 20th percentile, the weak-patient tail). \textbf{Pooled AP}: one AP over all held-out cells pooled across folds; it weights patients by cell count and is dominated by the high-prevalence patients.
\item \textbf{Transductive vs inductive}: the stored Stage~1 bases were fit on all 60,118 cells, including the held-out patient's (transductive), so LOPO on them is label transfer. An inductive basis is fit without the held-out patient.
\item \textbf{DR cell}: one panel at one Stage~1 method and $K$ (4 methods $\times$ 4 $K$, so up to 16 per panel). \textbf{$\eta^2$}: the share of a factor's within-label score variance explained by patient identity.
\end{itemize}
\end{explainbox}

\subsection{Q1. Shared or patient-specific?}

\textbf{Current answer.} Tested directly on one program only. On IFN-$\alpha$ the same axis (PC8, or the related FA7) carries held-out lift in five of seven supported patients, including two low-prevalence ones (P07 11.08\pp, P09 5.47\pp), and survives a basis fit without the held-out patient. Across the atlas, lift above 0 reaches every supported patient (69--84\% of the 640 biological DR cells), but lift above 5\pp{} occurs in 64--80\% of cells for P02, P03 and P04 and at most 23\% for the others, so shared transfer and prevalence are not yet separated. Patients' factor-score distributions overlap on 40 of 41 programs; that permits one basis but does not show that the same factors drive lift.

\textbf{Evidence only this section shows.}
\begin{itemize}
\item \emph{Transfer is common; strong transfer is not} (Figure~\ref{fig:08a}). Of the 783 atlas cells, 506 have lift above 0 in all seven patients, but only 84 reach lift above 5\pp{} in at least five patients and 2 in all seven. The median within-representation Spearman $\rho$ between lift and prevalence is 0.821 over 643 representations (nominal, not a test). \src{D1f, D6; 2026-09-11\_d1-atlas-783-d1e-d1f}
\item \emph{Low-support patients.} P01 and P06 have 157 and 232 within-patient DE genes, yet their best held-out lift is only 6.3 and 5.7 AP points. \src{D13}
\end{itemize}
\textbf{Elsewhere in this report.} The IFN-$\alpha$ axis and its per-patient lifts: the section on what PC8 is. Its survival without the held-out patient: the section on option~2. Overlap of factor-score distributions (40 of 41 programs aggregate; the only patient-private axes are P03's, in broad gene bags): the section on aggregating patients. Leverage of P02, P03 and P04 on the Stage~1 axis (IFN-$\alpha$, allograft rejection, apoptosis; PCA and FA): beyond composition-matched drops on allograft and apoptosis, mostly through P03, but on no panel beyond the contrast patients, and P01 moves the axes most (D26): sections~\ref{glance:02} and \ref{glance:02b}.

\begin{figure}[H]
\centering
\includegraphics[width=\linewidth]{figures/fig08a_patient_concentration.pdf}
\caption{Where the lift sits. For each supported patient, ordered by malignant prevalence $\pi$, the bars show the share of the 640 biological DR atlas cells with held-out lift above 0 (light) and above 5\pp{} (dark). The dark bars are 64--80\% for P03, P02 and P04, at most 23\% for the rest and under 1\% for P06 and P01. The light bars (69--84\%) show that some transfer above $\pi$ reaches every patient. Each cell uses one global $C$, chosen on pooled LOPO AP, on the transductive basis.}
\label{fig:08a}
\end{figure}

\textbf{What a publication claim still needs.} Prevalence- and support-matched lift, to separate shared transfer from prevalence. The same factor-level transfer test as IFN-$\alpha$ on other programs. Inductive bases for the selected representations beyond IFN-$\alpha$. The leverage test on FactoSig bases and on the other 38 panels (W2.2, W2.8). A calibrated overlap test. New patients as an external test.

\subsection{Q2. How strong is the evidence?}

\textbf{Current answer.} The effect sizes are good, with a clean zero point, and several robustness checks exist. Strength of evidence in the statistical sense does not: nothing has been tested against a null model, the atlas has no patient-level intervals, and the best configurations are chosen on the folds they are scored on. ``How much variation do shared and patient-specific factors explain'' has three different answers, one for each kind of variance (Table~\ref{tab:08-variance}). Report them separately.

\begin{table}[H]
\centering
\small
\caption{Three variance estimands, with IFN-$\alpha$ $K$ = 40 as the example and the cohort-wide value where it exists. They answer different questions and are never added or compared directly.}
\label{tab:08-variance}
\begin{tabularx}{\linewidth}{>{\raggedright\arraybackslash}p{3.0cm} >{\raggedright\arraybackslash}X >{\raggedright\arraybackslash}X}
\toprule
Estimand & What it measures & Value \\
\midrule
Label-free Stage~1 variance & Share of the panel's gene variance the basis captures (PCA), or a common-variance proxy (FA) & IFN-$\alpha$: PCA40 captures 54\% of the gene variance; its first 8 PCs carry about 21\% (38\% of the captured part); FA40 proxy 16\% \\
Per-factor patient effect ($\eta^2$), IFN-$\alpha$ & Share of one factor's within-label score variance explained by patient & PC8 0.010 (normal) / 0.031 (malignant); FA7 0.015 / 0.025; largest malignant $\eta^2$ 0.114 (PCA40), 0.121 (FA40) \\
Per-factor patient effect ($\eta^2$), all programs & The same, over all 41 programs, PCA and FA, $K$ = 5--40 (5,838 non-degenerate factors) & Median 0.024 (normal) and 0.051 (malignant); largest 0.352 and 0.619 \\
Stage~2 attribution & Held-out lift of one factor alone & IFN-$\alpha$ PC8: 11.08--16.27\pp{} in P02, P03, P04, P07; 5.47\pp{} in P09 \\
\bottomrule
\end{tabularx}
\end{table}

\textbf{Evidence.}
\begin{itemize}
\item \emph{Effect size.} Constant-score no-skill AP $=\pi$ holds exactly: all 757 all-zero (cell, fold) pairs have $|\lift|\leq5.6\times10^{-17}$. Median lift is above 0 in 670 of 783 atlas cells and P20 lift in 639. \src{D6, D1; 2026-10-08\_w7-atlas-strip-factor-explorer}
\item \emph{Significance: none yet.} The overlap test's permutation $p$ is at its floor ($1/26=0.038$, 25 permutations) in all 47 normal strata, so its calls rest on uncalibrated magnitude gates. The prevalence correlations are nominal. \src{D14, D6}
\item \emph{Selection optimism, where measured.} With $C$ tuned per training set, the full-LOPO medians are 8.71 (IFN-$\alpha$ PCA40), 10.26 (apoptosis) and 10.12\pp{} (allograft), against 11.08, 11.52 and 12.56\pp{} at the stored atlas $C$, which saw the held-out patients. FA40 goes the other way (7.03 vs 0.00\pp) because the atlas $C$ zeroes it. \src{2026-10-08\_d15b-full-design}
\item \emph{Robustness.} 528 of 783 atlas cells select the lowest $C$ on the grid. 136 L1/elastic-net folds in 49 cells sit within $\pm10\%$ of the entry threshold, and crossing it changes lift by $\geq5$\pp{} in 41 of 131 evaluable folds (maximum 39.06\pp). \src{D19; 2026-10-08\_d19-regularization-contract}
\end{itemize}
\textbf{Elsewhere in this report.} The only patient-level intervals so far (the 7-cluster bootstrap of the patient-count curve: gains from two to six training patients of 0.30--2.34\pp, CIs including 0 for three of four representations): the learning-curve section. Transductive optimism on IFN-$\alpha$ ($-1.29$ to $-0.05$\pp{} per patient at the median): the section on option~2. Why FA40 is all-zero at the atlas $C$: the IFN-$\alpha$ section.

\textbf{What a publication claim still needs.} A matched random gene-set null, which compares ``this program'' with ``any genes of this size and expression profile''. Patient-level uncertainty and an AUROC companion for the atlas. Nested selection of method, $K$ and $C$. Stage~1 seed and resample stability. A variance decomposition that keeps the estimands of Table~\ref{tab:08-variance} apart.

\subsection{Q3. What does dimensionality reduction add over direct genes?}

\textbf{Current answer.} For a few programs DR clearly adds something, and for broad gene bags it clearly does not. For most programs the verdict depends on the estimand and on how the winning DR cell is chosen (Table~\ref{tab:08-estimands}). Each panel has up to 16 DR cells (4 methods $\times$ 4 $K$; the three small panels have 8--12). ``Best DR'' is the best of these, chosen \emph{in-sample} on the held-out folds it is scored on, so every win count is optimistic (winner's curse).

\begin{table}[H]
\centering
\small
\caption{Biological programs (of 41) where best DR beats the direct-gene probe, by selection rule (rows) and estimand (columns). Quote the three estimands together.}
\label{tab:08-estimands}
\begin{tabularx}{\linewidth}{>{\raggedright\arraybackslash}X c c c}
\toprule
Best of up to 16 DR cells, chosen in-sample by \ldots & pooled LOPO AP & median lift & P20 lift \\
\midrule
pooled LOPO AP (the atlas rule) & 38/41 & 16/41 & 17/41 \\
median lift & not tabulated & 28/41 (30/47 with controls) & 24/41 \\
\bottomrule
\end{tabularx}
\end{table}

Under the pooled-AP rule, the median DR value added is $+0.0221$ pooled AP, $-0.33$\pp{} median lift and $-0.19$\pp{} P20 lift. The rule matters even within one program. For IFN-$\alpha$ the pooled-AP rule picks FA $K = 10$. That probe has 11.36\pp{} lift in P03 and is all-zero (lift 0 by construction) in the other six folds. Its pooled AP (0.2009, against 0.0486 for direct genes) is still the highest of the 16 DR cells, because P03 holds 758 of the 1,557 supported malignant cells and dominates pooled AP. The median-lift rule picks PCA $K = 40$ instead. \src{D9, D1e; 2026-08-31\_d6-d9-meeting-followup}

\begin{figure}[H]
\centering
\includegraphics[width=\linewidth]{figures/fig08b_dr_vs_direct.pdf}
\caption{Best DR versus direct genes, one point per Stage~0 panel, median lift over the seven supported patients; both axes in pp on the same scale. A point above the diagonal means that the best of the panel's DR cells (up to 16, chosen in-sample by median lift) beats the direct-gene probe on the same genes. Each value is at that cell's stored atlas $C$. The three labelled programs on the left edge show the largest gains: allograft rejection (12.56 vs 0.23\pp), apoptosis (11.52 vs 0.005\pp) and IFN-$\alpha$ (11.08 vs 1.27\pp). Most other programs sit near the diagonal (median gap $+0.93$\pp{} when the DR cell is chosen by median lift; $-0.33$\pp{} under the atlas's pooled-AP rule). Four of the six control bags (open squares) sit below it, three of them far below.}
\label{fig:08b}
\end{figure}

\textbf{Evidence.}
\begin{itemize}
\item \emph{Large gains are confined to a few programs and to the highest-prevalence patients} (Figure~\ref{fig:08b}). For allograft, selected DR minus direct genes is $+46.37$\pp{} in P03, $+18.48$ in P02, $+10.33$ in P04 and $+12.69$ in P07. It is $+0.09$ in P01, $+1.57$ in P06 and $+4.12$ in P09. A DR advantage is therefore entangled with prevalence. \src{F2/F4, D1DD5; 2026-09-17\_team-feedback-f1-f7, 2026-09-10\_d1a-descriptive-deep-dives}
\item \emph{Control bags favour direct genes.} On HVG~10,000, direct genes reach 37.9\pp{} and best DR 11.9\pp. On the full\_34 union the values are 27.1 and 12.9\pp. One plausible reading, not yet tested: with thousands of genes, the direct-gene probe keeps signal that a 40-factor basis compresses away. \src{D1e}
\item \emph{Fairness of the comparison.} The program-level comparison is not a relabelled HVG comparison: the 2,719-gene old34 union and HVG~3,000 have Jaccard 0.199, and 65.1\% of the union lies outside HVG~3,000. A mean-$z$ one-score proxy gives 1.68, 5.30 and 1.97\pp{} median lift for allograft, apoptosis and TNF-$\alpha$, against 12.56, 11.52 and 18.06\pp{} for selected DR; the proxy is not UCell, AUCell or ssGSEA, and its Stage~2 is not matched. The comparison is also tilted towards DR: direct-gene LOPO standardizes on the training fold only, which makes it effectively inductive for gene scaling, while the DR bases saw the held-out patient. \src{D12, D1DD9; 2026-08-25\_repository-audit}
\end{itemize}

\textbf{What a publication claim still needs.} Nested or prespecified selection of method, $K$ and $C$, with $C$ tuned per fold. Exact single-score baselines (UCell, AUCell, ssGSEA) with a matched Stage~2. Inductive DR bases for the selected representations. A matched random gene-set null and patient-level uncertainty, so that ``DR adds $x$\pp'' comes with an interval.

\begin{caveatbox}
\begin{itemize}
\item Everything here is descriptive, on a transductive Stage~1 basis (the inductive pilot covers IFN-$\alpha$ only), over seven supported patients. It does not establish inductive deployment, causal biology or recurrence risk.
\item ``Best DR'' is picked on the same held-out folds it is scored on, and the atlas $C$ was chosen on pooled LOPO AP with held-out outcomes. Win counts will shrink under nested selection, and fixed-$C$ atlas cells do not rank methods.
\item Lift ranks with prevalence. The highest-prevalence patients drive both the strong transfer (Q1) and the DR gains (Q3), so shared biology and prevalence are not separated.
\item There are no $p$-values and no patient-level intervals, except the coarse 7-cluster bootstrap of the patient-count curve.
\end{itemize}
\end{caveatbox}

\begin{teambox}
``On the current transductive basis and seven supported patients, transfer above prevalence is widespread, but strong transfer and DR's advantage over direct genes both concentrate in P02, P03 and P04, and none of it has yet been tested against a null.''

\textbf{Likely question:} ``So does DR beat direct genes or not?''\\
\textbf{Answer:} On pooled AP, usually (38/41). On patient-equal median and P20 lift it wins in fewer than half of the programs: 16/41 and 17/41 under the same rule, or 28/41 programs (30/47 with the control bags) on median lift if the winner is picked by median lift. The clear wins are allograft, apoptosis and IFN-$\alpha$, and the clear losses are the broad HVG bags. The winner is picked in-sample, so these counts are upper bounds until selection is nested.

\textbf{Likely question:} ``Isn't the concentration in P02--P04 just prevalence?''\\
\textbf{Answer:} We can't tell yet: lift ranks with prevalence (median $\rho$ 0.821 over 643 representations). Prevalence- and support-matched lift is the test.
\end{teambox}

\subsection*{Next step}
\begin{itemize}
\item Inductive Stage~1 for all 47 atlas-selected representations (W2.6), and the leverage test on the other 38 panels and on FactoSig bases (W2.8, W2.2; W2.4 is done: D26, section~\ref{glance:02b}).
\item Nested DR-vs-direct selection with tuned $C$, plus exact single-score baselines (W3.3, W3.4).
\item A matched random gene-set null, patient-level uncertainty and support-matched lift (W4.2, W4.1, W1.5), then the new patients as external points (W8.1).
\end{itemize}

\subsection{Central theme (draft paragraph for the paper)}

\emph{Draft.} We use classification as a biological probe of post-transplant MDS measurable residual disease. Curated gene programs restrict the gene space (Stage~0). Label-free dimensionality reduction learns low-dimensional bases on all cells (Stage~1). Regularized logistic probes then transfer malignant-versus-normal labels to a held-out patient (Stage~2). The question is at what biological resolution malignant cells show phenotypic convergence that transfers between patients: a curated program, a learned sub-axis of a program, or the individual patient. Across seven supported patients, many curated programs carry held-out signal above the constant-score baseline in most patients. Patients also share one factor basis without patient-private program factors. In the interferon-$\alpha$ program, the signal lives on a combination axis that contrasts two groups of interferon-stimulated genes, and it holds when the basis is learned without the held-out patient. Strong transfer, and the advantage of a learned representation over the program's genes, concentrate in the patients with the highest malignant prevalence. Turning this description into a claim needs four steps: rule out donor-versus-recipient genotype as the separating signal (58,499 of 58,549 normal cells are donor-derived), separate shared biology from prevalence, test against matched random gene sets, and validate on new patients.

\subsection{Claim boundary}

The canonical list is \texttt{research/THEME.md} \S6, and the labels S and N are its stable references.

{\footnotesize
\begin{longtable}{>{\raggedright\arraybackslash}p{1.7cm} >{\raggedright\arraybackslash}p{6.9cm} >{\raggedright\arraybackslash}p{6.4cm}}
\toprule
Question & Supported as descriptive & Not yet claimable \\
\midrule
\endhead
Q1 shared vs patient-specific &
Patients can be pooled on one transductive Stage~1 basis, and no program has a patient-private factor (S15). Lift ranks with prevalence and support across representations (S07). Strong patient $\times$ panel $\times$ method $\times$ K heterogeneity (S01). IFN-$\alpha$ transfer survives an inductive Stage~1 under tuned $C$ or L2, on one panel (S18). Low-support patients carry within-patient DE (S14). &
High lift reflects patient-specific biology rather than prevalence or support (N01). Score overlap is predictive sharedness (N12). The P03 broad-bag axis is a subtype (N11). The highest-prevalence patients shape the Stage~1 loadings (N14). LOPO equals inductive deployment (N07). \\
\addlinespace
Q2 strength of evidence &
Constant-score no-skill AP $=\pi$ (S06). On a frozen basis with fixed $C$, held-out AUPRC rises with Stage~2 cell depth where the probe is not collapsed (S16). Donor downsampling at fixed $C$ can collapse sparse probes to AP $=\pi$ (S11). FA7 is more stable than PC8 under donor downsampling (S12). &
Panels beat matched-size random gene sets (N05). Best configurations are unbiased estimates (N06). More patients help (N08). Cell-level lift is recurrence-risk improvement (N09). Rank-defined TP/FN are clinical misclassifications (N10). The model is invariant to donor versus recipient origin (N03). \\
\addlinespace
Q3 what DR adds (and the IFN-$\alpha$ lens) &
Direct-gene baselines exist for every panel (S08). DR usually improves pooled AP but not consistently median or P20 lift, with in-sample selection (S09). Selected DR reaches clearly positive median lift for allograft and apoptosis where direct genes are near zero (S03). A mean-$z$ proxy is weaker than selected DR, but this is not yet a UCell, AUCell or ssGSEA comparison (S04). The prior is not an HVG alias (S13). Larger K wins most panels, and no method dominates (S02). The IFN-$\alpha$ PCA/FA coverage gap is an L1-at-fixed-$C$ effect on related directions (S10, S17). &
DR universally adds value over direct genes (N04). FA is intrinsically inferior to PCA for IFN-$\alpha$ (N02). ``FA misses patients'' is a property of FA (N13). PC8 or FA7 is a unique, identifiable IFN-$\alpha$ sub-program (N15). \\
\bottomrule
\end{longtable}
}
% Uncaptioned longtable: undo its table-counter step so the numbering stays consecutive.
\addtocounter{table}{-1}

\subsection*{How this was computed}
{\small\raggedright This section synthesizes devlog Results and adds no new model fits. Figure~\ref{fig:08a} counts, per supported patient, the 640 biological DR atlas cells with held-out lift $>10^{-9}$ (all-zero folds score about $10^{-17}$) or $>0.05$. The source is \nolinkurl{analysis/scorecards/stage2_diagnostic/old34_d1_atlas_supported_patient_metrics.csv}. Figure~\ref{fig:08b} and the median-lift row of Table~\ref{tab:08-estimands} use \nolinkurl{old34_d1e_best_dr_vs_direct_gene_selection.csv} in the same directory (best of up to 16 DR cells by median lift, in-sample). The pooled-AP row uses \nolinkurl{meeting_followup_v1/old34_direct_gene_vs_best_dr_summary.csv}. The all-program $\eta^2$ row of Table~\ref{tab:08-variance} summarizes \nolinkurl{factor_patient_matrix_v1/factor_patient_matrix_factor_summary.csv} (non-degenerate PCA and FA factors; patients with $\geq 200$ normal or $\geq 20$ malignant cells). Every atlas cell has one global penalty and $C$, selected on pooled LOPO AP over the same folds (not nested), and uses the transductive Stage~1 scores. Patient-count and tuned-versus-atlas numbers come from \nolinkurl{patient_count_learning_curve_v1} (full design). The code is \nolinkurl{scripts/mrd_stage0_2/research/reports/2026-10-08_questions_report/figs_08_dr_vs_direct.py}; its \texttt{main()} regenerates both figures and prints every derived count. The upstream notebooks are \nolinkurl{stage2_old34_comprehensive_diagnostic_suite.ipynb} and \nolinkurl{stage2_old34_meeting_followup_metric_interactions.ipynb}. The claim boundary follows \texttt{research/THEME.md} \S2--\S6 and \texttt{research/CLAIMS.md}. \src{D1, D1e, D1f, D6, D9, D12--D19, F2/F4; devlogs 2026-09-11\_d1-atlas-783-d1e-d1f, 2026-08-31\_d6-d9-meeting-followup, 2026-10-08\_d15b-full-design, 2026-10-08\_d16-factor-patient-matrix, 2026-10-08\_research-spine-adopted}\par}
             % 8. The three scientific questions (Q1-Q3)
\labelnextsection{glance:09}\section{Remaining TODOs and decisions}
\label{sec:remaining-todos}

\begin{quote}
``Incorporate new patient MRD data from Yu-Hsiang.'' \ldots\ ``1.1 subset to top 4-5 patients with sufficient malignant cells'' \ldots\ ``3. Focus only on patients with large cell counts (2, 3, 4) and rerun the full pipeline on that subset.'' \ldots\ ``Use days since transplant and days before relapse per patient to plot variation across patients and think about how timing could be modeled as a covariate (low priority).'' \ldots\ ``Ankit: end-to-end semi-supervised analysis where factors are learned jointly to discriminate healthy vs malignant (later, with new data).'' \ldots\ ``Review single-cell ML literature for existing frameworks.''
\end{quote}

\begin{answerbox}
Most of what remains waits on Yu-Hsiang's new patients: the email and a pre-registered integration plan are drafted, but the email has not been sent. Option 3's supervision question has a first answer in the patient-count curve: training on two of P02/P03/P04 and scoring the third differs from full LOPO by a median $+0.13$\pp{} over 12 folds, but folds range from $-8.18$ to $+10.09$\pp{} and follow the held-out patient (P02 lower in all four representations, P04 higher in all four). With seven patients, timing cannot be separated from patient identity: median lift tracks prevalence (Spearman $\rho = 0.89$) more closely than days since transplant ($\rho = 0.57$, exact $p = 0.20$) or days before relapse ($\rho = 0.18$, $p = 0.71$), and with seven patients these correlations cannot be told apart. The literature review is done. New since version 1: your healthy-cell downsampling idea is written up as a pre-registered plan that runs with the new patients; new full runs leave out patients without malignant cells and, by your decision of 2026-10-10, keep only patients with at least 100 malignant cells (on the current cohort P02, P03, P04 and P09, so summaries over them run 1.8--2.7 times higher in median lift than today's seven-patient summaries); re-weighting will target the Stage~2 labels in those runs. All descriptive.
\end{answerbox}

\subsection*{What we did}

We collected the TODOs that the other sections do not answer and gave each a status and a next step (Table~\ref{tab:09status}). Two items have new descriptive results, both read from existing artifacts without new compute: Option 3, from the patient-count learning curve, and timing, from the per-sample metadata. \emph{Lift} is the held-out patient's AUPRC minus that patient's malignant prevalence $\pi$, in percentage points (pp); a constant score has AP $= \pi$, so lift 0 means no skill.

% Not a float: in the assembled report a floating table here was placed above this section's heading.
\begin{center}
\begin{minipage}{\linewidth}% keeps the caption and the table on one page
\centering
\small
\captionof{table}{Status of the remaining TODOs on 2026-10-10.}
\label{tab:09status}
\begin{tabularx}{\linewidth}{>{\raggedright\arraybackslash}p{2.7cm} >{\raggedright\arraybackslash}X >{\raggedright\arraybackslash}p{2.6cm} >{\raggedright\arraybackslash}X}
\toprule
TODO & Status & Waiting on & Next step \\
\midrule
New patients & Email and plan drafted; email not sent & You send the email & Lock the plan (11 open choices) \\
Which patients new runs use & Zero-malignant patients out; only patients with $\geq$100 malignant cells kept (both decided) & You (W8.2 lock) & Confirm: patients below 100 left out, not train-only \\
Healthy-cell downsampling & Plan drafted (pre-registered, not locked) & New patients; W8.5 & Runner and tests (W8.6), then lock with the plan \\
Option 1.1 (top 4--5) & Not run; subset rule = the $\geq$100 cut & New patients & Run with the new patients \\
Option 3 (P02/P03/P04) & Supervision part read from the patient-count curve & -- & Merge into the 1.1 subset pipeline \\
Timing & Two figures; descriptive & Manuscript plot & P09 per-draw score check \\
Ankit's joint factors & Later & New patients, locked plan & Pre-register the comparison \\
Literature review & Done (two reviews) & -- & Single-score baselines, matched nulls \\
\bottomrule
\end{tabularx}
\end{minipage}
\end{center}

\subsection{New MRD patients from Yu-Hsiang}

They are the study's only external test: patients that took no part in choosing the panel, the basis, the probe or its C. The draft email (\texttt{email\_yu-hsiang\_draft.md} in \texttt{scripts/\allowbreak mrd\_stage0\_2/\allowbreak workstreams/\allowbreak W8\_new\_cohort/}) is ready for you to edit and send. It asks for (1) the new patients' samples, assay and ADT panel, batch, genome reference, and one \texttt{.h5ad} with raw RNA and ADT counts and our obs columns (\texttt{CN.label}, \texttt{predicted.annotation}, both timing columns); (2) per patient, the evidence behind \texttt{CN.label} (Souporcell, Numbat CNVs, mtDNA, CD5/CD56, cluster reference), donor/recipient calls and sex; (3) what \texttt{Days.from.Relapse} means for P10--P12, which have no relapse sample in our data ($-1{,}329$, $-1{,}158$, $-1{,}382$ days; we checked these values against the cell metadata), and whether our labels are final; (4) the manuscript's timing plot.

\begin{explainbox}{Pre-registration, and transductive vs inductive}
\emph{Lock.} Every analysis choice and a frozen configuration are hashed into a decision entry before unblinding; until then no AP, AUROC or probe score may be computed on a new-patient cell, and anything outside the plan is labelled exploratory.

\emph{Transductive vs inductive.} Our stored Stage 1 bases are transductive: fit once, without labels, on all 60,118 cells including each held-out patient's, so LOPO on them is label transfer to a patient the basis has seen. An inductive basis is refit without the held-out patient, whose cells are then projected. In the plan's primary arm the original bases stay frozen and new cells are projected with the original per-gene statistics, so new patients are held out of every step; probes are fitted at lock time at the stored C and at C tuned by inner LOPO, with L1 and L2.
\end{explainbox}

The plan (\texttt{integration\_plan.md}, same directory, a draft) fixes the estimand (lift per new supported patient; median and 20th percentile, P20, across them, with pooled AP), the primary family (IFN-$\alpha$ PCA40 and FA40 plus the selected allograft-rejection, apoptosis and TNF-$\alpha$/NF-kB representations; Holm within it), the inclusion rules (since 2026-10-10, new full runs and new summarized patients need $\geq$100 malignant cells and both classes; the study's $\geq$10 rule stays for existing results; the 15,560-gene universe is unchanged) and patient-level uncertainty (patient-cluster bootstrap, exact sign test). FactoSig bases cannot yet project new cells, so each FactoSig-promax focus selection has a named PCA/FA fallback. Eleven choices are left to you before the lock, including whether the tuned or the stored-C probe is primary, whether to lock before the data arrive (the plan's preference), and three added on 2026-10-09: the patient-inclusion threshold (decided on 2026-10-10: 100 malignant cells), how the zero-malignant exclusion applies to the frozen probes, and the downsampling plan's own choices (next subsections).

\subsection{Options 1.1 and 3: patient subsets}

\textbf{Option 1.1} needs a rule for ``sufficient malignant cells''. At $\geq$100 it selects P03 (758), P09 (375), P02 (211) and P04 (102); the other supported patients have 22--45. Until new patients arrive it is a four-patient design. The 100 started as an illustrative cut, not a derived rule: it was first used in our review of the team report and in version~1 to make ``sufficient'' concrete, and any cut from 46 to 102 selects the same four patients, because the next supported patient has 45. On 2026-10-10 you chose it for the runs on the new data (W8.5; see \emph{Proposed for the new data} below). It is not the study's support rule, which stays at 10 malignant cells and both classes (\texttt{study\_config.json}).

\textbf{Option 3's premise needs one correction.} P02, P03 and P04 are the \emph{highest-prevalence} patients, but among the \emph{smallest} by cell count (1,462, 2,327 and 1,128 cells; P09 has 17,126 and P01 11,860). They hold 1,071 of 1,569 malignant cells (68.3\%) but 4,917 of 60,118 cells (8.2\%). A strict three-patient rerun would train each LOPO fold on the other two, 2,590--3,789 cells against 57,791--58,990 in the full folds. Because scikit-learn's C does not scale with $n$, under balanced L1 a standardized factor enters only if its training class-mean gap exceeds about $4/(C \cdot n_{\text{train}})$ SD; at the stored C that bar would be 16--22 times higher, so C must be retuned by an inner LOPO whose two folds each train on one patient.

\textbf{The supervision half has a first, descriptive answer.} The full patient-count curve contains, for each of P02, P03 and P04 held out, a training set made of exactly the other two plus the six unsupported patients (20,771 cells, including 12 malignant cells from P05 and P13), so 23,361--24,560 training cells with 325--981 malignant; frozen transductive basis; penalty and C tuned by two-fold inner LOPO. Against the full LOPO fold (Table~\ref{tab:09opt3}), the median difference is $+0.13$\pp{} with 6 of 12 folds higher, but single folds move by $-8.18$ to $+10.09$\pp, and the sign follows the held-out patient, not the method: P02 is lower in all four representations, P04 higher in all four.

\begin{table}[H]
\centering
\small
\caption{Held-out lift (pp) when training on the other two highest-prevalence patients (``pair'') vs on all six other supported patients (``LOPO''). Tuned penalty and C; frozen transductive basis; the six unsupported patients are always in training. diff = pair $-$ LOPO.}
\label{tab:09opt3}
\setlength{\tabcolsep}{3.6pt}
\begin{tabular}{l rrr rrr rrr}
\toprule
 & \multicolumn{3}{c}{Held out P02} & \multicolumn{3}{c}{Held out P03} & \multicolumn{3}{c}{Held out P04} \\
\cmidrule(lr){2-4}\cmidrule(lr){5-7}\cmidrule(lr){8-10}
Representation & pair & LOPO & diff & pair & LOPO & diff & pair & LOPO & diff \\
\midrule
IFN-$\alpha$ PCA40 & 15.94 & 16.74 & $-0.79$ & 18.27 & 15.57 & $+2.69$ & 17.04 & 15.48 & $+1.56$ \\
IFN-$\alpha$ FA40 & 6.88 & 15.05 & $-8.18$ & 13.25 & 16.66 & $-3.41$ & 19.99 & 18.93 & $+1.06$ \\
Apoptosis FactoSig-promax40 & 31.86 & 39.80 & $-7.94$ & 38.14 & 41.86 & $-3.72$ & 26.01 & 15.93 & $+10.09$ \\
Allograft FactoSig-promax40 & 21.99 & 27.76 & $-5.76$ & 39.00 & 37.52 & $+1.48$ & 20.35 & 12.65 & $+7.70$ \\
\bottomrule
\end{tabular}
\end{table}

The converse, leaving P02/P03/P04 out of the Stage~1 basis, is the leave-out arm of the basis-variant pilot (see the section on whether these patients shape the factors). What remains of Option 3 is a basis fit on these three patients only. We recommend running it, if at all, in one patient-subset pipeline with Option 1.1, with the subset rule fixed in advance and the new patients included.

\subsection{Proposed for the new data: healthy-cell downsampling}
\label{sec:09_downsampling}

Your reading notes proposed thinning the healthy cells, stratified by cell type and trajectory, to show that the healthy-to-malignant \emph{ratio}, not only the malignant \emph{count}, makes classification hard, and asked first of all \emph{why} any gain would happen. The ratio varies widely: among the supported patients there are about 2 healthy cells per malignant cell in P03 and about 263 in P01. As you decided, the run waits for the new patients. A draft pre-registration exists (\texttt{healthy\_downsampling\_plan.md} in \texttt{scripts/\allowbreak mrd\_stage0\_2/\allowbreak workstreams/\allowbreak W8\_new\_cohort/}; not locked; hypothesis H-S2-09, open); nothing has been computed.

\begin{explainbox}{Why fewer healthy cells could help, and why it might not}
\emph{Intuition.} Our probe already balances the classes: with \texttt{class\_weight='balanced'} each class carries half of the loss, however many healthy cells there are. Removing healthy cells at random therefore leaves the objective the same in expectation, only fitted on fewer cells, which in scikit-learn is the same problem at a smaller $C$ (the loss is a sum over cells). And AP and AUROC only rank the held-out cells, so the intercept shift that a new class ratio causes cannot move them.

\emph{So a gain can come from four places only.} (1) \emph{Stage~1 geometry}: with fewer healthy cells, malignant cells hold more of the fit's variance (for PCA a class-mean difference $\Delta$ adds $p(1-p)\lVert\Delta\rVert^2$ along $\Delta$, with $p$ the malignant fraction of the fit set), so the malignant direction moves to an earlier factor; this should help most at small $K$, and on panels whose axis is partly made by malignant cells (allograft PCA, section~\ref{glance:02b}). The supported patient with the most healthy cells per malignant cell, P01, also has the largest Stage~1 leverage on all three tested panels; whether that is its cell count, its ratio or its ncMono cells is untested (W2.9). (2) \emph{Which healthy cells define ``normal''}: thinning abundant types more than rare ones moves weight toward healthy cells that resemble malignant cells. (3) \emph{The effective penalty} $C\,n_{\text{train}}$: removed by retuning $C$ on an extended grid. (4) \emph{Sampling noise}. The plan is built to tell these apart. All four are arguments, not results; section~\ref{glance:02} makes the same case from the shaping analyses.
\end{explainbox}

\textbf{What your earlier donor downsampling showed (D10, D11).} D10 thinned the donor-healthy \emph{training} cells of Stage~2 only, stratified by patient $\times$ predicted cell type, on the stored transductive basis, with the held-out patient untouched and $C$ fixed at each representation's atlas value (47,670 fits over 681 targets). At 25\% of the donor cells the sparse probes mostly collapsed to all-zero (mean collapse 0.604 for L1, 0.303 for elastic net), so their AP returned to $\pi$ by construction; L2, which cannot collapse, barely moved (mean $\Delta$AP $-0.0002$, median $+0.0001$). D11 refit the PCA/FA bases on the retained cells: at 25\% the IFN-$\alpha$ supervised factor kept score $r$ 0.952 for FA40 and 0.854 for PCA40. So the Stage-2-only arm has been run once and showed no gain under L2, but with $C$ fixed, which mostly changes the L1 penalty (channel 3). Not yet done: $C$ retuned per condition, a scheme that changes the cell-type mix (channel 2), the malignant count varied apart from the ratio, and a Stage~1 refit with Stage~2 retuned (channel 1); the plan's placements isolate exactly these. \src{2026-08-21\_donor-sensitivity-d10-d11; plan \S2.1}

\textbf{As a paper.} Your catch line is recorded as candidate arc R12 in \texttt{research/\allowbreak paper/\allowbreak OUTLINE.md}, phrased as a question until W8.7 reports; both outcomes are reportable (plan \S10).

\textbf{The design, in brief} (plan \S5--\S8):
\begin{itemize}[leftmargin=1.4em]
\item \emph{Schemes}, compared at the same number of kept healthy cells: random; proportional within patient $\times$ predicted cell type $\times$ pseudotime bin (your ``representative'' scheme); type-equalizing; patient-capped (thins the largest patients, such as P09 and P01, most); sketching optional.
\item \emph{Where}: Stage~1 only, Stage~2 only, or both; Stage~1 transductive and inductive. Since your re-weighting decision (2026-10-10) Stage~2 only is the primary placement, and a Stage~2 sample-weight scheme (each training patient's malignant cells sum to the same weight, likewise its healthy cells; no cell dropped) joins the schemes there; the Stage~1 placements stay as read-outs of why a gain happens. The held-out patient is never thinned, so its $\pi$, and so its lift, is the same in every condition.
\item \emph{Count vs ratio}: malignant retention $\{1, \tfrac12, \tfrac14\}$ $\times$ healthy retention $\{1, \tfrac12, \tfrac14, \tfrac18\}$; per held-out patient, lift is regressed on $\log_2$ malignant count and $\log_2$ healthy-to-malignant ratio, which separates the two effects.
\item \emph{$C$}: retuned per condition by inner LOPO on a grid extended below the current floor, L1 and L2; diagnostic arms at fixed $C$ (should reproduce the earlier donor-downsampling collapse), at fixed $C\,n_{\text{train}}$, and the full-data $C$ path.
\item \emph{Why-read-outs}: malignant share of score energy and the rank of the most discriminative factor (Stage~1); entry margins and the probe direction (Stage~2); the retained cell-type mix and the share of healthy cells in the malignant cells' types.
\item \emph{Success and failure, fixed in advance}: a nested recipe must beat the full-data probe per patient (interval excluding 0, tuned L1 and L2), the ratio effect must be positive, no point of the full-data $C$ path may match the gain, and it must hold inductively, outside the highest-prevalence patients and, in direction, on the new patients. If nothing beats the full-data probe beyond seed noise with $C$ tuned, the idea is falsified, which is still a useful methods result.
\item \emph{Cost}: roughly 100--150 CPU-hours in phased screens; the runner (W8.6) is not written yet.
\end{itemize}

\textbf{Decided: no patients without malignant cells in new runs.} Your note (A-20261009-14) is recorded: P08, P10, P11 and P12 (10,638 cells, 17.7\% of 60,118, all normal) leave the Stage~1 fit set, Stage~2 training and the $C$-tuning folds of every new run, including the patient-subset pipeline and the downsampling arm. Training $n$ changes, so $C$ is retuned. The stored bases are not a new run and keep all 13 patients. One reading of the rule is left to you: whether the frozen probes for the external test also drop these patients (plan \S14, item 10; no recommendation). The allograft and apoptosis leverage run (D26), which refits variants of the stored bases, ran as built on all 13 patients, so it compares directly with the IFN-$\alpha$ leverage test (section~\ref{glance:02b}).

\textbf{Decided (2026-10-10): new full runs keep only patients with at least 100 malignant cells} (A-20261010-07; W8.5). On the current cohort that is P02, P03, P04 and P09; a new patient qualifies by the same rule, from its counts, before any score is computed. Patients below 100 are left out of every step (Stage~1 fit, Stage~2 training, summaries). That is our reading of your note together with your earlier ``subset the patients to do the whole stage0-2 pipeline'' (A-20261009-14); the alternative, keeping them in training only as P05 and P13 are now, stays possible until the plans are locked (W8.2), where you confirm. The study's support rule (10 malignant cells) and every existing result keep the seven supported patients, and a comparison between a new run and an existing result names both cohort definitions. Table~\ref{tab:09incl} lists what the threshold implies on the current cohort, including what it does to the headline summaries. \src{2026-10-10\_decisions-new-run-inclusion-reweighting-lfs}

\begin{table}[H]
\centering
\small
\caption{What each inclusion threshold implies on the current cohort, after the zero-malignant patients are excluded. Counts from the study configuration; ``inner folds'' are the training patients available to tune $C$ by inner LOPO in each outer fold; the sign-test entry is the smallest exact one-sided $p$ possible if every held-out patient has lift $> 0$. Median / P20 lift: the stored LOPO folds of the patient-count curve (tuned penalty and $C$, frozen transductive basis, nothing refit), summarized over the seven or over the four patients (promax40 = FactoSig-promax $K = 40$); computed for this report from \texttt{d15b\_full\_unit\_results.csv} by \texttt{figs\_09\_timing.py}.}
\label{tab:09incl}
\begin{tabularx}{\linewidth}{>{\raggedright\arraybackslash}p{4.2cm} >{\raggedright\arraybackslash}X >{\raggedright\arraybackslash}X}
\toprule
 & $\geq$10 (current rule) or $\geq$20 & $\geq$100 (your choice, 2026-10-10) \\
\midrule
Held-out, summarized patients & 7: P01, P02, P03, P04, P06, P07, P09 & 4: P02, P03, P04, P09 \\
Below the threshold (left out; train-only is the alternative) & P05, P13 & P01, P05, P06, P07, P13 \\
Inner folds for $C$ & 6 & 3 \\
Prevalence of the held-out patients & 0.38--32.57\% & 2.19--32.57\% \\
Highest-prevalence patients among them & 3 of 7 & 3 of 4 \\
Smallest exact sign-test $p$ & 1/128 & 1/16 (cannot reach 0.05) \\
Cells kept if they are left out & 39,347 (1,557 malignant) & 22,043 (1,446 malignant) \\
\addlinespace
\multicolumn{3}{@{}l}{\emph{Median / P20 lift (pp) over the summarized patients:}} \\
\quad IFN-$\alpha$ PCA40 & 8.71 / 1.37 & 15.52 / 11.54 \\
\quad IFN-$\alpha$ FA40 & 7.03 / 1.27 & 15.86 / 11.84 \\
\quad apoptosis promax40 & 10.26 / 2.25 & 27.86 / 10.55 \\
\quad allograft promax40 & 10.12 / 3.06 & 20.20 / 11.08 \\
\bottomrule
\end{tabularx}
\end{table}

The thresholds $\geq$10 and $\geq$20 select the same patients here (P07 has 22 malignant cells) and differ only for new patients with 10--19. At $\geq$100, with the patients below it left out, P01 also leaves the Stage~1 fit set, so the bases change; P01 is the patient with the largest leverage on all three tested panels (section~\ref{glance:02b}). Existing patient summaries all use the seven folds. \emph{What $\geq$100 does to the summaries:} median lift rises 1.8--2.7-fold and P20 lift from 1.3--3.1 to 10.5--11.8\pp, because three of the four summarized patients are the highest-prevalence ones and the low-lift patients P01 and P06 drop out. So the same probe reads much better over the $\geq$100 set without being better; summaries on the new data should say which set they cover, and the seven-patient numbers in this report do not compare with them. P07 is the median patient of the seven in the atlas IFN-$\alpha$ PCA40 summary (11.08\pp) and in three of the four representations above (FA40: P09), so dropping it moves the median.

\textbf{Where this report already shows P01, P06 and P07} (the evidence your note asks for; \texttt{integration\_plan.md} \S4.1 has the same pointers with devlog ids):
\begin{itemize}[leftmargin=1.4em]
\item IFN-$\alpha$ PCA40 at the atlas setting, patient by patient: the yardstick figure in \emph{The cohort and the yardstick} (lift as the gap above $\pi$; P01's probe is all-zero).
\item P07's IFN-$\alpha$ lift: \emph{IFN-$\alpha$: why does FA40 reach only P03 and P04?}, the entry-margin figure (P06's and P07's PC8 margins sit just above the entry bar) and the $C$-policy figure (every patient, P01, P06 and P07 included, gets a factor once $C$ is tuned per fold). Why P07 ranks its cells well: in \emph{What is PC8}, PC8 alone gives 11.08\pp{} in P07 and $+10.3$\pp{} inside HSC/MPP2 alone, and a cell-type oracle reaches only 8.17\pp{} there, so P07's lift is not just which cell types are malignant (unlike P02 and P03). P07's IFN-$\alpha$ score energy is not its cell-type mix either (PCA40 1.31 observed, 1.54 adjusted; section~\ref{glance:02}). It rests on 22 malignant cells: one cell is 4.5\% of them.
\item Across the whole atlas: \emph{The three scientific questions}, the figure ``Where the lift sits'' (each patient's share of DR representations with lift above 0 and above 5\pp).
\item What each patient gains when more training patients are added: \emph{Learning curves}, the per-patient paired-gain figure.
\item Each patient's median lift against prevalence and timing: Figure~\ref{fig:09b} below.
\item Proposed, as an input to the W8.2 confirmation: a cell-bootstrap interval on lift for P01, P06 and P07 (resampling cells within the held-out patient, on the stored scores); new compute, not run.
\end{itemize}

\subsection{Days since transplant and days before relapse}

Timing is recorded per sample. Each supported patient has one MRD draw except P09, which has four, so timing is a property of the patient: with seven patients it cannot be separated from anything else that differs between patients, such as prevalence, sequencing batch or label evidence.

\begin{figure}[H]
\centering
\includegraphics[width=\linewidth]{figures/fig09a_timeline.pdf}
\caption{Each supported patient from transplant (day 0) to relapse (tick). Filled dot: the MRD draw holding the patient's malignant cells, labelled with days to relapse; hollow dots: P09's earlier draws with their malignant counts. Right: malignant / all cells, pooled over draws. Look at the narrow window before relapse against the wide spread since transplant, and at P09.}
\label{fig:09a}
\end{figure}

Every malignant-bearing draw was taken 28--118 days before relapse (P02 closest), while days since transplant range from 76 (P06) to 694 (P03); relapse came 150 to 805 days after transplant (Figure~\ref{fig:09a}). P09 is the only within-patient series: 0, 1, 2 and 372 malignant cells at 586, 399, 294 and 112 days before relapse. In the last draw malignant cells are 8.4\% of cells (372 of 4,410); P09's pooled 2.2\% averages it with three nearly malignant-free draws. Among unsupported patients, the draws with few malignant cells came early: P05 (8 malignant) at 30 days after transplant, P13 (4) at 35.

\begin{figure}[H]
\centering
\includegraphics[width=\linewidth]{figures/fig09b_timing.pdf}
\caption{Per-patient median lift (median over 643 representations) against days since transplant and days before relapse of the malignant-bearing draw, and against prevalence for reference. Bold: highest-prevalence P02, P03, P04. Hollow dots: P09's earlier draws. Exact permutation $p$ over all 5,040 orderings. Look at how the right panel orders the patients more closely than the timing panels; with seven points none of the three correlations can be told apart from the others.}
\label{fig:09b}
\end{figure}

Spearman $\rho$ between median lift and days since transplant is 0.57 (exact $p = 0.20$) and 0.18 for days before relapse ($p = 0.71$), against 0.89 for prevalence (exact $p = 0.012$; Figure~\ref{fig:09b}). With seven patients these correlations cannot be told apart, and prevalence is itself a known correlate of lift (the cohort section), so the comparison is partly circular. Placing P09 at an earlier draw gives 0.54--0.57 and 0.14; without P03, which has both the latest draw and the highest lift, the first falls to 0.31.

\textbf{How timing could be modelled.} In a cell-level probe, a per-patient timing value adds the same amount to every cell of the held-out patient, so by itself it cannot change that patient's ranking or AUPRC. It could act only indirectly, by absorbing between-patient differences while the factor weights are fitted: an adjustment learned from six training patients. Timing belongs in a patient-level model of lift once there are many more patients, and in P09's four draws, the only place where it varies within a patient.

\subsection{Ankit's end-to-end semi-supervised factors}

The literature review turns the idea into a ladder: supervised PCA, then sparse PLS-DA, then a linear supervised autoencoder minimizing reconstruction loss plus $\lambda$ times logistic loss ($\lambda = 0$ is close to the current pipeline, large $\lambda$ is end-to-end; $\lambda$ chosen by inner LOPO). It waits for three reasons. The IFN-$\alpha$ bottleneck was the Stage~2 penalty, not the basis: with C tuned per fold, FA40 has positive lift in 7/7 supported patients under L1, L2 and elastic net (see the IFN-$\alpha$ section). A label-aware Stage 1 must be refit inside every LOPO fold, and choosing $\lambda$ by nested LOPO over seven patients is fragile. Per-fold bases also lose the shared coordinates we use to compare discovery, LOPO and patient-specific models.

\textbf{The bar it must clear}, on the same folds, with pooled AP, median lift and P20 lift read together: the inductive basis with C tuned per fold (the inductive pilot), for IFN-$\alpha$ under tuned L1 PCA40 0.1164 / 10.55 / 1.29\pp{} and FA40 0.0869 / 6.69 / 1.42\pp{} (pooled AP / median / P20); plus the direct-gene probe on the same panel and, once they exist, single-score baselines. The comparison must be pre-registered.

\subsection{Single-cell ML literature: what to adopt}

Two agent-run reviews (2026-10-08) live in \texttt{scripts/\allowbreak mrd\_stage0\_2/\allowbreak workstreams/\allowbreak W9\_literature/}, with each of their 32 recommended actions routed to a task. Table~\ref{tab:09lit} lists the six frameworks from the single-cell review that change what we do.

\begin{table}[H]
\centering
\small
\caption{Frameworks from the single-cell ML review and what we adopt.}
\label{tab:09lit}
\begin{tabularx}{\linewidth}{>{\raggedright\arraybackslash}p{4.4cm} >{\raggedright\arraybackslash}X >{\raggedright\arraybackslash}p{1.4cm}}
\toprule
Framework & What we adopt & Task \\
\midrule
UCell, AUCell, singscore, scanpy \texttt{score\_genes}: one per-cell score per gene set (inductive by construction) & One-feature Stage 2 comparators on the same LOPO folds: a second baseline next to direct genes for ``what does DR add'' & D24 \\
Matched random gene-set nulls (Venet 2011; Goeman and B\"uhlmann 2007) & Random sets matched on size, expression and detection, put through the same selection (``is this prior special?''); within-patient label permutation (``is there any signal?'') & D21, W4.3 \\
MOFA+ multi-group factors (Argelaguet 2020) & Groups = patients: factors active in some patients only and variance explained per factor $\times$ patient, the variance split Q2 asks for & W1.3 \\
Spectra, f-scLVM, PLIER (Kunes 2024) & Gene-set-guided, label-free factors that let a Hallmark set split; to name IFN-$\alpha$ sub-programs (PC8/FA7) & W5.6 \\
Nested selection and patient-level uncertainty (Varma and Simon 2006; Bates 2024; McDermott 2024) & Method, $K$ and C chosen inside inner LOPO; patient-cluster intervals and exact tests; AUROC and AUPRC/$\pi$ next to lift & D23, D22 \\
The cohort paper and label tools (Chen 2025; Souporcell, Numbat) & 58,499 of 58,549 normal cells are donor cells and labels rest on per-patient evidence: audit HLA, ERAP2, chrY and XIST genes, sex match and recipient-normal cells (the review's top priority) & D25 \\
\bottomrule
\end{tabularx}
\end{table}

\begin{caveatbox}
\begin{itemize}
\item \textbf{Timing is $n = 7$}, one value per patient except P09; $\rho$ and exact $p$ are descriptive, P09's position is a choice, and median lift pools 643 representations, including all-zero folds (lift 0 by construction). No causal or recurrence-risk reading.
\item \textbf{Option 3 is read, not rerun:} the 12 pair folds keep the six unsupported patients in training (20,771 cells, 12 of them malignant) and the basis fit on all cells, and cover three held-out patients; single folds move by $-8.18$ to $+10.09$\pp.
\item \textbf{Plans are drafts:} the email is unsent, the integration and downsampling plans are not locked, no new data exist on disk. Nothing about downsampling has been computed; the design's channels and predictions are arguments, not results.
\item \textbf{The literature table summarizes agent-run reviews;} their numbers count only after a routed task reproduces them.
\end{itemize}
\end{caveatbox}

\begin{teambox}
``Most of what is left waits on Yu-Hsiang's patients, and the plan for them is written before we see their labels. Training on two of P02/P03/P04 transfers to the third with a median change of $+0.13$\pp{} from full LOPO, though single folds move by $-8$ to $+10$\pp{} depending on which patient is held out; and with seven patients, timing cannot be told apart from patient identity.''

\textbf{Likely question:} ``Why not just rerun everything on P02, P03 and P04?''\\
\textbf{Answer:} They are among the smallest patients (1,128--2,327 cells), so each fold would train on 2,590--3,789 cells with C tuned on one-patient inner folds. The supervision question already has a first answer inside the patient-count curve (Table~\ref{tab:09opt3}); only a basis fit on the trio is new, and it is better run with Option 1.1 once the new patients arrive.
\end{teambox}

\subsection*{Next step}
\begin{itemize}
\item You: send the Yu-Hsiang email (W8.1); confirm, at the plan lock, that patients below 100 malignant cells are left out of every step rather than kept in training (W8.2; the threshold itself is decided, W8.5); settle the remaining open choices so the integration and downsampling plans can be locked together (W8.2).
\item The downsampling runner and its tests (W8.6), tested on toy data and on labels permuted within patients only; the run itself (W8.7) waits for the new patients.
\item New-cohort readiness code (W8.3) and a descriptive P09 per-draw score check (W1.7).
\end{itemize}

\subsection*{How this was computed}
{\small\raggedright
\emph{Timing:} per-sample days and malignant/normal counts from the D1DD13 table \texttt{analysis/\allowbreak scorecards/\allowbreak stage2\_diagnostic/\allowbreak d1a\_deep\_dive\_20260910/\allowbreak d1dd\_rank\_matched\_\allowbreak confusion\_by\_\allowbreak sample\_collection.csv}, checked against \texttt{stage2\_supervised/\allowbreak multiobjective/\allowbreak discovery\_full\_cohort/\allowbreak cell\_metadata.csv.gz} (EXP-relative). \texttt{d1dd\_supported\_\allowbreak patient\_\allowbreak collection\_\allowbreak context.csv} was not used: its P09 days before relapse do not pair with its days since transplant, whereas in the per-sample table every P09 draw gives relapse at day 723. Prevalence and median raw lift per patient (643 representations, LOPO on the stored transductive bases at each representation's selected penalty and C) and $\rho = 0.89$ from D6 (\texttt{meeting\_followup\_v1/}). Timing $\rho$: Spearman, exact two-sided permutation $p$; nothing refit, so no C policy. \emph{Option 3:} tuned arm (penalty in \{L1, L2\}, C by inner LOPO over the training patients) of \texttt{patient\_count\_\allowbreak learning\_curve\_v1/\allowbreak d15b\_full\_\allowbreak unit\_results.csv}, $k = 2$ units whose subset is the other two of P02/P03/P04, against the $k = 6$ unit for the same patient and representation. Status: workstream READMEs, the email draft, the plans and the review. \emph{Inclusion table and the zero-malignant exclusion:} per-patient malignant and total counts in \texttt{scripts/\allowbreak mrd\_stage0\_2/\allowbreak research/\allowbreak study\_config.json} (checked against the cell metadata), as tabulated in \texttt{integration\_plan.md} \S4.0--\S4.1; healthy-per-malignant ratio $=$ (total $-$ malignant)/malignant. \emph{Downsampling design:} \texttt{healthy\_downsampling\_plan.md} (draft, not locked). Figures, $\rho$ values and Table~\ref{tab:09opt3}: \texttt{figs\_09\_timing.py}. \src{2026-09-10\_d1a-descriptive-deep-dives; 2026-08-31\_d6-d9-meeting-followup; 2026-10-08\_d15b-full-design; 2026-10-08\_d17-basis-variant-pilot; 2026-10-08\_d19-regularization-contract; 2026-10-08\_literature-reviews; 2026-10-08\_team-update-and-wave2-integration; 2026-10-08\_team-report-devlog-pivot; 2026-10-09\_report-reading-notes; 2026-10-09\_w8-healthy-downsampling-plan}\par}
                        % 9. Remaining TODOs and decisions

% ---------------------------------------------------------------- appendices (unnumbered)
\FloatBarrier
\appendix
\section*{Appendix A: Glossary}
\phantomsection
\addcontentsline{toc}{section}{Appendix A: Glossary}
\label{app:glossary}

Terms used anywhere in this report, in alphabetical order. Each entry gives the intuition first and then the exact quantity. The examples refer to the IFN-$\alpha$ panel unless stated otherwise.

\begin{description}[leftmargin=0pt,itemsep=3pt,font=\normalfont\bfseries]

\item[ADT (CLR).] Antibody-derived tags: the CITE-seq surface-protein counts, 161 markers plus 9 isotype controls. Here they are raw counts transformed by a per-cell centred log-ratio (CLR) over all 170 features, without dsb normalization.

\item[Aggregation calls (aggregate, location shift, separate-subtype candidate).] The pre-set verdicts of the patient factor-overlap test for one representation. \emph{Aggregate}: normal weighted overlap $\geq 0.75$, every patient's normal patient-vs-rest overlap $\geq 0.60$, malignant weighted overlap $\geq 0.45$ and malignant gap to the null $<0.15$. \emph{Separate-subtype candidate}: a normal patient-vs-rest overlap below 0.40, or a malignant weighted overlap below 0.35 with a gap of at least 0.15. \emph{Location shift}: everything in between. The subtype gates are checked first, and the malignant cells have no per-patient gate. The gates are magnitude rules; they are not calibrated against a null.

\item[All-zero probe.] An L1 or elastic-net probe whose coefficients are all zero at the chosen $C$. Every held-out cell then gets the same score, so AP $=\pi$ and lift $=0$ \emph{by construction}. It means no factor cleared the entry threshold. It does not mean the representation carries no signal.

\item[AP (average precision).] scikit-learn's step-wise area under the precision--recall curve, $\mathrm{AP}=\sum_n (R_n-R_{n-1})\,P_n$ over the distinct score thresholds. This is the number the report calls AUPRC.

\item[Atlas cell.] One stored panel $\times$ Stage~1 configuration (method and $K$), evaluated on the seven supported patients at a single regularization setting. That setting was chosen per cell by pooled LOPO AP, on the same folds, so the choice is not nested.

\item[AUPRC.] Area under the precision--recall curve. In this report it is always the AP of one held-out patient's cells.

\item[AUROC.] Area under the ROC curve: the probability that a random malignant cell outscores a random normal cell. Its floor is 0.5 whatever the prevalence. It is not the headline metric here; it is planned as a companion column.

\item[Axis loss.] How much of a reference axis (for example PC8 or FA7 of the stored basis) disappears when Stage~1 is refit on a variant, such as one that leaves patients out. It equals $1-R^2$, where $R^2$ is the many-to-one $R^2$ of the reference scores on all factors of the refit basis, computed on the cells both fits share. A value of 0 means the axis is fully recoverable from the refit.

\item[$C$.] scikit-learn's inverse regularization strength. The objective is $C\sum_i s_i\,\mathrm{logloss}_i+\mathrm{penalty}(w)$ with class-balancing weights $s_i$, so a larger $C$ means a weaker penalty. The loss is a \emph{sum} over cells, so the same $C$ acts as a stronger penalty when $n_{\text{train}}$ shrinks. $C$ must therefore be retuned whenever the training size changes.

\item[Cell-type-adjusted ratio.] A patient's score-energy over-representation ratio divided by the ratio its predicted cell-type mix alone would give, if every one of its cells carried its cell type's average score energy. This is an indirect standardization, like a standardized mortality ratio: 1 means the cell-type mix fully accounts for the patient's ratio. ``Accounts for'' is arithmetic on predicted labels, not a cause: batch, collection time and biology still differ between patients of the same mix. The reference (all cells, or only the other patients' cells) matters when one patient holds most of a cell type.

\item[Claim levels (descriptive, inferential, publication).] How far a result may be carried beyond these data (THEME \S6; AGENTS R6). \emph{Descriptive}: an effect size measured on these patients and bases, with no calibrated null, no patient-level interval and no outside cohort, so it describes them and claims nothing beyond them. \emph{Inferential} adds a calibrated null, patient-level uncertainty and nested selection. \emph{Publication} also needs every reviewer alternative addressed, every blocking task done and evidence from committed code. Every result in this report is descriptive.

\item[Composition-matched control.] A random drop of cells that removes the same number of cells \emph{and} the same number of malignant cells as the leave-out set being tested; the leave-out analysis used 20 per set. It separates ``these patients'' from ``this many malignant cells''. Compare size-matched control.

\item[Constant-score no-skill AP.] The AP of a score that is the same for every cell. It equals the prevalence $\pi$ exactly, because all cells form one tie group at recall 1 and precision $\pi$.

\item[Degenerate factor.] A factor with zero score variance (and so no usable scores), which some FA fits produce at large $K$. Degenerate factors are excluded from effective $K$ and from the variance weights.

\item[Determinacy (factor-score determinacy, $\rho$).] For an FA factor: the correlation between the latent factor and its best linear estimate from the genes, $\rho_k=\sqrt{[\Lambda^{\top}\Sigma^{-1}\Lambda]_{kk}}$. Below 0.707 ($\rho^2<0.5$), equally valid score sets can be uncorrelated with each other, so the factor should not be named as a unique program. FA7 has 0.649.

\item[Direct-gene representation (direct-gene probe).] No Stage~1: the Stage~2 probe is fit directly on the panel's standardized genes. This is the no-DR baseline against which DR is compared.

\item[Effective $K$.] The number of non-degenerate factors in a Stage~1 fit; a degenerate factor has zero score variance. IFN-$\alpha$ FA40 has an effective $K$ of 36.

\item[Eigengap (relative).] For a factor with variance $v_k$ (explained variance for PCA, $\Psi$-whitened loading norm for FA): $\min(|v_k-v_{k-1}|,|v_k-v_{k+1}|)/v_k$. At 0.10 or less, the factor is nearly tied with a neighbour and any rotation within that pair fits about equally well, so the individual axis is not identifiable. PC8 has 0.055.

\item[Elastic net (EN).] A penalty that mixes L1 and L2, with \texttt{l1\_ratio} setting the L1 share.

\item[Entry threshold, entry margin.] With balanced class weights $s_i$, the gradient of the data term $C\sum_i s_i\ell_i$ at $w=0$ for a standardized feature $j$ is $C\,(n_{\text{train}}/4)$ times its training class-mean gap. The feature can therefore become nonzero only if $|\text{gap}_j|>4\alpha/(C\,n_{\text{train}})$, with $\alpha=1$ for L1 and $\alpha=$ \texttt{l1\_ratio} for EN. The \emph{entry margin} is $|\text{gap}_j|$ divided by that threshold. If every margin is at most 1, the probe is all-zero. At $C=10^{-4}$ the threshold runs from 0.678 SD (P04 held out) to 0.930 SD (P09 held out).

\item[Eta-squared ($\eta^2$).] The fraction of a factor's score variance explained by patient identity, computed within one label (normal cells only, or malignant cells only). Near 0, the patients look alike on that factor; near 1, the factor mostly separates patients. It is computed over patients with at least 200 normal or at least 20 malignant cells. These are analysis strata, not the support rule.

\item[Factor card.] A one-page record per factor (from the factor-card analysis). It lists the top genes on both loading poles, the orientation, whether the probe selects the factor, the univariate held-out lift per patient, matches in other bases (matched $|r|$, many-to-one $R^2$, Tucker $\phi$), determinacy, eigengap, $\eta^2$ and surface-protein correlations. The cards can be browsed in the factor explorer app.

\item[Factor names (PC8, FA7).] Factor 8 of the PCA $K=40$ basis and factor 7 of the FA $K=40$ basis. The numbers are indices within one fit, so the same number in two fits does not mean the same factor.

\item[Grid floor.] The smallest $C$ on the stored grid, $10^{-4}$, i.e.\ the strongest penalty. Many atlas cells select it, and an L1 probe at the floor is the most likely to be all-zero.

\item[Haufe pattern (activation pattern).] How each input moves with a linear probe's score $s$: the forward pattern $\mathrm{cov}(X,s)/\mathrm{var}(s)$. Unlike the weights, it does not depend on how the inputs are rotated, so probes on different bases can be compared in gene space.

\item[Healthy-to-malignant ratio.] For one patient, (cells $-$ malignant cells)/malignant cells: healthy cells per malignant cell. Prevalence is $\pi = 1/(1+\text{ratio})$. Among the supported patients it runs from about 2 (P03) to about 263 (P01).

\item[Highest-prevalence patients.] P03, P02 and P04 (32.6\%, 14.4\% and 9.0\% malignant). They are not the patients with the most malignant cells (P09 has more than P02 or P04) and not the largest patients.

\item[Identifiable (a single factor).] The data pin the factor's direction down uniquely, so a refit on slightly different cells returns the same factor. A small relative eigengap (PCA) or a determinacy below 0.707 (FA) means it does not: the span of the nearly tied factors is stable, but how that span is cut into single factors is not. Such a factor is named by its two poles and its subspace, not as a unique program.

\item[Inductive (Stage~1).] A basis fit without the held-out patient's cells; those cells are then projected onto it. This is what deployment to a new patient would look like. Compare transductive.

\item[Isotype reference.] For one factor card and label: the largest $|$median within-patient $\rho|$ between the factor and any of the 9 isotype-control antibodies, which bind nothing specific. A marker association at or below this level cannot be told apart from non-specific binding and CLR centring.

\item[$K$.] The number of factors requested from Stage~1: 5, 10, 20 or 40.

\item[Knife-edge.] A fold whose largest entry margin is within $\pm$10\% of 1. A small change in $C$ flips the probe between all-zero and active, and its lift can jump. The regularization analysis counted 136 knife-edge folds in 49 atlas cells.

\item[L1, L2.] The lasso penalty ($\sum_j|w_j|$) and the ridge penalty ($\tfrac12\sum_j w_j^2$). L1 gives sparse probes, and which factor it keeps depends on the basis. L2 keeps every factor with shrunken weights, so a dense L2 probe cannot be read factor by factor.

\item[Leverage test.] An in-silico perturbation of Stage~1: refit the basis without a set of patients, measure how far a reference axis moves (axis loss, matched $|r|$), and rank that against random drops of the same number of cells (size-matched controls) and of the same number of cells \emph{and} malignant cells (composition-matched controls). The decision rule calls it leverage beyond composition when the set's axis loss exceeds all 20 composition-matched controls in at least 4 of the 6 method $\times$ $K$ cells; support for H-S1-05 (the highest-prevalence patients shape the factors) also needs the set's ratio to its controls to exceed every contrast set's (P01, P06, P09, P01+P09) in at least 4 cells. The rule was written after the IFN-$\alpha$ results were known, so it is post hoc there and pre-specified for allograft rejection and apoptosis (D26, section 2b).

\item[Lift.] $\mathrm{AP}-\pi$ for one held-out patient, in percentage points. A lift of 0 is no better than a constant score; the ceiling is $1-\pi$.

\item[LOPO.] Leave-one-patient-out: hold out one patient, train the Stage~2 probe on the cells of the other 12, and rank the held-out patient's cells. On a transductive basis this is label transfer, not deployment.

\item[Many-to-one $R^2$.] The $R^2$ of regressing one reference factor's scores on \emph{all} factors of another basis. A value of 1 means the direction lies in the other basis's span, even if it is split across several factors. PC8 has 0.920 on the FA40 basis.

\item[Matched $|r|$.] The absolute score correlation between a reference factor and its one-to-one (Hungarian) match in another basis. A low matched $|r|$ together with a high many-to-one $R^2$ means the axis was rotated into several factors, not lost.

\item[Nested selection.] Choosing the method, $K$, penalty or $C$ inside the training folds only, so that the held-out patient never influences the choice. The stored atlas is \emph{not} nested: its $C$ and its best-DR cells were picked on the same LOPO folds they are scored on, which inflates the winner's score (the winner's curse).

\item[Normalized lift.] $(\mathrm{AP}-\pi)/(1-\pi)$: the share of the available headroom the probe uses. It is secondary to lift.

\item[Option B.] The 2026-07-03 decision to summarize per-patient lifts with every supported patient weighted equally, using the median and the P20.

\item[Orientation.] The sign that makes a factor point toward malignant cells: the sign of its Stage~2 coefficient, provided that sign is the same in every supported fold where the factor is active ($+1$: malignant cells score higher). Loading signs are arbitrary. PC8's orientation is $-1$, so its malignant-leaning pole is the negative-loading gene group.

\item[Over-representation ratio.] A patient's share of a basis's score energy divided by the patient's share of cells (see score energy share). A value of 1 means the patient contributes to the factor spread exactly in proportion to its size. The section on the highest-prevalence patients uses this name; the factor-by-patient matrix artifacts call it enrichment.

\item[Overlap coefficient.] How much two patients' score distributions on one factor coincide, within one label. With $p_b$ and $q_b$ the two normalized histograms over 32 shared bins (edges at the 1st and 99th percentiles of that label's, i.e.\ that stratum's, scores), the overlap is $\sum_b\min(p_b,q_b)$. A value of 1 means identical distributions and 0 means disjoint ones. The aggregation test summarizes a basis by the variance-weighted mean, over factors, of the median pairwise overlap.

\item[P20 lift.] The 20th percentile of lift over the seven supported patients (pandas linear interpolation). It falls between the second- and third-lowest patient, close to the second-lowest, so it summarizes the lower tail.

\item[Paired gain.] Used in the patient-count learning curve. For each held-out patient, its lift in its LOPO fold ($k = 6$ supported training patients) minus its mean lift over the 15 two-patient training subsets, with the same tuning protocol at both ends.

\item[Patient-cluster bootstrap.] Intervals from resampling held-out \emph{patients}, not cells, with replacement: 2,000 draws, percentile 95\% interval.

\item[Patient-vs-rest overlap.] The overlap coefficient between one patient's score histogram and the pooled histogram of all other patients, within one label. The minimum over patients flags the most distinct patient.

\item[Pooled LOPO AP.] AP computed once over the concatenated held-out predictions of all folds. It is cell-weighted, so the large patients dominate it, and it also depends on how comparable the scores of different fold models are. The stored atlas used it to choose $C$.

\item[pp.] Percentage points: lift $\times$ 100.

\item[Pre-registration.] Writing down the estimand, the comparison and the decision rule before the data that test them are seen. It matters most for the new patients, which are the only external test of the shared axes.

\item[Prevalence ($\pi$).] The fraction of a patient's cells that are malignant. It is also the constant-score no-skill AP.

\item[Principal angle.] The angles between two subspaces, here the spans of two loading matrices; $0^\circ$ means the same subspace. The leave-out analysis reports the largest of the top-8 angles.

\item[Rank $p$.] The one-sided empirical rank of an observed statistic among $n$ random controls: $(1+\#\{\text{controls}\geq\text{observed}\})/(1+n)$. With 20 controls its floor is $1/21\approx0.048$. It is descriptive and not corrected for multiple comparisons.

\item[Robust association.] The rule fixed before computing for factor $\times$ surface-protein correlations: a non-isotype marker whose $|$median within-patient $\rho|$ over the seven supported patients is at least 0.20, with at least 5 of the 7 of the same sign.

\item[Score energy share, enrichment.] How much of a factor's spread comes from one patient. For factor $f$ and patient $p$,
\[\mathrm{SS}_{pf}=\textstyle\sum_{\text{cells of }p}(\text{score}-\text{cohort mean})^2 .\]
The share is $\mathrm{SS}_{pf}/\sum_{p'}\mathrm{SS}_{p'f}$, and the enrichment is that share divided by the patient's share of cells (1 means proportional). Enrichment splits exactly into a location part (squared standardized mean shift) and a spread part (squared SD ratio). At the representation level, SS is pooled over the non-degenerate factors.

\item[Size-matched control.] A random drop of the same number of cells as the leave-out set, ignoring labels; the leave-out analysis used 3 per set. The composition-matched control is the stricter comparator.

\item[Stage~0.] The prior: restrict the genes to one panel (41 curated biological programs or 6 broad control bags) and z-score them within the panel.

\item[Stage~1.] Label-blind dimensionality reduction of the panel: PCA, FA, FactoSig or FactoSig-promax at $K=5,10,20,40$, seed 42, or no reduction (direct genes). The stored bases are transductive.

\item[Stage~2.] A regularized logistic probe (L1, L2 or EN; balanced class weights) on standardized Stage~1 scores, evaluated by LOPO. The discovery and patient-specific views are in-sample.

\item[Supported patient.] A patient with at least 10 malignant cells and both classes: P01, P02, P03, P04, P06, P07 and P09. The other six patients stay in every training fold but are never summarized as held-out patients.

\item[Top patient (a patient ``tops'' a factor).] The patient with the largest over-representation ratio on that factor, among all 13 patients. On IFN-$\alpha$ $K = 40$, P02 tops 31 of 40 PCA factors and 27 of 36 FA factors, P03 none.

\item[Transductive (Stage~1).] A basis fit on all 60,118 cells, including the held-out patient's; labels are never used. LOPO on such a basis is label transfer on a shared atlas.

\item[Transductive optimism.] Lift on a transductive basis minus lift on the matching inductive basis.

\item[Tucker $\phi$ (congruence coefficient).] The uncentered cosine between two loading vectors, $\phi=\sum_g a_g b_g/\sqrt{\sum_g a_g^2\sum_g b_g^2}$, after orienting signs. Literature bands: 0.85--0.94 is fair similarity and 0.95 or above means equal. PC8 and FA7 have $|\phi|=0.691$.

\item[Tuned $C$ (per-fold tuned).] $C$, and in some arms the penalty, chosen separately for each outer LOPO fold by an inner LOPO over the supported training patients, with median inner lift as the criterion. Unlike the stored atlas $C$, this choice never sees the held-out patient.

\item[Univariate held-out lift.] The lift of one factor used alone as the score on the held-out patient, with its sign fixed on the training folds. It measures what that factor carries, whether or not the L1 probe keeps it.

\item[Within-patient $\rho$.] The Spearman $\rho$ between a factor score and a marker, computed inside one patient's cells of one label, so that differences between patients (and their batches) cannot create it. It is summarized as the median over the supported patients together with the number of patients of the same sign. The pooled $\rho$ over all cells of a label mixes within- and between-patient variation.

\end{description}

\src{THEME.md \S1 and \S8; AGENTS.md R5--R7; devlog 2026-10-02\_d14-patient-factor-overlap, 2026-10-08\_d15b-patient-count-curve, 2026-10-08\_d16-factor-patient-matrix, 2026-10-08\_d17-basis-variant-pilot, 2026-10-08\_d17-composition-matched-controls, 2026-10-08\_d18-factor-cards, 2026-10-08\_d19-regularization-contract, 2026-10-09\_d26-leverage-allograft-apoptosis, 2026-10-09\_d27-factor-modality-assoc, 2026-10-09\_d28-celltype-score-shares; workstreams/W9\_literature/2026-10-08\_factor\_interpretability\_review.md}
                                                       % Appendix A: glossary
\section*{Appendix B: Where everything lives}
\phantomsection
\addcontentsline{toc}{section}{Appendix B: Where everything lives}
\label{app:where}

Every number in this report traces to a devlog entry and an artifact. Code and documentation paths are relative to the repository root (\texttt{sc\_classification/}). Artifact paths are relative to the experiment directory EXP (\nolinkurl{experiments/20260525_060508_stage0_mrd_old34_broad_screen_82db5093/}, about 199 GB, gitignored, on the lab server). Devlog entries are \nolinkurl{scripts/mrd_stage0_2/research/devlog/2026/<id>.md}; \texttt{python tools/research\_index.py show <ID>} prints any devlog, analysis (D/F), hypothesis (H-), task (W\#.n) or ask (A-) ID with everything that cites it.

\subsection*{Section by section}

In Table~\ref{tab:where}, code is under \nolinkurl{scripts/mrd_stage0_2/}, artifact directories are under \nolinkurl{EXP/analysis/scorecards/stage2_diagnostic/}, and notebooks are under \nolinkurl{scripts/mrd_stage0_2/notebooks/stage0_2/}. Each section's figures are regenerated by \texttt{figs\_NN\_<topic>.py} in this report's directory (section 2b: \texttt{figs\_02b\_d26.py}).

{\footnotesize
\begin{longtable}{>{\raggedright\arraybackslash}p{2.6cm} >{\raggedright\arraybackslash}p{4.0cm} >{\raggedright\arraybackslash}p{8.4cm}}
\caption{Where each section's evidence lives.}\label{tab:where}\\
\toprule
Section & Devlog entries & Code, artifact dirs, notebook \\
\midrule
\endfirsthead
\toprule
Section & Devlog entries & Code, artifact dirs, notebook \\
\midrule
\endhead
\bottomrule
\endfoot
1. Cohort and yardstick (D6, D8) & 2026-08-31\_d6-d9-meeting-followup & \emph{Code:} \nolinkurl{stage2_supervised/run_stage2_mrd_multiobjective_scorecard.py}\newline \emph{Artifacts:} \nolinkurl{meeting_followup_v1/}\newline \emph{Notebook:} \nolinkurl{stage2_old34_meeting_followup_metric_interactions.ipynb} \\
\addlinespace
2. Do the highest-prevalence patients shape the factors? (D16, D17, D28) & 2026-10-08\_d16-factor-patient-matrix; 2026-10-08\_d17-basis-variant-pilot; 2026-10-08\_d17-composition-matched-controls; 2026-10-09\_d28-celltype-score-shares & \emph{Code:} \nolinkurl{workstreams/W1_patient_aggregation/run_d16_factor_patient_matrix.py}; \nolinkurl{workstreams/W2_stage1_basis_variants/run_d17_basis_variant_pilot.py}; \nolinkurl{workstreams/W1_patient_aggregation/run_d28_celltype_score_shares.py}\newline \emph{Artifacts:} \nolinkurl{factor_patient_matrix_v1/}; \nolinkurl{basis_variant_pilot_ifna_v1/}; \nolinkurl{factor_patient_matrix_celltype_v1/}\newline \emph{Notebook:} -- \\
\addlinespace
2b. What did we do with allograft rejection and apoptosis? (D26) & 2026-10-09\_d26-leverage-allograft-apoptosis (run entry, decision rule); 2026-10-10\_d26-leverage-allograft-apoptosis-done (results, verdict) & \emph{Code:} \nolinkurl{workstreams/W2_stage1_basis_variants/run_d26_leverage_panels.py} (and .sh for screen \texttt{W2\_d26}); this report's \nolinkurl{figs_02b_d26.py}\newline \emph{Artifacts:} \nolinkurl{leverage_allograft_apoptosis_v1/} (scorecards); the run's own figures under \nolinkurl{EXP/analysis/figures/stage2_diagnostic/leverage_allograft_apoptosis_v1/}\newline \emph{Notebook:} -- \\
\addlinespace
3. Should we aggregate patients? (D14) & 2026-10-02\_d14-patient-factor-overlap & \emph{Code:} \nolinkurl{stage2_supervised/patient_factor_overlap.py}\newline \emph{Artifacts:} \nolinkurl{patient_factor_overlap_v1/}\newline \emph{Notebook:} \nolinkurl{stage2_old34_patient_factor_overlap.ipynb} \\
\addlinespace
4. IFN-$\alpha$: why FA40 reaches only P03 and P04 (D19, D17, D16, F3, F6, F9) & 2026-10-08\_d19-regularization-contract; 2026-10-08\_d16-factor-patient-matrix; 2026-09-17\_team-feedback-f1-f7; 2026-10-02\_f9-ifna-pc8-loadings & \emph{Code:} \nolinkurl{workstreams/W3_dr_value_regularization/regularization_contract.py}; \nolinkurl{workstreams/W3_dr_value_regularization/run_d19_regularization_contract.py}\newline \emph{Artifacts:} \nolinkurl{regularization_contract_v1/}; \nolinkurl{factor_patient_matrix_v1/}; \nolinkurl{feedback_deep_dive_20260917/}\newline \emph{Notebook:} \nolinkurl{stage2_old34_feedback_deep_dives_20260917.ipynb} \\
\addlinespace
5. What is PC8? (D18, D27, F9) & 2026-10-08\_d18-factor-cards; 2026-10-09\_d27-factor-modality-assoc; 2026-10-02\_f9-ifna-pc8-loadings; 2026-10-08\_w7-atlas-strip-factor-explorer & \emph{Code:} \nolinkurl{workstreams/W5_factor_interpretation/run_d18_factor_cards.py}; \nolinkurl{workstreams/W5_factor_interpretation/run_d27_factor_modality_assoc.py} (library \nolinkurl{factor_modality_assoc.py}); \nolinkurl{apps/old34_factor_explorer.py}\newline \emph{Artifacts:} \nolinkurl{factor_cards_v1/}; \nolinkurl{factor_modality_assoc_v1/}; \nolinkurl{feedback_deep_dive_20260917/}\newline \emph{Notebook:} \nolinkurl{stage2_old34_feedback_deep_dives_20260917.ipynb} (section 9) \\
\addlinespace
6. Option 2: inductive Stage 1 (D17) & 2026-10-08\_d17-basis-variant-pilot & \emph{Code:} \nolinkurl{workstreams/W2_stage1_basis_variants/run_d17_basis_variant_pilot.py}; \nolinkurl{workstreams/W2_stage1_basis_variants/stage1_refit.py}\newline \emph{Artifacts:} \nolinkurl{basis_variant_pilot_ifna_v1/}\newline \emph{Notebook:} -- \\
\addlinespace
7. Learning curves (D15, D15b) & 2026-10-02\_d15-cell-count-learning-curve; 2026-10-08\_d15b-patient-count-curve; 2026-10-08\_d15b-full-design & \emph{Code:} \nolinkurl{stage2_supervised/cell_count_learning_curve.py}; \nolinkurl{workstreams/W6_learning_curves/run_d15b_patient_count_curve.py} (and .sh for screen)\newline \emph{Artifacts:} \nolinkurl{cell_count_learning_curve_v1/}; \nolinkurl{patient_count_learning_curve_v1/}\newline \emph{Notebook:} \nolinkurl{stage2_old34_patient_factor_overlap.ipynb} (D15) \\
\addlinespace
8. The three scientific questions (D1, D1e, D1f, D3, D6, D9, D12--D19, F2/F4) & 2026-09-11\_d1-atlas-783-d1e-d1f; 2026-08-31\_d6-d9-meeting-followup; 2026-09-10\_d1a-descriptive-deep-dives; 2026-09-17\_team-feedback-f1-f7; 2026-09-03\_d12-d13-gsea-hvg-lowsupport & \emph{Code:} \nolinkurl{stage2_supervised/stage2_sharedness_plotting.py}\newline \emph{Artifacts:} the directory itself (\nolinkurl{old34_*} tables); \nolinkurl{meeting_followup_v1/}; \nolinkurl{d1a_deep_dive_20260910/}; \nolinkurl{deg_hvg_lowsupport_v1/}\newline \emph{Notebook:} \nolinkurl{stage2_old34_comprehensive_diagnostic_suite.ipynb}; \nolinkurl{stage2_old34_meeting_followup_metric_interactions.ipynb} \\
\addlinespace
9. Remaining TODOs and decisions (D1DD13, D6, D15b) & 2026-09-10\_d1a-descriptive-deep-dives; 2026-10-08\_d15b-full-design; 2026-10-08\_literature-reviews; 2026-10-08\_team-update-and-wave2-integration; 2026-10-09\_report-reading-notes; 2026-10-09\_w8-healthy-downsampling-plan & \emph{Code:} \nolinkurl{workstreams/W8_new_cohort/} (email draft, integration plan with the patient-inclusion section, healthy-cell downsampling plan); \nolinkurl{workstreams/W9_literature/} (reviews)\newline \emph{Artifacts:} \nolinkurl{d1a_deep_dive_20260910/}; \nolinkurl{meeting_followup_v1/}; \nolinkurl{patient_count_learning_curve_v1/}\newline \emph{Notebook:} \nolinkurl{stage2_old34_d1a_descriptive_deep_dives_20260910.ipynb} \\
\end{longtable}
}

\subsection*{The project's context layer}

\begin{tabularx}{\linewidth}{>{\raggedright\arraybackslash}p{5.6cm} >{\raggedright\arraybackslash}X}
\toprule
File & What it is \\
\midrule
\texttt{AGENTS.md} & Start here: the repository in brief and the hard rules R1--R12 (long jobs only in your screen, retune C whenever training $n$ changes, claim language) \\
\nolinkurl{scripts/mrd_stage0_2/research/STATUS.md} & Generated dashboard: cohort, workstreams, open tasks by priority, hypotheses and current verdicts, analysis registry, campaigns \\
\texttt{research/THEME.md} & The central question, Q1--Q3, and the canonical claim boundary (\S6) \\
\texttt{research/CLAIMS.md}, \texttt{hypotheses.md} & Hypotheses with their evidence and verdicts \\
\texttt{research/devlog/} & One entry per result: question, estimand, code, artifacts, Results, claim gate, next steps \\
\texttt{research/analyses.csv}, \texttt{campaigns.csv} & Analysis registry (ID $\to$ code, notebook, artifact dir) and compute campaigns with their resume commands \\
\texttt{research/asks/}, \texttt{team\_updates/} & Questions as asked and dated team updates. This report answers A-20261008-01 to -32 (\texttt{asks/\allowbreak 2026-10-08\_team-report.md}); your reading notes on version 1 are A-20261009-01 to -16 (\texttt{asks/\allowbreak 2026-10-09\_report-reading-notes.md}), with a table of where version 2 takes each up \\
\texttt{research/paper/} & The living paper context, drafted alongside the runs: \texttt{README.md} (rules), \texttt{OUTLINE.md} (working titles, results arcs mapped to hypotheses, analyses and devlog ids, methods, experiments still needed) and \texttt{FIGURES.md} (the panel plan). Number-free; THEME \S6 stays the claim boundary \\
\nolinkurl{scripts/mrd_stage0_2/workstreams/} (W0--W9) & Each workstream's charter, backlog (W\#.n tasks), code and conventions \\
\bottomrule
\end{tabularx}

\subsection*{Regenerating this report}

From the repository root, in the lab env (\nolinkurl{sc-classification-rapids-py311}) with the experiment directory available:
\begin{quote}\raggedright
\texttt{bash} \nolinkurl{scripts/mrd_stage0_2/research/reports/2026-10-08_questions_report/build.sh} \texttt{-{}-full}
\end{quote}
This writes \texttt{report\_full.pdf}; without \texttt{-{}-full} the script builds the short core report, \texttt{report\_core.pdf}, and \texttt{-{}-all} builds both plus the team version (\texttt{README.md} in the report directory).
Each figure module can also be run alone, for example \texttt{python} \nolinkurl{scripts/mrd_stage0_2/research/reports/2026-10-08_questions_report/figs_09_timing.py}: its \texttt{main()} rewrites that section's \texttt{figures/<name>.pdf}, plus a \texttt{.png} preview, from read-only tables in EXP. The modules share \texttt{report\_style.py} (palette, type sizes, output paths) and resolve EXP through \nolinkurl{sc_classification.utils.paths}, so the report directory holds no absolute paths. Sections are self-contained: to share a subset, remove the \texttt{\textbackslash input} lines you do not want from \texttt{report.tex} and rebuild. \texttt{build.sh -{}-team} writes \texttt{report\_team.pdf}, which leaves out Appendix~C (your repository decisions).

\subsection*{Browsing the results interactively}

Two read-only Dash apps live in \nolinkurl{scripts/mrd_stage0_2/apps/}. Both bind to 127.0.0.1 by default; to share from the lab server, forward the ports over SSH (for example \texttt{ssh -L 8050:localhost:8050 -L 8051:localhost:8051 <server>}).
\begin{itemize}
\item \textbf{D1 atlas explorer}, \texttt{old34\_d1\_atlas\_explorer.py} (port 8050): the 47 Stage 0 panels $\times$ 17 Stage 1 configurations $\times$ 7 supported patients as patient-specific heatmaps of lift or threshold counts. A cell's view shows the per-fold strip, with all-zero folds marked ``lift = 0 by construction'', the selected penalty and C, and an ``Open factor cards'' link.
\item \textbf{Factor explorer}, \texttt{old34\_factor\_explorer.py} (port 8051): the 1,810 D18 factor cards over 56 representations, with both gene poles, per-patient univariate lift, the fold coefficient strip, and ADT and cell-type tables. URLs address a card directly, for example \nolinkurl{?panel=single_geneset__hallmark_interferon_alpha_response&method=pca&k=40&factor=8} for PC8.
\item \textbf{Checks:} \texttt{-{}-check} validates either app's inputs against the stored tables and prints a pass/fail report without serving.
\item \textbf{Without EXP:} \texttt{-{}-export-bundle DIR} writes a small shareable copy to a new directory (about 2.2 MB of tables for the atlas, about 7.0 MB for the factor cards); a teammate runs the same script with \texttt{-{}-data-dir DIR}. For everything else, \texttt{STATUS.md}, \texttt{CLAIMS.md} and the devlog are in git.
\item \textbf{Long-running servers} go in your screen sessions (\texttt{d1\_atlas} and \texttt{W7\_factor\_explorer}), as recorded in devlog 2026-10-08\_w7-atlas-strip-factor-explorer.
\end{itemize}
                                              % Appendix B: where everything lives
\ifdefined\TEAMBUILD\else
\clearpage
% Appendix C: repository decisions for the project lead only. report.tex leaves this file out of the
% team build (bash build.sh --team), so the shared PDF never carries it.
\section*{Appendix C: For you only (repository decisions)}
\phantomsection
\addcontentsline{toc}{section}{Appendix C: For you only (repository decisions)}
\label{app:foryou}

These repository decisions are yours and do not concern the team's questions. The team build of this report (\texttt{bash build.sh -{}-team}) leaves this appendix out. Status as of 2026-10-10, from git and the lab env; the first two items are now done.

\textbf{1. The Git LFS lines (done, 2026-10-10).} After your ``Let's change the git LFS thing'' (A-20261010-10), all five \texttt{filter=lfs} lines in \texttt{.gitattributes} (\texttt{*.ipynb}, \texttt{*.png}, \texttt{*.pdf}, \texttt{*.pkl}, \texttt{*.h5ad}) were removed; the file now says why in a comment. What that means in practice: nothing in history changes, because git-lfs was never installed here and none of the 341 tracked files is an LFS pointer; the 59 tracked notebooks stay plain JSON and versioned as before; and a collaborator who has git-lfs no longer gets ``should have been pointers'' warnings or turns a notebook into a pointer on their next commit. The change is uncommitted, like the rest of this round. One related choice is new (W0.15): \texttt{*.png} and \texttt{*.pdf} are also listed in \texttt{.gitignore}, so this report's PDF and figures and the team-update figures are not in a clone; decide which of them the repository should carry. \src{2026-10-10\_decisions-new-run-inclusion-reweighting-lfs}

\textbf{2. pytest in the lab env (done, 2026-10-09).} You installed pytest in the lab env, the conda env \texttt{sc-classification-\allowbreak rapids-py311} (Python 3.11.15) that runs every Stage~0--2 runner, the analyses and both apps; this closes W0.4. Tests now run there: \texttt{CONDA\_PREFIX=...\allowbreak/sc-classification-rapids-py311 python -m pytest -q tests} with that env's python (the exact line is in \texttt{AGENTS.md} \S8). The full suite on 2026-10-10, after your \texttt{W2\_d26} run, gave 296 passed and none skipped.

\textbf{3. Pushing.} \texttt{sc\_classification\_gpu} is 3 commits ahead of \texttt{origin}: b60a6ee, the d387eda work-in-progress checkpoint (squash it if you prefer) and f5a8f4d. Decision 1 no longer blocks it: commit this round, then push (W0.3); this unlocks the CI workflow and clean-tree reruns for anything meant for publication, and lets a collaborator clone a repository that starts from \texttt{AGENTS.md} (for agents) and \texttt{README.md} (for people). The umbrella repository already records f5a8f4d locally and has no remote. This report's directory is untracked.

\textbf{4. Updating the repository gradually} (W0.14, your 2026-10-09 note). Each ask gets the smallest change that answers it, as one reviewable diff named in that round's devlog entry; nothing is reorganized wholesale.

\textbf{Next step.} Commit this round (a proposed message comes with each round), push (W0.3), and decide which figures and PDFs to track (W0.15).
                                                   % Appendix C: for you only (left out of the team build)
\fi

\end{document}
